% note.tex — 標準模型と CKM 行列：MS-bar の繰り込みと世代を変える過程（FJ のタドポール、一般の R_xi ゲージ、d = 4）
% コンパイル: lualatex note.tex（目次と参照の解決のため2〜3回）
\documentclass[a4paper,10pt]{ltjsarticle}
\usepackage[margin=20mm]{geometry}
\usepackage[haranoaji,deluxe]{luatexja-preset}
\ltjsetparameter{jacharrange={-3}}
\usepackage{amsmath,amssymb,slashed}
\usepackage{booktabs,tabularx,array}
\usepackage[table]{xcolor}
\usepackage{tikz,ifthen}
\usetikzlibrary{arrows.meta,calc,positioning,decorations.markings,decorations.pathmorphing}
\usepackage[font=small,labelfont=bf]{caption}
\usepackage{graphicx}
\usepackage{placeins}
\usepackage{xurl}
\usepackage{hyperref}
\hypersetup{colorlinks=true,linkcolor=blue!50!black,citecolor=green!40!black,urlcolor=blue!60!black,
  pdftitle={The Standard Model with the CKM matrix: MS-bar renormalization and flavour-changing processes with Fleischer-Jegerlehner tadpoles},pdfauthor={}}

\setlength{\emergencystretch}{3em}
\renewcommand{\jsParagraphMark}{\textcolor{black!65}{\rule[-0.02\zw]{0.62\zw}{0.62\zw}}\hspace{0.35\zw}}
\newcommand{\dd}{\mathrm{d}}
\newcommand{\MSbar}{\overline{\mathrm{MS}}}
\newcommand{\DUV}{\Delta_{\mathrm{UV}}}
\newcommand{\tr}{\mathrm{tr}}
\newcommand{\sw}{s_W}
\newcommand{\cw}{c_W}
\newcommand{\re}{\mathrm{Re}\,}
\newcommand{\im}{\mathrm{Im}\,}

\tikzset{
  photon/.style={thick,decorate,decoration={snake,amplitude=1.1pt,segment length=4.6pt}},
  gluon/.style={thick,decorate,decoration={coil,amplitude=2.3pt,segment length=3.6pt,aspect=0.75}},
  fermion/.style={thick,postaction={decorate},
    decoration={markings,mark=at position 0.58 with {\arrow{Stealth[length=1.9mm]}}}},
  scalar/.style={thick,dashed},
  ghost/.style={thick,densely dotted,postaction={decorate},
    decoration={markings,mark=at position 0.58 with {\arrow{Stealth[length=1.9mm]}}}},
  vtx/.style={circle,fill=black,inner sep=1.3pt},
  lab/.style={font=\small,align=center},
  slab/.style={font=\scriptsize,align=center},
  mom/.style={-{Stealth[length=1.6mm]},thin},
  ctv/.style={circle,draw=black,fill=white,inner sep=0pt,minimum size=7pt,path picture={
      \draw (path picture bounding box.north west) -- (path picture bounding box.south east)
            (path picture bounding box.south west) -- (path picture bounding box.north east);}},
  blob/.style={circle,draw=black,thick,fill=black!12,inner sep=0pt},
}

\title{標準模型と CKM 行列：$\MSbar$ の繰り込みと世代を変える過程\\（FJ のタドポール、一般の $R_\xi$ ゲージ、$d=4$）}
\author{}
\date{2026年10月17日}

\begin{document}
\maketitle

\begin{abstract}
標準模型のノート~\cite{smnote}（以下 SM ノート）は、CKM 行列を単位行列に置いた点で電弱の1ループの繰り込みを調べ、CKM 行列を入れたときの変化をそのノートの量について見積もった（相対 $10^{-7}$ 以下）。本ノートでは CKM 行列を模型に入れ、SM ノートでは扱えなかった CKM 行列そのものの繰り込みと走り、$V=1$ では0になる世代を変える過程を、SM ノートと同じ処方（$\MSbar$、Fleischer--Jegerlehner（FJ）のタドポール、一般の $R_\xi$ ゲージ、$d=4$）と同じ道具で調べる。模型は SM ノートの模型の湯川結合を世代の $3\times3$ の複素行列にしたもので、質量の基底では CKM 行列 $V$（標準のパラメータ表示の角度3個と位相1個）が $W$ と $G^\pm$ のクォークの結合にだけ現れ、ほかの93個の頂点のクラスの式は SM ノートと同じである。手計算では、Feynman 則、$|V_{ij}|^2$ を含む木の崩壊幅、湯川行列の $\MSbar$ の $\beta$ 関数から質量の基底への回転の摂動で導いた $V$ の走り（$K\,\dd V/\dd\ln\mu$ の閉じた式。$K=(4\pi)^2$）、クォークの自己エネルギーの世代の非対角の成分の紫外の極を求めた。$V$ の走りの式は、行列の $\beta$ 関数を数値で積分して対角化した値と相対 $4\times10^{-8}$ 以内で一致し、基準点では $|V_{ub}|$、$|V_{cb}|$、$|V_{td}|$、$|V_{ts}|$ が $K\,\dd\ln|V|/\dd\ln\mu\simeq\frac32(y_t^2+y_b^2)=1.38$ で、Jarlskog 不変量がその2倍で増える。$\MSbar$ の CKM の対抗項は $\delta V=\frac{\DUV}2\dd V/\dd\ln\mu$ で、湯川行列の $\beta$ 関数だけで決まるのでゲージによらない（$V$ は $v$ によらないので FJ のタドポールも入らない）。自己エネルギーの非対角の紫外の極は、GIM の機構で $W$ のループの部分が打ち消し、$G^\pm$ のループの $m_{u_k}^2$ に比例する部分だけが残るので $\xi$ によらない。その値から作った Denner--Sack の on-shell の対抗項の紫外の極は $\MSbar$ の $\delta V$ に等しい。検算と応用の十の手順を定め、手順 (1)〜(7) を行った。模型ファイル \texttt{SM\_CKM.fr} の FeynRules 2.3.49 の93個の頂点のクラスは、$V$ を一般の複素行列とした手計算の表と9個の世代の成分すべてで記号として一致し（わざと5か所を誤らせた表では5個の差）、標準のパラメータ表示を入れた値はユニタリー性の規則によらず、$V=1$ では SM ノートの模型の頂点と93個すべてで一致した。MadGraph 3.8.1 は35個の過程（SM ノートの21個と、$V$ の非対角の成分を要する14個）の木の図を、ユニタリーゲージと Feynman ゲージの両方で手で数えた数だけ作り、世代を変える中性の5個の過程には図を作らなかった。手順 (2) では、$W^+\to u_i\bar d_j$、$t\to d_jW^+$ などの13個の振幅の2乗が FeynCalc で手計算の式と記号として一致し、木のレベルで CP が破れないこと、$u\bar u\to W^+W^-$ が標準のパラメータ表示で $V=1$ の値に戻ること、$t\bar s\to W^+H$ が一般の $R_\xi$ で $\xi_W$ によらず $|V_{ts}|^2$ に比例すること（$W$ と $G^+$ の頂点の CKM 行列の整合）を確かめた。明示的な行列（$t\bar s\to W^+H$ はユニタリーゲージ）とは相対 $6\times10^{-15}$ 以内、MadGraph の12個の部分幅とは $10^{-15}$、断面積とは統計誤差の0.9倍以内で一致した。手順 (3) では、クォークの自己エネルギーを $3\times3$ の世代の行列として一般の $R_\xi$ ゲージで求めた。非対角の紫外の極は手計算の式に記号として一致して $\xi$ によらず、対角の極は CKM 行列を含む質量の走りに一致し、FeynArts の \texttt{SMQCD.mod}（\texttt{\$CKM = True}）の値とも記号として一致した。Denner--Sack の対抗項の紫外の極は $\MSbar$ の $\delta V$ と30桁で一致し、有限部分は $\xi_W$ に依存するが、その $|V_{ij}|$ への効果は $2\times10^{-6}$ 以下で、二つの処方の差 $2.3\times10^{-3}$（$|V_{cb}|$ など、$\mu=M_Z$）より3桁小さい。$V$ を入れた $W$ とトップの自己エネルギーから求めた SM ノートの量の変化は、SM ノートの見積もりと相対 $3\times10^{-12}$ 以内で一致した。手順 (4) では、$\bar u_id_jG^+$ の頂点の極と場の繰り込みから湯川行列の $\MSbar$ の対抗項を本ノートの模型で解き、式~\eqref{eq:betaY} の18個の成分と記号として一致させた（$\xi$ によらない）。頂点の補正は $V_{ij}$ に比例するだけで、世代を混ぜる構造はすべてクォークの場の繰り込みの非対角の成分から来る。一般の公式を複素の湯川行列で組んだ実装とも $2\times10^{-15}$ で一致し、質量、$|V_{ij}|$、$J$ の走りの式は行列の方程式を積分した値と一致した（$|V_{ij}|$ は50桁の計算で $5\times10^{-15}$）。手順 (5) では、$W^+\to u_i\bar d_j$ の1ループの電弱の振幅を一般の $R_\xi$ ゲージで求め、$\MSbar$ の $V$ と外線の世代を混ぜる on-shell の場の定数で組んだ。紫外の極は記号として消え（非対角の場の定数を除くと残る）、$\xi_Z$、$\xi_W$ を変えたときの差は木の振幅の $2\times10^{-11}$ 以下だった（$\xi_A$ の差は光子の質量の2乗に比例して0に向かう）。$\delta V$ を Denner--Sack の対抗項に替えると、差は $u\bar b$、$c\bar b$ で $1.5\times10^{-6}$ になる。FeynArts の \texttt{SM.mod} と FormCalc（on-shell の処方、Feynman ゲージ）とは、頂点の寄与が $2\times10^{-15}$、対抗項の寄与が $1\times10^{-9}$ で一致した。手順 (6) では、外線の運動量を0にした $d_j\to u_ie\bar\nu_e$ と $\mu$ 崩壊の1ループの振幅の比を一般の $R_\xi$ ゲージで求め、紫外で有限でゲージによらないこと、トップの質量を0にした極限で Sirlin の $\frac\alpha\pi\ln\frac{M_Z}{M_A}$ にちょうど一致することを確かめた。トップによる部分は、すべてのチャンネルで一つのゲージによらない関数 $G(\mu)$ と CKM の因子の積になり（$G(M_Z)=2.30\times10^{-3}$）、Denner--Sack の on-shell の $V$ とそのゲージ依存の範囲で一致する。PDG の値から求めた $\MSbar$ の $|V_{ub}|$、$|V_{cb}|$ は、$M_Z$ で $0.23\%$ 小さい。手順 (7) では、$b\to s\gamma$ と $Z\to b\bar s$ の on-shell の1ループの振幅を一般の $R_\xi$ ゲージで求め、外線の世代の混合を入れて、紫外で有限なこと、電磁気のゲージ不変性、$\xi$ によらないこと、GIM の機構を確かめた。双極子と運動量0の $Z$ ペンギンは、Inami--Lim の関数に $m_b^2/M_W^2$ の補正の範囲で一致した。右巻きの非対角の場の定数は、$\Sigma^S$ の係数が左巻きと違う（$2m_im_j$）。手順 (1) で FeynRules の問題を一つ見つけた。CKM 行列のある模型で FeynArts の模型を書き出した後に同じカーネルで UFO を書き出すと、UFO のための Feynman 則の計算が、電荷を保存しない壊れたクラス $\bar duW$（添字が1個の CKM 行列）を94個目として作る。導出は付録に載せる。
\end{abstract}

\tableofcontents

\section{目的と問い}\label{sec:intro}

本ノートの問いは三つである。\textbf{第一に、CKM 行列は $\MSbar$ と FJ のタドポールの処方でどう繰り込まれ、その対抗項と走りはゲージによらないか。第二に、$V=1$ では0になる世代を変える過程の1ループの振幅は本ノートの模型でどうなり、GIM の機構とゲージによらないことがどう現れるか。第三に、CKM 行列と Jarlskog 不変量は高いスケールまでどう走るか。}

第一の問いを立てる理由は、SM ノート~\cite{smnote} が CKM 行列を単位行列に置き、その繰り込みを扱っていないからである。SM ノートは、$V=1$ が繰り込み群で保たれ、それ自体でくりこみ可能な理論であることを前提の理由とし、$V$ を入れたときの変化をそのノートの量について見積もった。$V\ne1$ では湯川結合が世代の行列になり、質量の基底への回転そのものが繰り込みを受ける。on-shell の処方では、外線の場の繰り込みの反エルミートの部分で CKM の対抗項を決める Denner と Sack の処方~\cite{ds90} がよく使われるが、一般の $R_\xi$ ゲージでは物理的な振幅をゲージに依存させることが指摘され~\cite{ggm99,yamada01}、ゲージによらない on-shell の処方がいくつも提案された~\cite{dkr04,ks06}。$\MSbar$ では、CKM の対抗項は湯川行列の $\beta$ 関数から決まり、$V$ は質量の基底に回した湯川行列の関数として走る~\cite{babu87,bhmp99}。SM ノートの処方（$\MSbar$ のパラメータは破れていないラグランジアンのもの、タドポールは FJ）に CKM 行列がどう入るか、自己エネルギーのどの部分がゲージに依存し、どの部分が依存しないかを、SM ノートと同じく一般の $R_\xi$ ゲージで具体的に確かめる。

第二の問いを立てる理由は、世代を変える中性の過程（$b\to s\gamma$、$Z\to b\bar s$、$K$ と $B$ 中間子の混合）が、$V\ne1$ で初めて現れ、しかも1ループで初めて現れるからである。これらの振幅は内線の上型のクォークについての和 $\sum_kV_{ki}^*V_{kj}F(m_{u_k}^2)$ で、ユニタリー性で $F$ の定数の部分が打ち消す（GIM の機構~\cite{gim70}）。個々の図は $\xi$ に依存するので、on-shell の振幅がゲージによらないことは、模型とツールの強い検算になる。クォークの CP の破れ（Jarlskog 不変量 $J$~\cite{jarlskog85}）が現れるのもこれらの過程である。第三の問いは、SM ノートの手順 (8) の走りを CKM 行列に広げるものである。

求める量は五つである。(i) CKM 行列を含む Feynman 則。(ii) $|V_{ij}|^2$ を含む木の崩壊幅と断面積。(iii) 湯川行列の $\MSbar$ の対抗項と $\beta$ 関数、そこから導く質量と $V$ と $J$ の走り、クォークの自己エネルギーの行列の紫外の極、CKM の対抗項。(iv) $W\to u_i\bar d_j$ と、セミレプトニックな崩壊を通じた $|V_{ij}|$ と $\MSbar$ の $V(\mu)$ の関係。(v) 世代を変える1ループの過程（$b\to s\gamma$、$Z\to b\bar s$、$\Delta F=2$ の箱の図）と、高いスケールまでの $V$ と $J$ の走り。$V$ によらない SM ノートの結果（タドポール、$Z$ と光子と $H$ の自己エネルギー、ゲージ結合と $\lambda$ の1ループの $\beta$ 関数など）は求め直さず、引用する。$V$ が入る SM ノートの量（$W$ とトップの自己エネルギー、湯川結合の $\beta$ 関数）は本ノートの模型で $V$ を入れたまま求め直し、SM ノートの見積もりと比べる。一致の基準は SM ノートと同じである。

本ノートの構成は次の通りである。第~\ref{sec:setup} 節で模型と前提を定める。第~\ref{sec:results} 節で (i)〜(iii) の手計算の結果と、(iv)、(v) で求める量の形を述べる。第~\ref{sec:plan} 節でそれらをツールと突き合わせる計画を定め、第~\ref{sec:checks} 節で手順 (1) の結果を示す。第~\ref{sec:discussion} 節で現時点の結果と残る手順をまとめる。導出は\ref{app:rules}〜\ref{app:rg} に、検算の細部は\ref{app:details} にある。

\section{設定}\label{sec:setup}

模型は SM ノートの模型の湯川の項を、世代の行列で書いたものに置き換えたものである。ゲージ固定とゴーストの項を除いたラグランジアン密度は、SM ノートの式~(1) の湯川の項を
\begin{equation}
\mathcal L_Y=-\bigl(\bar Q_LY_d\Phi d_R+\bar Q_LY_u\tilde\Phi u_R+\bar L_LY_\ell\Phi\ell_R+\mathrm{h.c.}\bigr)
\label{eq:Y}
\end{equation}
に置き換えたものである。$Y_u$、$Y_d$、$Y_\ell$ は $3\times3$ の複素行列で、世代の添字は $\bar Q_{Li}(Y_d)_{ij}\Phi d_{Rj}$ のように縮約する。ゲージ群、場と電荷、Higgs 二重項の書き方 $\Phi=(-iG^+,(v+H+iG^0)/\sqrt2)$、ゲージ場の組み替え、ゲージ固定、Faddeev--Popov の項、伝播関数、裸の量の定義は SM ノートの式~(1)〜(8) と同じで、共変微分の符号と電荷の割り当て（$D_\mu=\partial_\mu-ieQA_\mu$、$e>0$、$Q_e=-1$）も同じである。

\paragraph{質量の基底と CKM 行列。}湯川行列を特異値分解して $Y_u=U^u_L\hat y_uU^{u\dagger}_R$、$Y_d=U^d_L\hat y_dU^{d\dagger}_R$ と書く（$\hat y$ は正の対角行列）。ゲージ群の変換と交換する世代の回転 $Q_L\to U^u_LQ_L$、$u_R\to U^u_Ru_R$、$d_R\to U^d_Rd_R$ で運動項とゲージ相互作用は変わらないので、場をこう回すと
\begin{equation}
Y_u=\hat y_u,\qquad Y_d=V\hat y_d,\qquad V\equiv U^{u\dagger}_LU^d_L,\qquad
Q_L=\begin{pmatrix}u_L\\ Vd_L\end{pmatrix}
\label{eq:QL}
\end{equation}
になる。$u$、$d$ は質量の固有状態で、$m_f=y_fv/\sqrt2$ は SM ノートの式~(4) と同じである。$V$ は CKM 行列~\cite{cabibbo63,km73} で、行が上型、列が下型のクォークの世代である。ニュートリノは質量0なので、$L_L$ と $\ell_R$ を回して $Y_\ell$ を対角にでき、レプトンは混合しない。パラメータを数えると、$Y_u$、$Y_d$ の36個の実数のうち、三つの $U(3)$ の回転の27個からバリオン数（$Y$ を変えない共通の位相）を除いた26個が場の定義に吸収され、6個の質量と CKM の4個（角度3個と位相1個）が残る。$V$ はクォークの場の位相の付け替え $u_i\to e^{i\alpha_i}u_i$、$d_j\to e^{i\beta_j}d_j$ で $V_{ij}\to e^{-i\alpha_i}V_{ij}e^{i\beta_j}$ と変わる。物理量はこれによらない $|V_{ij}|$ と $J$ だけで書ける。本ノートでは標準のパラメータ表示~\cite{ck84,pdg24}
\begin{equation}
V=\begin{pmatrix}
c_{12}c_{13} & s_{12}c_{13} & s_{13}e^{-i\delta}\\
-s_{12}c_{23}-c_{12}s_{23}s_{13}e^{i\delta} & c_{12}c_{23}-s_{12}s_{23}s_{13}e^{i\delta} & s_{23}c_{13}\\
s_{12}s_{23}-c_{12}c_{23}s_{13}e^{i\delta} & -c_{12}s_{23}-s_{12}c_{23}s_{13}e^{i\delta} & c_{23}c_{13}
\end{pmatrix}
\label{eq:Vstd}
\end{equation}
を使い（$c_{ij}=\sqrt{1-s_{ij}^2}$）、Jarlskog 不変量を
\begin{equation}
J=\im\bigl(V_{us}V_{cb}V_{ub}^*V_{cs}^*\bigr)=c_{12}c_{23}c_{13}^2s_{12}s_{23}s_{13}\sin\delta
\label{eq:J}
\end{equation}
とする。一般には $\im(V_{ij}V_{kl}V_{il}^*V_{kj}^*)=J\sum_{m,n}\epsilon_{ikm}\epsilon_{jln}$ である。$J\ne0$ がクォークの CP の破れの条件である~\cite{jarlskog85}。

\paragraph{前提。}SM ノートの前提 1〜7 をそのまま使い、2 だけを次のように変える。
\begin{enumerate}
\setcounter{enumi}{1}
\item フェルミオンは3世代で、クォークは CKM 行列 \eqref{eq:QL}、\eqref{eq:Vstd} で混合する。ニュートリノは質量0の左巻きだけの場で、レプトンは混合しない。$V$ のパラメータ $(s_{12},s_{23},s_{13},\delta)$ も $\MSbar$ のパラメータで、破れていないラグランジアンの湯川行列 $Y_u$、$Y_d$ の関数である。
\end{enumerate}
数値の基準点は SM ノートの基準点（$\MSbar$ のパラメータで $\alpha^{-1}=128$、$M_W=80$~GeV、$M_Z=91$~GeV、$M_H=125$~GeV、$m_t=165$~GeV など、$\mu=M_Z$）に、$\mu=M_Z$ の $\MSbar$ の値として
\begin{equation}
s_{12}=0.22501,\qquad s_{23}=0.04183,\qquad s_{13}=0.003732,\qquad \delta=1.147
\label{eq:bench}
\end{equation}
を加えたものとする。値は PDG 2024 の大域的な当てはめ~\cite{pdg24} から取ったが、$\MSbar$ の値として決めたものではない（その関係は手順 (6) で求めた。表~\ref{tab:pdgms}）。このとき
\begin{equation*}
|V|=\begin{pmatrix}0.974350 & 0.225008 & 0.003732\\ 0.224876 & 0.973489 & 0.041830\\ 0.008583 & 0.041109 & 0.999118\end{pmatrix},\qquad J=3.1170\times10^{-5}
\end{equation*}
である。本ノートで新たに使う記号を表~\ref{tab:notation} にまとめる。それ以外の記号は SM ノートの表~1 と同じである。

\begin{table}[tbp]
\centering\small
\caption{本ノートで新たに使う記号（ほかは SM ノートの表~1 と同じ）。}
\label{tab:notation}
\begin{tabularx}{\linewidth}{>{\raggedright\arraybackslash}p{4.1cm}X}
\toprule
\textbf{記号} & \textbf{意味} \\
\midrule
$Y_u$, $Y_d$, $Y_\ell$ & 湯川行列（式~\eqref{eq:Y}）。質量の基底では $Y_u=\hat y_u$、$Y_d=V\hat y_d$ \\
$V$, $V_{ij}$ & CKM 行列（模型ファイル \texttt{Src/SM\_CKM.fr} の \texttt{CKM[i, j]}、$i$ は上型、$j$ は下型の世代）。$V_{ud}$ などとも書く \\
$s_{ij}$, $c_{ij}$, $\delta$ & 標準のパラメータ表示の角度の正弦と余弦、位相（\texttt{s12}、\texttt{s23}、\texttt{s13}、\texttt{dCKM}）。式~\eqref{eq:Vstd} \\
$J$ & Jarlskog 不変量（式~\eqref{eq:J}） \\
$D_u$, $D_d$ & $\mathrm{diag}(y_u^2,y_c^2,y_t^2)$、$\mathrm{diag}(y_d^2,y_s^2,y_b^2)$ \\
$\lambda^{(ij)}_k$ & $V_{ki}^*V_{kj}$（下型の $d_j\to d_i$ の遷移で、内線の上型のクォーク $u_k$ に付く因子） \\
$x_k$ & $m_{u_k}^2/M_W^2$ \\
$\Sigma^L_{ij}$, $\Sigma^R_{ij}$, $\Sigma^S_{ij}$ & クォークの自己エネルギーの世代の行列（式~\eqref{eq:Gammaq}） \\
$\delta V$, $\delta V^{\mathrm{DS}}$ & $\MSbar$ の CKM の対抗項（式~\eqref{eq:dV}）と Denner--Sack の on-shell の対抗項（式~\eqref{eq:dVds}） \\
\bottomrule
\end{tabularx}
\end{table}

\section{結果}\label{sec:results}

\subsection{Feynman 則}\label{sec:rules}

式~\eqref{eq:QL} を SM ノートのラグランジアンに入れると、頂点のクラスは SM ノートと同じ93個で、そのうち CKM 行列が入るのは表~\ref{tab:rules-ckm} の4個だけである（\ref{app:rules}）。荷電のクォークの電流 $\bar Q_L\gamma^\mu W^\pm Q_L$ は上と下の成分を結ぶので $V$ を1個含み、湯川の項の $G^\pm$ の部分は $\bar u_LV\hat y_dd_R$ と $\bar d_LV^\dagger\hat y_uu_R$ を含む。中性の電流（$A$、$Z$、$G$）と $H$、$G^0$ の結合は、二重項の下の成分どうしを結ぶので $V^\dagger V=1$ を含み、世代について対角になる。ほかの89個のクラスの式は SM ノートの表~2〜5 と同じである。したがって木のレベルには世代を変える中性の過程がない。図~\ref{fig:ckmdiag} の (a)、(b) が表の頂点、(c)〜(e) が第~\ref{sec:fcnc} 節の1ループの図の型である。

\begin{table}[tbp]
\centering\small\renewcommand{\arraystretch}{1.75}
\caption{CKM 行列を含む頂点（4個）。$i$ は上型、$j$ は下型のクォークの世代で、色の $\delta$ は省いた。ほかの89個のクラス（レプトンの荷電の頂点を含む）は SM ノートの表~2〜5 と同じで、世代について対角である。}
\label{tab:rules-ckm}
\begin{tabular}{ll}
\toprule
頂点 & Feynman 則 \\
\midrule
$\bar u_id_jW_\mu$ & $\dfrac{ie}{\sqrt2\sw}V_{ij}\gamma^\mu P_L$ \\
$\bar d_ju_i\bar W_\mu$ & $\dfrac{ie}{\sqrt2\sw}V_{ij}^*\gamma^\mu P_L$ \\
$\bar u_id_jG^+$ & $\dfrac{\sqrt2}{v}V_{ij}\bigl(m_{u_i}P_L-m_{d_j}P_R\bigr)$ \\
$\bar d_ju_iG^-$ & $\dfrac{\sqrt2}{v}V_{ij}^*\bigl(m_{d_j}P_L-m_{u_i}P_R\bigr)$ \\
\bottomrule
\end{tabular}
\end{table}

\begin{figure}[tbp]
\centering
\begin{tikzpicture}[scale=0.9]
% (a) W vertex
\begin{scope}[shift={(0,0)}]
  \draw[fermion] (-0.85,-0.6) -- (0,0); \draw[fermion] (0,0) -- (-0.85,0.6); \draw[photon] (0,0) -- (1.1,0);
  \node[vtx] at (0,0) {};
  \node[font=\scriptsize] at (-1.05,-0.75) {$d_j$}; \node[font=\scriptsize] at (-1.05,0.75) {$u_i$}; \node[font=\scriptsize] at (1.35,0.25) {$W$};
  \node[slab] at (0.15,-1.35) {(a) $\bar u_id_jW$\\$\propto V_{ij}$};
\end{scope}
% (b) G+ vertex
\begin{scope}[shift={(3.2,0)}]
  \draw[fermion] (-0.85,-0.6) -- (0,0); \draw[fermion] (0,0) -- (-0.85,0.6); \draw[scalar] (0,0) -- (1.1,0);
  \node[vtx] at (0,0) {};
  \node[font=\scriptsize] at (-1.05,-0.75) {$d_j$}; \node[font=\scriptsize] at (-1.05,0.75) {$u_i$}; \node[font=\scriptsize] at (1.35,0.25) {$G^+$};
  \node[slab] at (0.15,-1.35) {(b) $\bar u_id_jG^+$\\$\propto V_{ij}(m_{u_i}P_L-m_{d_j}P_R)$};
\end{scope}
% (c) off-diagonal self-energy
\begin{scope}[shift={(6.6,0)}]
  \draw[fermion] (-1.15,0) -- (-0.55,0); \draw[fermion] (-0.55,0) -- (0.55,0); \draw[fermion] (0.55,0) -- (1.15,0);
  \draw[photon] (-0.55,0) arc (180:0:0.55);
  \node[vtx] at (-0.55,0) {}; \node[vtx] at (0.55,0) {};
  \node[font=\scriptsize] at (-1.0,0.22) {$d_j$}; \node[font=\scriptsize] at (1.0,0.22) {$d_i$};
  \node[font=\scriptsize] at (0,-0.22) {$u_k$}; \node[font=\scriptsize] at (0,0.82) {$W^-$、$G^-$};
  \node[slab] at (0,-1.0) {(c) $\Sigma_{ij}$：$\lambda^{(ij)}_k=V_{ki}^*V_{kj}$};
\end{scope}
% (d) penguin
\begin{scope}[shift={(10.2,0)}]
  \draw[fermion] (-1.25,0) -- (-0.65,0); \draw[fermion] (-0.65,0) -- (0.65,0); \draw[fermion] (0.65,0) -- (1.25,0);
  \draw[photon] (-0.65,0) arc (180:0:0.65);
  \draw[photon] (0,0) -- (0,-0.75);
  \node[vtx] at (-0.65,0) {}; \node[vtx] at (0.65,0) {}; \node[vtx] at (0,0) {};
  \node[font=\scriptsize] at (-1.1,0.22) {$b$}; \node[font=\scriptsize] at (1.1,0.22) {$s$};
  \node[font=\scriptsize] at (-0.35,-0.22) {$u_k$}; \node[font=\scriptsize] at (0,0.92) {$W^-$、$G^-$};
  \node[font=\scriptsize] at (0.35,-0.65) {$\gamma$、$Z$};
  \node[slab] at (0,-1.25) {(d) $b\to s\gamma$、$Z\to b\bar s$};
\end{scope}
% (e) box
\begin{scope}[shift={(13.9,0)}]
  \draw[fermion] (-1.1,0.55) -- (-0.55,0.55); \draw[fermion] (-0.55,0.55) -- (0.55,0.55); \draw[fermion] (0.55,0.55) -- (1.1,0.55);
  \draw[fermion] (1.1,-0.55) -- (0.55,-0.55); \draw[fermion] (0.55,-0.55) -- (-0.55,-0.55); \draw[fermion] (-0.55,-0.55) -- (-1.1,-0.55);
  \draw[photon] (-0.55,0.55) -- (-0.55,-0.55); \draw[photon] (0.55,0.55) -- (0.55,-0.55);
  \node[vtx] at (-0.55,0.55) {}; \node[vtx] at (0.55,0.55) {}; \node[vtx] at (-0.55,-0.55) {}; \node[vtx] at (0.55,-0.55) {};
  \node[font=\scriptsize] at (-1.0,0.78) {$b$}; \node[font=\scriptsize] at (1.0,0.78) {$s$};
  \node[font=\scriptsize] at (-1.0,-0.78) {$s$}; \node[font=\scriptsize] at (1.0,-0.78) {$b$};
  \node[font=\scriptsize] at (0,0.78) {$u_k$}; \node[font=\scriptsize] at (0,-0.78) {$u_l$};
  \node[slab] at (0,-1.25) {(e) $\Delta F=2$ の箱\\$W$、$G^\pm$ の組};
\end{scope}
\end{tikzpicture}
\caption{CKM 行列の入る頂点（(a)、(b)、表~\ref{tab:rules-ckm}）と、世代を変える1ループの図の型。(c) は下型のクォークの自己エネルギーの世代の非対角の成分（$d_j$ が入り $d_i$ が出る）で、$W^\pm$、$G^\pm$ と内線の上型のクォーク $u_k$ のループに $\lambda^{(ij)}_k=V_{ki}^*V_{kj}$ が掛かる。(d) は $b\to s\gamma$ と $Z\to b\bar s$ の頂点の図の一つで、光子と $Z$ はループの $u_k$ にも $W^\pm$、$G^\pm$ にも付く。(e) は $B_s$ の混合の箱の図（$b\bar s\to s\bar b$。下の線は反クォークで、左から $\bar s$ が入り右へ $\bar b$ が出る）で、縦の線は $W^\pm$ か $G^\pm$ である。運動量は左から右へ、矢印はフェルミオン数の流れである。}
\label{fig:ckmdiag}
\end{figure}

\subsection{木のレベルの崩壊幅と断面積}\label{sec:tree}

木のレベルの振幅は、荷電の頂点ごとに $V_{ij}$ か $V_{ij}^*$ を1個持つ。$W$ が1本の過程の振幅の2乗は $|V_{ij}|^2$ に比例し、場の位相の付け替えによらない。SM ノートの式~(11) の $W$ とトップの崩壊幅は（\ref{app:tree}）
\begin{equation}
\Gamma(W^+\to u_i\bar d_j)=|V_{ij}|^2\,\Gamma_{V=1}(W^+\to u_i\bar d_j),\qquad
\Gamma(t\to d_jW^+)=|V_{tj}|^2\,\Gamma_{V=1}(t\to d_jW^+)
\label{eq:widthsckm}
\end{equation}
になる。$\Gamma_{V=1}$ は SM ノートの式~(11) に終状態の質量 $m_{u_i}$、$m_{d_j}$ を入れたものである。下型のクォークの質量を0にすると、行のユニタリー性 $\sum_j|V_{ij}|^2=1$ から $\sum_j\Gamma(t\to d_jW^+)$ は $V=1$ の $\Gamma(t\to bW^+)$ に等しく、$W$ のハドロンへの崩壊幅は、開いている $u$、$c$ の2行で $2N_c\Gamma(W\to\ell\nu)$ になる。クォークとレプトンの質量を0にすると、$s$ チャンネルの $W$ の $u_i(p_1)\bar d_j(p_2)\to\nu_e(k_1)e^+(k_2)$ と、$t$ チャンネルの $W$ の $u(p_1)\bar u(p_2)\to d(k_1)\bar s(k_2)$ は（\ref{app:tree}）
\begin{equation}
\begin{gathered}
\sum|\mathcal M|^2(u_i\bar d_j\to\nu_ee^+)=N_cg^4|V_{ij}|^2\frac{u^2}{|s-M_W^2+iM_W\Gamma_W|^2},\\
\sum|\mathcal M|^2(u\bar u\to d\bar s)=N_c^2g^4|V_{ud}|^2|V_{us}|^2\frac{u^2}{(t-M_W^2)^2}
\end{gathered}
\label{eq:qqbar}
\end{equation}
である（$t=(p_1-k_1)^2$、$u=(p_1-k_2)^2$。和はスピンと色について）。前者のスピンと色の平均の断面積は $\sigma=g^4|V_{ij}|^2s/(576\pi|s-M_W^2+iM_W\Gamma_W|^2)$ である。$e^+$ は $\bar d_j$ の向きに出やすい（$u^2\propto(1+\cos\theta)^2$、$\theta$ は $u_i$ と $\nu_e$ の角）。$t$ チャンネルに下型のクォークが入る $u\bar u\to W^+W^-$ では、振幅が $\sum_k|V_{uk}|^2A(m_{d_k})$ になり、$d_k$ が質量0なら $V=1$ の振幅に等しい。$u\bar u\to d\bar s$ のような $t$ チャンネルの $W$ の交換は $V_{ud}^*V_{us}$ に比例する（入る $u$ を消して $d$ を作るのは $\bar d u\bar W$ の頂点で $V_{ud}^*$、入る $\bar u$ を消して $\bar s$ を作るのは $\bar usW$ の頂点で $V_{us}$。$t\to d_jW^+$ の振幅は $V_{tj}^*$ に比例する）。木のレベルでは、CP で移り合う過程（$t\to d_jW^+$ と $\bar t\to\bar d_jW^-$）の振幅の2乗が等しく、CP の破れは現れない。崩壊幅の CP の非対称には、位相の違う二つの振幅と吸収的な部分（ループの虚部）が要る。これらの過程の図と、振幅に掛かる CKM の因子を図~\ref{fig:treeckm} に示す。

\begin{figure}[tbp]
\centering
\begin{tikzpicture}[scale=0.85]
% (a) W+ -> u_i dbar_j
\begin{scope}[shift={(0.6,0)}]
  \draw[photon] (-1.0,0) -- (0,0); \draw[fermion] (0,0) -- (0.85,0.6); \draw[fermion] (0.85,-0.6) -- (0,0);
  \node[vtx] at (0,0) {};
  \node[font=\scriptsize] at (-0.6,0.3) {$W^+$};
  \node[font=\scriptsize,anchor=west] at (0.85,0.6) {$u_i$}; \node[font=\scriptsize,anchor=west] at (0.85,-0.6) {$\bar d_j$};
  \node[slab] at (0,-1.2) {(a) $W^+\to u_i\bar d_j$\\$\propto V_{ij}$};
\end{scope}
% (b) t -> d_j W+
\begin{scope}[shift={(4.0,0)}]
  \draw[fermion] (-1.0,0) -- (0,0); \draw[fermion] (0,0) -- (0.85,0.6); \draw[photon] (0,0) -- (0.85,-0.6);
  \node[vtx] at (0,0) {};
  \node[font=\scriptsize] at (-0.6,0.25) {$t$};
  \node[font=\scriptsize,anchor=west] at (0.85,0.6) {$d_j$}; \node[font=\scriptsize,anchor=west] at (0.85,-0.6) {$W^+$};
  \node[slab] at (0,-1.2) {(b) $t\to d_jW^+$\\$\propto V_{tj}^*$};
\end{scope}
% (c) u_i dbar_j -> nu_e e+
\begin{scope}[shift={(7.3,0)}]
  \draw[fermion] (-0.9,-0.65) -- (0,0); \draw[fermion] (0,0) -- (-0.9,0.65);
  \draw[photon] (0,0) -- (1.3,0);
  \draw[fermion] (1.3,0) -- (2.2,0.65); \draw[fermion] (2.2,-0.65) -- (1.3,0);
  \node[vtx] at (0,0) {}; \node[vtx] at (1.3,0) {};
  \node[font=\scriptsize,anchor=east] at (-0.9,-0.65) {$u_i$}; \node[font=\scriptsize,anchor=east] at (-0.9,0.65) {$\bar d_j$};
  \node[font=\scriptsize,anchor=west] at (2.2,0.65) {$\nu_e$}; \node[font=\scriptsize,anchor=west] at (2.2,-0.65) {$e^+$};
  \node[font=\scriptsize] at (0.65,0.32) {$W^+$};
  \node[slab] at (0.65,-1.2) {(c) $u_i\bar d_j\to\nu_ee^+$\\$\propto V_{ij}^*$};
\end{scope}
% (d) u ubar -> d sbar (t-channel W)
\begin{scope}[shift={(1.2,-3.4)}]
  \draw[fermion] (-1.0,0.55) -- (0,0.55); \draw[fermion] (0,0.55) -- (1.0,0.55);
  \draw[fermion] (0,-0.55) -- (-1.0,-0.55); \draw[fermion] (1.0,-0.55) -- (0,-0.55);
  \draw[photon] (0,0.55) -- (0,-0.55);
  \node[vtx] at (0,0.55) {}; \node[vtx] at (0,-0.55) {};
  \node[font=\scriptsize,anchor=east] at (-1.0,0.55) {$u$}; \node[font=\scriptsize,anchor=west] at (1.0,0.55) {$d$};
  \node[font=\scriptsize,anchor=east] at (-1.0,-0.55) {$\bar u$}; \node[font=\scriptsize,anchor=west] at (1.0,-0.55) {$\bar s$};
  \node[font=\scriptsize] at (0.38,0) {$W$};
  \node[slab] at (0,-1.2) {(d) $u\bar u\to d\bar s$\\$\propto V_{ud}^*V_{us}$};
\end{scope}
% (e1) u ubar -> W+ W-: s-channel gamma, Z
\begin{scope}[shift={(4.6,-3.4)}]
  \draw[fermion] (-0.9,-0.65) -- (0,0); \draw[fermion] (0,0) -- (-0.9,0.65);
  \draw[photon] (0,0) -- (1.2,0);
  \draw[photon] (1.2,0) -- (2.0,0.65); \draw[photon] (1.2,0) -- (2.0,-0.65);
  \node[vtx] at (0,0) {}; \node[vtx] at (1.2,0) {};
  \node[font=\scriptsize,anchor=east] at (-0.9,-0.65) {$u$}; \node[font=\scriptsize,anchor=east] at (-0.9,0.65) {$\bar u$};
  \node[font=\scriptsize,anchor=west] at (2.0,0.65) {$W^+$}; \node[font=\scriptsize,anchor=west] at (2.0,-0.65) {$W^-$};
  \node[font=\scriptsize] at (0.6,0.32) {$A$、$Z$};
  \node[slab] at (0.6,-1.2) {(e) $u\bar u\to W^+W^-$\\$s$ の $\gamma$、$Z$};
\end{scope}
% (e2) t-channel d_k
\begin{scope}[shift={(9.0,-3.4)}]
  \draw[fermion] (-0.8,-0.95) -- (0,-0.5); \draw[fermion] (0,-0.5) -- (0,0.5); \draw[fermion] (0,0.5) -- (-0.8,0.95);
  \draw[photon] (0,-0.5) -- (1.0,-0.95); \draw[photon] (0,0.5) -- (1.0,0.95);
  \node[vtx] at (0,-0.5) {}; \node[vtx] at (0,0.5) {};
  \node[font=\scriptsize,anchor=east] at (-0.8,-0.95) {$u$}; \node[font=\scriptsize,anchor=east] at (-0.8,0.95) {$\bar u$};
  \node[font=\scriptsize,anchor=west] at (1.0,-0.95) {$W^+$}; \node[font=\scriptsize,anchor=west] at (1.0,0.95) {$W^-$};
  \node[font=\scriptsize] at (0.3,0) {$d_k$};
  \node[slab] at (0.1,-1.45) {$t$ の $d_k$（$k=d,s,b$）\\$\propto|V_{uk}|^2$};
\end{scope}
% (f1) t sbar -> W+ H: s-channel W
\begin{scope}[shift={(0.6,-6.9)}]
  \draw[fermion] (-0.9,-0.65) -- (0,0); \draw[fermion] (0,0) -- (-0.9,0.65);
  \draw[photon] (0,0) -- (1.1,0);
  \draw[photon] (1.1,0) -- (1.85,0.65); \draw[scalar] (1.1,0) -- (1.85,-0.65);
  \node[vtx] at (0,0) {}; \node[vtx] at (1.1,0) {};
  \node[font=\scriptsize,anchor=east] at (-0.9,-0.65) {$t$}; \node[font=\scriptsize,anchor=east] at (-0.9,0.65) {$\bar s$};
  \node[font=\scriptsize,anchor=west] at (1.85,0.65) {$W^+$}; \node[font=\scriptsize,anchor=west] at (1.85,-0.65) {$H$};
  \node[font=\scriptsize] at (0.55,0.32) {$W^+$};
  \node[slab] at (0.55,-1.2) {(f) $t\bar s\to W^+H$\\$s$ の $W$（$\xi_W$）};
\end{scope}
% (f2) s-channel G+
\begin{scope}[shift={(5.1,-6.9)}]
  \draw[fermion] (-0.9,-0.65) -- (0,0); \draw[fermion] (0,0) -- (-0.9,0.65);
  \draw[scalar] (0,0) -- (1.1,0);
  \draw[photon] (1.1,0) -- (1.85,0.65); \draw[scalar] (1.1,0) -- (1.85,-0.65);
  \node[vtx] at (0,0) {}; \node[vtx] at (1.1,0) {};
  \node[font=\scriptsize,anchor=east] at (-0.9,-0.65) {$t$}; \node[font=\scriptsize,anchor=east] at (-0.9,0.65) {$\bar s$};
  \node[font=\scriptsize,anchor=west] at (1.85,0.65) {$W^+$}; \node[font=\scriptsize,anchor=west] at (1.85,-0.65) {$H$};
  \node[font=\scriptsize] at (0.55,0.3) {$G^+$};
  \node[slab] at (0.55,-1.2) {$s$ の $G^+$（$\xi_W$）\\$R_\xi$ だけ};
\end{scope}
% (f3) t-channel top
\begin{scope}[shift={(9.4,-6.9)}]
  \draw[fermion] (-0.9,0.55) -- (0,0.55); \draw[scalar] (0,0.55) -- (0.9,0.55);
  \draw[fermion] (0,0.55) -- (0,-0.55);
  \draw[fermion] (0,-0.55) -- (-0.9,-0.55); \draw[photon] (0,-0.55) -- (0.9,-0.55);
  \node[vtx] at (0,0.55) {}; \node[vtx] at (0,-0.55) {};
  \node[font=\scriptsize,anchor=east] at (-0.9,0.55) {$t$}; \node[font=\scriptsize,anchor=west] at (0.9,0.55) {$H$};
  \node[font=\scriptsize,anchor=east] at (-0.9,-0.55) {$\bar s$}; \node[font=\scriptsize,anchor=west] at (0.9,-0.55) {$W^+$};
  \node[font=\scriptsize] at (0.22,0) {$t$};
  \node[slab] at (0,-1.2) {$t$ チャンネルの $t$};
\end{scope}
% (f4) u-channel s
\begin{scope}[shift={(12.4,-6.9)}]
  \draw[fermion] (-0.9,0.55) -- (0,0.55); \draw[photon] (0,0.55) -- (0.9,0.55);
  \draw[fermion] (0,0.55) -- (0,-0.55);
  \draw[fermion] (0,-0.55) -- (-0.9,-0.55); \draw[scalar] (0,-0.55) -- (0.9,-0.55);
  \node[vtx] at (0,0.55) {}; \node[vtx] at (0,-0.55) {};
  \node[font=\scriptsize,anchor=east] at (-0.9,0.55) {$t$}; \node[font=\scriptsize,anchor=west] at (0.9,0.55) {$W^+$};
  \node[font=\scriptsize,anchor=east] at (-0.9,-0.55) {$\bar s$}; \node[font=\scriptsize,anchor=west] at (0.9,-0.55) {$H$};
  \node[font=\scriptsize] at (0.22,0) {$s$};
  \node[slab] at (0,-1.2) {$u$ チャンネルの $s$};
\end{scope}
\end{tikzpicture}
\caption{手順 (2) で求めた木の過程の図と、振幅に掛かる CKM の因子。(a)、(b) は式~\eqref{eq:widthsckm} の崩壊、(c)、(d) は式~\eqref{eq:qqbar}、(e) は $u\bar u\to W^+W^-$ の5個の図で、$t$ チャンネルの3個は $|V_{uk}|^2$ を持ち、$d_k$ が質量0なら和が $V=1$ の1個の図に等しい。(f) は $t\bar s\to W^+H$ の4個の図で、どれも $V_{ts}^*$ を持つ。一般の $R_\xi$ ゲージでは $s$ チャンネルの $W$ の伝播関数の縦の部分と $G^+$ の図が $\xi_W$ を含み、和で打ち消す（$H$ の線が $s$ に付く $u$ チャンネルの図は $m_s$ に比例する）。矢印はフェルミオン数の流れで、左の端で外向きの矢印は入射する反粒子（$\bar d_j$、$\bar u$、$\bar s$）である。}
\label{fig:treeckm}
\end{figure}

\subsection{湯川行列の $\MSbar$ の繰り込みと CKM 行列の走り}\label{sec:rg}

\paragraph{行列の $\beta$ 関数。}湯川行列の裸の量を $Y_0=\mu^\epsilon(Y+\delta Y)$ と書く。SM ノートの式~(14) と同じく、1ループの $\MSbar$ では $\delta Y=\frac{\DUV}2\beta_Y$ で、$\beta_Y$ は4次元の1ループの $\beta$ 関数
\begin{equation}
\begin{aligned}
K\beta_{Y_u}&=\Bigl[\frac32\bigl(Y_uY_u^\dagger-Y_dY_d^\dagger\bigr)+\mathcal T-8g_s^2-\frac94g^2-\frac{17}{12}g'^2\Bigr]Y_u,\\
K\beta_{Y_d}&=\Bigl[\frac32\bigl(Y_dY_d^\dagger-Y_uY_u^\dagger\bigr)+\mathcal T-8g_s^2-\frac94g^2-\frac5{12}g'^2\Bigr]Y_d
\end{aligned}
\label{eq:betaY}
\end{equation}
である~\cite{mv2,arason92}（\ref{app:rg}。手順 (4) で本ノートの模型から導いた）。$\mathcal T=\tr(3Y_uY_u^\dagger+3Y_dY_d^\dagger+Y_\ell Y_\ell^\dagger)$ で、ゲージ結合と $\mathcal T$ の項は単位行列に比例する。場の繰り込みを世代のエルミートな行列にとったもので、$\MSbar$ の $\beta$ 関数なのでゲージによらない。$V=1$ では SM ノートの式~(15) の $\beta_{y_u}$、$\beta_{y_d}$ になる。ゲージ結合、$\lambda$、$\mu^2$ の1ループの $\beta$ 関数は湯川行列を $\tr Y_uY_u^\dagger=\sum y_u^2$、$\tr(Y_uY_u^\dagger)^2=\sum y_u^4$ などの形でだけ含むので $V$ によらず、SM ノートの式~(15) のままである（2ループでは $\tr(Y_uY_u^\dagger Y_dY_d^\dagger)=\sum_{ij}|V_{ij}|^2y_{u_i}^2y_{d_j}^2$ などが入る）。

\paragraph{質量と CKM 行列の走り。}$H_u=Y_uY_u^\dagger$、$H_d=Y_dY_d^\dagger$ は二重項の世代の空間のエルミートな行列で、上型の質量の基底では $H_u=D_u$、$H_d=VD_dV^\dagger$ である。式~\eqref{eq:betaY} から $H_u$、$H_d$ の走りを求め、固有値と固有ベクトルの摂動を使うと（\ref{app:rg}）
\begin{equation}
\begin{aligned}
K\frac{\dd\ln y_{u_i}}{\dd\ln\mu}&=\frac32y_{u_i}^2-\frac32\sum_k|V_{ik}|^2y_{d_k}^2+\mathcal T-8g_s^2-\frac94g^2-\frac{17}{12}g'^2,\\
K\frac{\dd\ln y_{d_j}}{\dd\ln\mu}&=\frac32y_{d_j}^2-\frac32\sum_k|V_{kj}|^2y_{u_k}^2+\mathcal T-8g_s^2-\frac94g^2-\frac5{12}g'^2
\end{aligned}
\label{eq:yrun}
\end{equation}
と
\begin{equation}
K\frac{\dd V_{ij}}{\dd\ln\mu}=\frac32\Biggl[\sum_{k\ne i}\frac{y_{u_i}^2+y_{u_k}^2}{y_{u_k}^2-y_{u_i}^2}\bigl(VD_dV^\dagger\bigr)_{ik}V_{kj}
-\sum_{k\ne j}V_{ik}\frac{y_{d_k}^2+y_{d_j}^2}{y_{d_j}^2-y_{d_k}^2}\bigl(V^\dagger D_uV\bigr)_{kj}\Biggr]
\label{eq:Vrun}
\end{equation}
を得る。ゲージ結合と $\mathcal T$ の項は湯川行列を単位行列の倍に変えるだけなので $V$ を回さず、$V$ の走りは $\frac32(H_u-H_d)$ の項だけから来る。式~\eqref{eq:Vrun} は固有ベクトルの位相を変えない約束（回転の生成子の対角の成分を0）によるもので、ほかの約束とは場の位相の付け替えだけ違い、$|V_{ij}|$ と $J$ の走りは同じである。$V=1$ では右辺が0になり、$V=1$ が保たれる。質量の階層 $y_t\gg y_c,y_u$、$y_b\gg y_s,y_d$ で主要な項をとると、
\begin{equation}
K\frac{\dd\ln|V_{ub}|}{\dd\ln\mu}\simeq K\frac{\dd\ln|V_{cb}|}{\dd\ln\mu}\simeq K\frac{\dd\ln|V_{td}|}{\dd\ln\mu}\simeq K\frac{\dd\ln|V_{ts}|}{\dd\ln\mu}\simeq\frac32\bigl(y_t^2+y_b^2\bigr),\qquad
K\frac{\dd\ln J}{\dd\ln\mu}\simeq3\bigl(y_t^2+y_b^2\bigr)
\label{eq:Vhier}
\end{equation}
で、ほかの成分の走りはこれより $|V_{cb}|^2$ 程度小さい。標準模型ではこれらの成分は高いスケールで大きくなる（式~\eqref{eq:betaY} の $-\frac32Y_dY_d^\dagger$ の符号による。超対称な模型では符号が逆になる~\cite{babu87,bhmp99}）。基準点では $y_t=0.9588$、$y_b=0.01662$ で、式~\eqref{eq:Vrun} から $K\,\dd\ln|V_{ub}|/\dd\ln\mu=1.37798$、$K\,\dd\ln|V_{cb}|/\dd\ln\mu=1.37786$、$K\,\dd\ln|V_{td}|/\dd\ln\mu=1.37926$、$K\,\dd\ln|V_{ts}|/\dd\ln\mu=1.37780$、$K\,\dd\ln J/\dd\ln\mu=2.7546$ で、式~\eqref{eq:Vhier} の $1.37935$、$2.7587$ とよく合う。$|V_{us}|$ の走りは $K\,\dd\ln|V_{us}|/\dd\ln\mu=1.55\times10^{-3}$ である。式~\eqref{eq:Vrun} は、式~\eqref{eq:betaY} を $Y_u$、$Y_d$ の行列のまま数値で積分し、各点で対角化して求めた $|V_{ij}|$ と $J$ の微分と、相対 $4\times10^{-8}$ 以内（中心差分の誤差）で一致した（湯川結合を大きくして全項が効くようにした点で。\texttt{Src/ckm\_rg\_handcheck.py}、この段落の数値もこのスクリプトの出力）。

\paragraph{CKM の対抗項。}$V$ は $Y_u$、$Y_d$ だけの関数で、$v$ によらない。裸の $V_0$ は裸の湯川行列 $Y_0=Y+\frac{\DUV}2\beta_Y$ から同じ回転で作るので、1ループでは
\begin{equation}
V_0=V+\delta V,\qquad \delta V=\frac{\DUV}2\,\frac{\dd V}{\dd\ln\mu}
\label{eq:dV}
\end{equation}
で、$\dd V/\dd\ln\mu$ は式~\eqref{eq:Vrun} である。これは SM ノートの式~(18) の質量の対抗項 $\delta M^2=\frac{\DUV}2\dd M^2/\dd\ln\mu$ と同じ形で、FJ の処方のタドポールは $v$ にしか入らないので $\delta V$ に入らない。$\delta V$ は $\MSbar$ の $\beta$ 関数だけで決まるので、ゲージによらない。場の位相の付け替えの自由度は、場の繰り込みの反エルミートな対角の部分と打ち消し合う。

\subsection{クォークの自己エネルギーの行列と CKM の対抗項}\label{sec:se}

\paragraph{世代の行列。}SM ノートの式~(24) を世代の行列に広げ、クォーク $q$（$u$ 型または $d$ 型）の1PI の2点関数の1ループの部分を、$q_j$ が入り $q_i$ が出るものとして
\begin{equation}
\Gamma^q_{ij}\supset i\bigl[\slashed pP_L\Sigma^L_{ij}(p^2)+\slashed pP_R\Sigma^R_{ij}(p^2)+\bigl(m_iP_L+m_jP_R\bigr)\Sigma^S_{ij}(p^2)\bigr]
\label{eq:Gammaq}
\end{equation}
と書く（Denner の形~\cite{denner93}。本ノートの模型の1ループでは、スカラーの部分がこの形になる）。$d$ 型のクォークでは、中性の場（$A$、$Z$、$G$、$H$、$G^0$）のループとタドポールの付いた図は世代について対角で、非対角の成分は $W^\pm$、$G^\pm$ と内線の $u_k$ のループ（図~\ref{fig:ckmdiag} (c)）だけから来て、$\sum_k\lambda^{(ij)}_kF(m_{u_k}^2)$ の形になる。$u$ 型では $V$ と $V^\dagger$ の役割を入れ替える。図の型を図~\ref{fig:seckm} に示す。

\begin{figure}[tbp]
\centering
\begin{tikzpicture}[scale=0.88]
\foreach \x/\sty/\inl/\arc/\lbl in {
    0/photon/u_k/{$W^-$}/{(a) $W^\pm$：$\lambda^{(ij)}_k$\\極は $m_{u_k}$ によらない},
    3.1/scalar/u_k/{$G^-$}/{(b) $G^\pm$：$\lambda^{(ij)}_k$\\極は $\lambda^{(ij)}_ky_{u_k}^2$},
    6.2/photon/d_i/{$A$、$Z$、$G$}/{(c) 中性のゲージ場\\対角（$i=j$）},
    9.3/scalar/d_i/{$H$、$G^0$}/{(d) $H$、$G^0$\\対角（$i=j$）}} {
  \begin{scope}[shift={(\x,0)}]
    \draw[fermion] (-1.1,0) -- (-0.55,0); \draw[fermion] (-0.55,0) -- (0.55,0); \draw[fermion] (0.55,0) -- (1.1,0);
    \draw[\sty] (-0.55,0) arc (180:0:0.55);
    \node[vtx] at (-0.55,0) {}; \node[vtx] at (0.55,0) {};
    \node[font=\scriptsize] at (0,-0.22) {$\inl$};
    \node[font=\scriptsize] at (-0.95,0.22) {$d_j$}; \node[font=\scriptsize] at (0.95,0.22) {$d_i$};
    \node[font=\scriptsize] at (0,0.82) {\arc};
    \node[slab] at (0,-0.85) {\lbl};
  \end{scope}}
% (e) tadpole
\begin{scope}[shift={(12.2,0)}]
  \draw[fermion] (-1.1,0) -- (0,0); \draw[fermion] (0,0) -- (1.1,0);
  \draw[scalar] (0,0) -- (0,0.55);
  \node[blob,minimum size=0.5cm] at (0,0.8) {$T$};
  \node[vtx] at (0,0) {};
  \node[font=\scriptsize] at (0.25,0.3) {$H$};
  \node[font=\scriptsize] at (-0.95,0.22) {$d_j$}; \node[font=\scriptsize] at (0.95,0.22) {$d_i$};
  \node[slab] at (0,-0.85) {(e) タドポール\\対角（$i=j$）};
\end{scope}
% (f) counterterm
\begin{scope}[shift={(15.0,0)}]
  \draw[fermion] (-1.1,0) -- (0,0); \draw[fermion] (0,0) -- (1.1,0);
  \node[ctv] at (0,0) {};
  \node[font=\scriptsize] at (-0.95,0.22) {$d_j$}; \node[font=\scriptsize] at (0.95,0.22) {$d_i$};
  \node[slab] at (0,-0.85) {(f) 対抗項\\$\delta Z^{L,R}_{ij}$、$\delta M_{ij}$};
\end{scope}
\end{tikzpicture}
\caption{下型のクォークの自己エネルギーの行列 $\Sigma_{ij}$（式~\eqref{eq:Gammaq}、$d_j$ が入り $d_i$ が出る）の図。世代を変えるのは $W^\pm$ と $G^\pm$ のループ (a)、(b) だけで、内線の $u_k$ について $\lambda^{(ij)}_k=V_{ki}^*V_{kj}$ を掛けて足す。(a) の紫外の極は $m_{u_k}$ によらず、ユニタリー性で非対角の成分が消える。非対角の紫外の極は (b) の $m_{u_k}^2$ に比例する部分だけで、$\xi_W$ によらない（式~\eqref{eq:uvoff}）。$\xi_W$ はどちらの有限部分にも入る。中性の場のループ (c)、(d) とタドポールの付いた図 (e) は対角で、SM ノートの図~4 のトップの (i)〜(n) と同じ形である。(f) は $\MSbar$ の対抗項で、非対角の成分は場の繰り込みと質量の対抗項 $\delta M_{ij}=-\frac34\frac{\DUV}K(V^\dagger D_uV)_{ij}m_{d_j}$（\ref{app:rg}）からなる。上型のクォークでは $V$ と $V^\dagger$、$u$ と $d$ を入れ替える。}
\label{fig:seckm}
\end{figure}

\paragraph{紫外の極。}$\MSbar$ のエルミートな場の繰り込みと $\delta Y$（\ref{app:rg}）から、非対角の成分（$i\ne j$）の紫外の極は
\begin{equation}
K\Sigma^L_{ij}\big|_{\mathrm{UV}}=\frac12X_{ij}\DUV,\qquad K\Sigma^R_{ij}\big|_{\mathrm{UV}}=0,\qquad K\Sigma^S_{ij}\big|_{\mathrm{UV}}=-X_{ij}\DUV,\qquad
X=\begin{cases}V^\dagger D_uV&(d\ \text{型})\\ VD_dV^\dagger&(u\ \text{型})\end{cases}
\label{eq:uvoff}
\end{equation}
になるはずである（手順 (3) で記号として確かめた）。$X_{ij}=\sum_k\lambda^{(ij)}_ky_{u_k}^2$（$d$ 型）は、$G^\pm$ のループの $m_{u_k}^2$ に比例する部分で、これは $\xi$ によらない（$\xi_W$ は $G^\pm$ の伝播関数の質量にだけ入る）。$W$ のループの紫外の極は内線の質量によらないので、ユニタリー性 $\sum_k\lambda^{(ij)}_k=\delta_{ij}$ で非対角の成分が0になる。$W$ の伝播関数の縦の部分 $(1-\xi_W)k_\mu k_\nu/[(k^2-M_W^2)(k^2-\xi_WM_W^2)]$ の極も $1/M_W^2$ を持たないので質量によらず、同じく打ち消す。したがって非対角の紫外の極はゲージによらず、$\xi$ に依存するのは有限部分だけである。$\Sigma^R$ に非対角の極がないのは、$d_R$ の場の繰り込み $Y_d^\dagger Y_d=D_d$ が質量の基底で対角だからである。対角の成分の極は、SM ノートの式~(26) の条件（$\Sigma^V+\Sigma^S$ の極が $\frac14\dd\ln m^2/\dd\ln\mu$）を式~\eqref{eq:yrun} で満たす（手順 (3) で確かめた）。

\paragraph{on-shell の対抗項との関係。}on-shell の処方では、外線の世代の混合を消す非対角の場の定数
\begin{equation}
\delta Z^L_{ij}=\frac2{m_i^2-m_j^2}\,\widetilde{\re}\bigl[m_j^2\Sigma^L_{ij}(m_j^2)+m_im_j\Sigma^R_{ij}(m_j^2)+(m_i^2+m_j^2)\Sigma^S_{ij}(m_j^2)\bigr]\qquad(i\ne j)
\label{eq:dZos}
\end{equation}
（$\widetilde{\re}$ はループの積分の実部だけをとる~\cite{denner93}）を使う。右巻きの場の定数は、同じ条件 $\hat\Sigma_{ij}(p)u_j(p)=0$（$p^2=m_j^2$）の右巻きの部分から
\begin{equation}
\delta Z^R_{ij}=\frac2{m_i^2-m_j^2}\,\widetilde{\re}\bigl[m_j^2\Sigma^R_{ij}(m_j^2)+m_im_j\Sigma^L_{ij}(m_j^2)+2m_im_j\Sigma^S_{ij}(m_j^2)\bigr]\qquad(i\ne j)
\label{eq:dZRos}
\end{equation}
で、$\Sigma^S$ の係数が $\delta Z^L$ の $m_i^2+m_j^2$ と違い $2m_im_j$ である（手順 (7)）。Denner と Sack~\cite{ds90} は CKM の対抗項を
\begin{equation}
\delta V^{\mathrm{DS}}=\frac14\Bigl[\bigl(\delta Z^{u,L}-\delta Z^{u,L\dagger}\bigr)V-V\bigl(\delta Z^{d,L}-\delta Z^{d,L\dagger}\bigr)\Bigr]
\label{eq:dVds}
\end{equation}
とした。式~\eqref{eq:uvoff} を式~\eqref{eq:dZos}、\eqref{eq:dVds} に入れると、$\delta V^{\mathrm{DS}}$ の紫外の極は式~\eqref{eq:dV} の $\MSbar$ の $\delta V$（式~\eqref{eq:Vrun} の位相の約束）にちょうど等しい（\ref{app:rg}）。違いは有限部分で、$\delta V^{\mathrm{DS}}$ の有限部分は自己エネルギーの非対角の有限部分を通じて $\xi_W$ に依存する~\cite{ggm99,yamada01}（手順 (3)、表~\ref{tab:ckmct}）。$\MSbar$ では、外線の世代の混合は、$\MSbar$ で繰り込んだ自己エネルギーから式~\eqref{eq:dZos} の形で作る有限の場の定数で扱い、それと頂点の補正の和は $\xi$ によらない（手順 (5) で確かめた）。$W^+\to u_i\bar d_j$ のどの図に世代の混合が入るかを図~\ref{fig:wud} に示す。

\begin{figure}[tbp]
\centering
\begin{tikzpicture}[scale=0.88]
% (a) neutral exchange between the outgoing quarks
\begin{scope}[shift={(0.6,0)}]
  \draw[photon] (-1.0,0) -- (0,0);
  \draw[fermion] (0,0) -- (0.6,0.42); \draw[fermion] (0.6,0.42) -- (1.2,0.84);
  \draw[fermion] (1.2,-0.84) -- (0.6,-0.42); \draw[fermion] (0.6,-0.42) -- (0,0);
  \draw[photon] (0.6,0.42) -- (0.6,-0.42);
  \node[vtx] at (0,0) {}; \node[vtx] at (0.6,0.42) {}; \node[vtx] at (0.6,-0.42) {};
  \node[font=\scriptsize] at (-0.6,0.3) {$W^+$};
  \node[font=\scriptsize,anchor=west] at (1.2,0.84) {$u_i$}; \node[font=\scriptsize,anchor=west] at (1.2,-0.84) {$\bar d_j$};
  \node[font=\scriptsize,anchor=west] at (0.72,0) {$\gamma$、$Z$、$G$};
  \node[slab] at (0.2,-1.35) {(a) 中性の場の交換\\（$H$、$G^0$ も）：$V_{ij}$};
\end{scope}
% (b) boson triangle
\begin{scope}[shift={(4.1,0)}]
  \draw[photon] (-1.0,0) -- (0,0);
  \draw[photon] (0,0) -- (0.75,0.5); \draw[photon] (0,0) -- (0.75,-0.5);
  \draw[fermion] (1.3,-0.85) -- (0.75,-0.5); \draw[fermion] (0.75,-0.5) -- (0.75,0.5); \draw[fermion] (0.75,0.5) -- (1.3,0.85);
  \node[vtx] at (0,0) {}; \node[vtx] at (0.75,0.5) {}; \node[vtx] at (0.75,-0.5) {};
  \node[font=\scriptsize] at (-0.6,0.3) {$W^+$};
  \node[font=\scriptsize,anchor=west] at (1.3,0.85) {$u_i$}; \node[font=\scriptsize,anchor=west] at (1.3,-0.85) {$\bar d_j$};
  \node[font=\scriptsize] at (1.0,0) {$d_j$};
  \node[font=\scriptsize] at (0.1,0.55) {$W^+$};
  \node[font=\scriptsize] at (0.1,-0.58) {$\gamma$、$Z$};
  \node[slab] at (0.2,-1.35) {(b) ボソンの三角\\（$G^+$、$H$ も）：$V_{ij}$};
\end{scope}
% (c) off-diagonal self-energy on the u leg
\begin{scope}[shift={(7.6,0)}]
  \draw[photon] (-1.0,0) -- (0,0);
  \draw[fermion] (0,0) -- (0.55,0.385); \draw[fermion] (0.75,0.525) -- (1.25,0.875);
  \draw[fermion] (1.25,-0.875) -- (0,0);
  \node[blob,minimum size=0.32cm] at (0.65,0.455) {};
  \node[vtx] at (0,0) {};
  \node[font=\scriptsize] at (-0.6,0.3) {$W^+$};
  \node[font=\scriptsize] at (0.12,0.42) {$u_k$};
  \node[font=\scriptsize,anchor=west] at (1.25,0.875) {$u_i$}; \node[font=\scriptsize,anchor=west] at (1.25,-0.875) {$\bar d_j$};
  \node[slab] at (0.2,-1.35) {(c) 外線の混合\\$\Sigma^u_{ik}V_{kj}$};
\end{scope}
% (d) off-diagonal self-energy on the dbar leg
\begin{scope}[shift={(11.0,0)}]
  \draw[photon] (-1.0,0) -- (0,0);
  \draw[fermion] (0,0) -- (1.25,0.875);
  \draw[fermion] (1.25,-0.875) -- (0.75,-0.525); \draw[fermion] (0.55,-0.385) -- (0,0);
  \node[blob,minimum size=0.32cm] at (0.65,-0.455) {};
  \node[vtx] at (0,0) {};
  \node[font=\scriptsize] at (-0.6,0.3) {$W^+$};
  \node[font=\scriptsize] at (0.12,-0.45) {$d_k$};
  \node[font=\scriptsize,anchor=west] at (1.25,0.875) {$u_i$}; \node[font=\scriptsize,anchor=west] at (1.25,-0.875) {$\bar d_j$};
  \node[slab] at (0.2,-1.35) {(d) 外線の混合\\$V_{ik}\Sigma^d_{kj}$};
\end{scope}
% (e) W self-energy on the external W
\begin{scope}[shift={(14.3,0)}]
  \draw[photon] (-1.2,0) -- (-0.75,0); \draw[photon] (-0.35,0) -- (0,0);
  \node[blob,minimum size=0.4cm] at (-0.55,0) {};
  \draw[fermion] (0,0) -- (0.85,0.6); \draw[fermion] (0.85,-0.6) -- (0,0);
  \node[vtx] at (0,0) {};
  \node[font=\scriptsize] at (-0.95,0.3) {$W^+$};
  \node[font=\scriptsize,anchor=west] at (0.85,0.6) {$u_i$}; \node[font=\scriptsize,anchor=west] at (0.85,-0.6) {$\bar d_j$};
  \node[slab] at (-0.1,-1.35) {(e) $\Sigma^{WW}$\\$V_{ij}$};
\end{scope}
% (f) counterterm
\begin{scope}[shift={(17.0,0)}]
  \draw[photon] (-1.0,0) -- (0,0);
  \draw[fermion] (0,0) -- (0.85,0.6); \draw[fermion] (0.85,-0.6) -- (0,0);
  \node[ctv] at (0,0) {};
  \node[font=\scriptsize] at (-0.6,0.3) {$W^+$};
  \node[font=\scriptsize,anchor=west] at (0.85,0.6) {$u_i$}; \node[font=\scriptsize,anchor=west] at (0.85,-0.6) {$\bar d_j$};
  \node[slab] at (0,-1.35) {(f) 対抗項\\$\delta V_{ij}$ と $\delta Z$};
\end{scope}
\end{tikzpicture}
\caption{手順 (5) で求める $W^+\to u_i\bar d_j$ の1ループの図の型。(a)、(b) の頂点の補正では、$W^+$ に付く荷電の頂点が一つで、残りは中性の対角な結合なので、$V_{ij}$ に比例する（荷電のボソンを2本含む頂点の補正はない）。(e) の $W$ の自己エネルギーも $V_{ij}$ に比例する。世代の混合が入るのは、外線の自己エネルギーの非対角の成分 (c)、(d)（図~\ref{fig:seckm} (a)、(b)）と CKM の対抗項 (f) だけで、(c) と (d) は内線の世代 $k$ について足す。$\MSbar$ の振幅では、(c)、(d) は $\MSbar$ で繰り込んだ自己エネルギーから作る有限の場の定数（式~\eqref{eq:dZos} の形）になり、それと (a)、(b)、(e) の $\xi$ 依存の和は $\xi$ によらない（手順 (5)）。$\delta V$ を $\delta V^{\mathrm{DS}}$ に替えると $\xi$ 依存が残る~\cite{ggm99,yamada01}（手順 (5)、表~\ref{tab:wud}）。}
\label{fig:wud}
\end{figure}

\subsection{世代を変える1ループの過程}\label{sec:fcnc}

$d_j\to d_i$（$i\ne j$）の中性の遷移は1ループで初めて現れ、振幅は
\begin{equation}
\mathcal A(d_j\to d_i)=\sum_k\lambda^{(ij)}_kF(x_k)=\sum_k\lambda^{(ij)}_k\bigl[F(x_k)-F(0)\bigr]
\label{eq:gim}
\end{equation}
の形になる（$x_k=m_{u_k}^2/M_W^2$。2つ目の等号は $\sum_k\lambda^{(ij)}_k=0$）。$F$ の定数の部分は打ち消し、質量の差だけが残る（GIM の機構~\cite{gim70}）。基準点の $b\to s$ では $\lambda_t=-0.04107+0.00076i$、$\lambda_c=0.04072$、$|\lambda_u|=8.4\times10^{-4}$ で、ほぼ $\lambda_t[F(x_t)-F(x_c)]$ になる。$\Delta F=2$ の箱の図（図~\ref{fig:ckmdiag} (e)）は二つの内線の和 $\sum_{k,l}\lambda_k\lambda_lS(x_k,x_l)$ になる。本ノートで求めるのは、(a) $b\to s\gamma$ の on-shell の双極子の振幅、(b) on-shell の $Z\to b\bar s$ の振幅（外線の世代の混合を含む）、(c) 外線の運動量を0にした $\Delta F=2$ の箱の図の和である。どれも個々の図は $\xi$ に依存し、和が $\xi$ によらないはずである。文献の Inami--Lim の関数~\cite{il81} は、本ノートの図から求めた値と比べる参照にとどめる。(c) から $K$、$B$ 中間子の質量差と $\epsilon_K$ を出すには、ハドロンの行列要素（格子 QCD）と QCD の補正の因子が要り、それらは文献の値を入力として明示して使う。それぞれの図の型を図~\ref{fig:fcncdiag} に示す。

\begin{figure}[tbp]
\centering
\begin{tikzpicture}[scale=0.85]
% (a), (b): photon from the internal u_k with a W or a G loop
\foreach \x/\sty/\arc/\lbl in {0/photon/{$W^-$}/{(a) $\gamma$ が $u_k$ に付く\\$W$ のループ}, 3.3/scalar/{$G^-$}/{(b) $\gamma$ が $u_k$ に付く\\$G^\pm$ のループ}} {
  \begin{scope}[shift={(\x,0)}]
    \draw[fermion] (-1.25,0) -- (-0.65,0); \draw[fermion] (-0.65,0) -- (0,0); \draw[fermion] (0,0) -- (0.65,0); \draw[fermion] (0.65,0) -- (1.25,0);
    \draw[\sty] (-0.65,0) arc (180:0:0.65);
    \draw[photon] (0,0) -- (0,-0.7);
    \node[vtx] at (-0.65,0) {}; \node[vtx] at (0.65,0) {}; \node[vtx] at (0,0) {};
    \node[font=\scriptsize] at (-1.1,0.22) {$b$}; \node[font=\scriptsize] at (1.1,0.22) {$s$};
    \node[font=\scriptsize] at (-0.36,-0.22) {$u_k$}; \node[font=\scriptsize] at (0,0.92) {\arc};
    \node[font=\scriptsize] at (0.25,-0.62) {$\gamma$};
    \node[slab] at (0,-1.25) {\lbl};
  \end{scope}}
% (c), (d): photon from the boson line
\foreach \x/\sty/\arc/\lbl in {6.6/photon/{$W^-$}/{(c) $\gamma W^+W^-$}, 9.9/scalar/{$G^-$}/{(d) $\gamma G^+G^-$}} {
  \begin{scope}[shift={(\x,0)}]
    \draw[fermion] (-1.25,0) -- (-0.65,0); \draw[fermion] (-0.65,0) -- (0.65,0); \draw[fermion] (0.65,0) -- (1.25,0);
    \draw[\sty] (-0.65,0) arc (180:0:0.65);
    \draw[photon] (0,0.65) -- (0,1.25);
    \node[vtx] at (-0.65,0) {}; \node[vtx] at (0.65,0) {}; \node[vtx] at (0,0.65) {};
    \node[font=\scriptsize] at (-1.1,0.22) {$b$}; \node[font=\scriptsize] at (1.1,0.22) {$s$};
    \node[font=\scriptsize] at (0,-0.22) {$u_k$}; \node[font=\scriptsize] at (-0.72,0.62) {\arc};
    \node[font=\scriptsize] at (0.25,1.15) {$\gamma$};
    \node[slab] at (0,-1.25) {\lbl};
  \end{scope}}
% (e): the gamma W G vertex
\begin{scope}[shift={(13.2,0)}]
  \draw[fermion] (-1.25,0) -- (-0.65,0); \draw[fermion] (-0.65,0) -- (0.65,0); \draw[fermion] (0.65,0) -- (1.25,0);
  \draw[photon] (-0.65,0) arc (180:90:0.65); \draw[scalar] (0,0.65) arc (90:0:0.65);
  \draw[photon] (0,0.65) -- (0,1.25);
  \node[vtx] at (-0.65,0) {}; \node[vtx] at (0.65,0) {}; \node[vtx] at (0,0.65) {};
  \node[font=\scriptsize] at (-1.1,0.22) {$b$}; \node[font=\scriptsize] at (1.1,0.22) {$s$};
  \node[font=\scriptsize] at (0,-0.22) {$u_k$};
  \node[font=\scriptsize] at (-0.75,0.6) {$W^-$}; \node[font=\scriptsize] at (0.75,0.6) {$G^-$};
  \node[font=\scriptsize] at (0.25,1.15) {$\gamma$};
  \node[slab] at (0,-1.25) {(e) $\gamma W^\mp G^\pm$\\（逆の組も）};
\end{scope}
% (f): photon from the external leg and the off-diagonal self-energy
\begin{scope}[shift={(0.3,-3.5)}]
  \draw[fermion] (-1.3,0) -- (-0.7,0); \draw[fermion] (-0.7,0) -- (0.05,0); \draw[fermion] (0.45,0) -- (1.1,0);
  \node[blob,minimum size=0.4cm] at (0.25,0) {};
  \draw[photon] (-0.7,0) -- (-0.7,-0.7);
  \node[vtx] at (-0.7,0) {};
  \node[font=\scriptsize] at (-1.15,0.22) {$b$}; \node[font=\scriptsize] at (-0.3,0.22) {$b$}; \node[font=\scriptsize] at (0.95,0.22) {$s$};
  \node[font=\scriptsize] at (0.25,0.4) {$\Sigma_{sb}$};
  \node[font=\scriptsize] at (-0.45,-0.62) {$\gamma$};
  \node[slab] at (0,-1.25) {(f) 外線の混合\\（$s$ の側に $\gamma$ も）};
\end{scope}
% (g) box with vertical bosons
\begin{scope}[shift={(4.4,-3.3)}]
  \draw[fermion] (-1.1,0.55) -- (-0.55,0.55); \draw[fermion] (-0.55,0.55) -- (0.55,0.55); \draw[fermion] (0.55,0.55) -- (1.1,0.55);
  \draw[fermion] (1.1,-0.55) -- (0.55,-0.55); \draw[fermion] (0.55,-0.55) -- (-0.55,-0.55); \draw[fermion] (-0.55,-0.55) -- (-1.1,-0.55);
  \draw[photon] (-0.55,0.55) -- (-0.55,-0.55); \draw[photon] (0.55,0.55) -- (0.55,-0.55);
  \node[vtx] at (-0.55,0.55) {}; \node[vtx] at (0.55,0.55) {}; \node[vtx] at (-0.55,-0.55) {}; \node[vtx] at (0.55,-0.55) {};
  \node[font=\scriptsize] at (-1.0,0.78) {$b$}; \node[font=\scriptsize] at (1.0,0.78) {$s$};
  \node[font=\scriptsize] at (-1.0,-0.78) {$s$}; \node[font=\scriptsize] at (1.0,-0.78) {$b$};
  \node[font=\scriptsize] at (0,0.78) {$u_k$}; \node[font=\scriptsize] at (0,-0.78) {$u_l$};
  \node[slab] at (0,-1.45) {(g) $\Delta F=2$ の箱\\縦の線が $W$ か $G^\pm$};
\end{scope}
% (h) box with horizontal bosons
\begin{scope}[shift={(8.2,-3.3)}]
  \draw[fermion] (-1.1,0.55) -- (-0.55,0.55); \draw[fermion] (-0.55,0.55) -- (-0.55,-0.55); \draw[fermion] (-0.55,-0.55) -- (-1.1,-0.55);
  \draw[fermion] (1.1,-0.55) -- (0.55,-0.55); \draw[fermion] (0.55,-0.55) -- (0.55,0.55); \draw[fermion] (0.55,0.55) -- (1.1,0.55);
  \draw[photon] (-0.55,0.55) -- (0.55,0.55); \draw[photon] (-0.55,-0.55) -- (0.55,-0.55);
  \node[vtx] at (-0.55,0.55) {}; \node[vtx] at (0.55,0.55) {}; \node[vtx] at (-0.55,-0.55) {}; \node[vtx] at (0.55,-0.55) {};
  \node[font=\scriptsize] at (-1.0,0.78) {$b$}; \node[font=\scriptsize] at (1.0,0.78) {$s$};
  \node[font=\scriptsize] at (-1.0,-0.78) {$s$}; \node[font=\scriptsize] at (1.0,-0.78) {$b$};
  \node[font=\scriptsize] at (-0.8,0) {$u_k$}; \node[font=\scriptsize] at (0.8,0) {$u_l$};
  \node[slab] at (0,-1.45) {(h) もう一つの箱\\横の線が $W$ か $G^\pm$};
\end{scope}
% (i) Z -> b sbar: Z coupling to the W loop and to u_k
\begin{scope}[shift={(12.2,-3.5)}]
  \draw[photon] (-1.2,0) -- (-0.55,0);
  \draw[fermion] (-0.55,0) -- (0.25,0.6); \draw[fermion] (0.25,0.6) -- (0.95,1.0);
  \draw[fermion] (0.95,-1.0) -- (0.25,-0.6); \draw[fermion] (0.25,-0.6) -- (-0.55,0);
  \draw[photon] (0.25,0.6) -- (0.25,-0.6);
  \node[vtx] at (-0.55,0) {}; \node[vtx] at (0.25,0.6) {}; \node[vtx] at (0.25,-0.6) {};
  \node[font=\scriptsize] at (-0.95,0.25) {$Z$};
  \node[font=\scriptsize,anchor=west] at (0.95,1.0) {$b$}; \node[font=\scriptsize,anchor=west] at (0.95,-1.0) {$\bar s$};
  \node[font=\scriptsize] at (-0.3,0.45) {$u_k$}; \node[font=\scriptsize] at (-0.3,-0.45) {$u_k$};
  \node[font=\scriptsize,anchor=west] at (0.32,0) {$W$、$G^\pm$};
  \node[slab] at (0.1,-1.45) {(i) $Z\to b\bar s$ の一つ\\（$Z$ が $W$、$G^\pm$ に付く図も）};
\end{scope}
\end{tikzpicture}
\caption{世代を変える1ループの過程の図（第~\ref{sec:fcnc} 節）。(a)〜(f) は $b\to s\gamma$ の図の型で、内線の $u_k$ について $\lambda^{(sb)}_k=V_{ks}^*V_{kb}$ を掛けて足す。$\gamma$ はループの $u_k$（(a)、(b)）、荷電のボソン（(c)、(d)）、$\gamma W^\mp G^\pm$ の頂点（(e)、$eM_W$ に比例）に付き、外線の $b$、$s$ に付く図 (f) には非対角の自己エネルギー（図~\ref{fig:seckm} (a)、(b)）が入る。$\xi_W$ は個々の図に入り、on-shell の振幅の和で打ち消すはずである（手順 (7)）。(g)、(h) は $B_s$ の混合（$b\bar s\to s\bar b$）の二つの位相の箱の図で、縦（(g)）または横（(h)）の2本が $(W,W)$、$(W,G^\pm)$、$(G^\pm,W)$、$(G^\pm,G^\pm)$ の4通りあり、$\sum_{k,l}\lambda_k\lambda_l$ を掛ける（手順 (8)）。(i) は $Z\to b\bar s$ の図の一つで、$Z$ は内線の $u_k$ のほか、$W$ と $G^\pm$ にも付き、外線の混合の図も加わる（全体は図~\ref{fig:zbsdiag}）。$K$ と $B_d$ では $s,b$ を $d,s$、$d,b$ に替える。}
\label{fig:fcncdiag}
\end{figure}

\section{検算の計画}\label{sec:plan}

第~\ref{sec:results} 節の各量を、どの手段で計算して突き合わせるかを表~\ref{tab:plan} に示す。手計算（付録）を基準とし、ツールの出力がそれと一致するかを確かめる。ツールの版と配置は計算環境の記録~\cite{env} にあり、SM ノートと同じである。一般の $R_\xi$ ゲージの計算は、本ノートの模型ファイルから書き出した FeynArts の模型（\texttt{GaugeXi} を残す）と FeynCalc で行う。FeynArts の \texttt{SM.mod} は CKM 行列と Denner--Sack の対抗項を持つので、Feynman ゲージではそれと FormCalc とも、外線の因子で規約を直して比べる。

\begin{table}[tbp]
\centering\small
\caption{各量と、それを計算して突き合わせる手段。「—」はその手段を使わないことを表す。}
\label{tab:plan}
\begin{tabularx}{\linewidth}{>{\raggedright\arraybackslash}p{3.4cm}>{\raggedright\arraybackslash}p{1.3cm}>{\raggedright\arraybackslash}X>{\raggedright\arraybackslash}p{1.9cm}>{\raggedright\arraybackslash}p{2.1cm}}
\toprule
\textbf{量} & \textbf{手計算} & \textbf{FeynCalc + FeynArts、FormCalc} & \textbf{数値} & \textbf{MadGraph} \\
\midrule
Feynman 則（表~\ref{tab:rules-ckm}） & \ref{app:rules} & 自前の模型の頂点と表の差。ユニタリー性の規則、$V=1$ で SM ノートの頂点 & — & UFO と図の数 \\
木の幅と断面積 \eqref{eq:widthsckm} & \ref{app:tree} & 振幅の2乗を記号として & 陽な行列 & 幅と断面積 \\
湯川行列の $\beta$ 関数 \eqref{eq:betaY}、$V$ の走り \eqref{eq:Vrun} & \ref{app:rg} & 2点関数の極。一般の公式の陽な実装 & 行列のまま積分 & — \\
自己エネルギーの行列 \eqref{eq:uvoff}、$\delta V$ \eqref{eq:dV}、$\delta V^{\mathrm{DS}}$ \eqref{eq:dVds} & \ref{app:rg} & 一般の $R_\xi$ で。\texttt{SM.mod}（CKM） & LoopTools & — \\
$W\to u_i\bar d_j$、$|V_{ij}|$ と $V(\mu)$ & 手順 (5)、(6) & 一般の $R_\xi$。FormCalc & 数値 & MadLoop \\
$b\to s\gamma$、$Z\to b\bar s$、$\Delta F=2$ \eqref{eq:gim} & 手順 (7)、(8) & 一般の $R_\xi$ & LoopTools、Package-X & MadLoop \\
$V$ と $J$ の高いスケールまでの走り & 手順 (10) & — & 常微分方程式 & — \\
\bottomrule
\end{tabularx}
\end{table}

\paragraph{一致の基準。}SM ノートと同じである。発散部分と対抗項は $1/\epsilon$ の係数が記号として完全に一致することを、有限部分は基準点で倍精度の丸め誤差の範囲（相対 $10^{-12}$ 程度）で一致することを求める。$\xi$ によらないはずの量は、記号として $\xi$ を含まないこと、または $\xi=0,1,3$ で同じ値になることを求める。$V$ を含む式は、$V$ を一般の複素行列の記号のまま比べるか、標準のパラメータ表示 \eqref{eq:Vstd} を入れてユニタリー性を使って比べる。応用の手順では、比べる値の精度の範囲で一致することを求め、その精度を各手順に示す。

\paragraph{手順。}次の順に進める。(1) 模型ファイル \texttt{SM\_CKM.fr} を SM ノートの \texttt{SM.fr} から作り、FeynArts の模型と UFO を書き出す。93個の頂点を $V$ を一般の複素行列とした表と比べ、ユニタリー性の規則を標準のパラメータ表示で確かめ、$V=1$ で SM ノートの頂点と比べる。MadGraph の図の数を、$V$ の非対角の成分を要する過程を含めて手で数えた数と比べる。(2) 木の振幅の2乗を FeynCalc で記号として確かめ、式~\eqref{eq:widthsckm} と、ユニタリー性による和（$\sum_j\Gamma(t\to d_jW)$、$W$ のハドロンへの幅、$u\bar u\to W^+W^-$）を確かめ、MadGraph の幅と断面積と比べる。(3) クォークの自己エネルギーを世代の行列として一般の $R_\xi$ ゲージで $D$ 次元で求め（図~\ref{fig:seckm}）、非対角の紫外の極が式~\eqref{eq:uvoff} で $\xi$ によらないこと、$V=1$ で SM ノートの値になること、有限部分の $\xi$ 依存を確かめる。$\delta V^{\mathrm{DS}}$ の紫外の極が $\MSbar$ の $\delta V$ に等しいことと、その有限部分が $\xi_W$ に依存することを求める。$V$ を厳密に入れた $W$ とトップの極質量と $\hat\Delta$ の変化を、SM ノートの見積もり（\texttt{ckm\_estimate.wl}）と比べる。(4) 湯川行列の $\beta$ 関数 \eqref{eq:betaY} を、2点関数（と $\bar u dG^+$ の頂点）の極と、一般の公式~\cite{mv1,mv2,luo03} を行列のまま組んだ実装の二つで確かめ、式~\eqref{eq:yrun}、\eqref{eq:Vrun} を導く。(5) $W^+\to u_i\bar d_j$ の1ループの電弱の振幅（図~\ref{fig:wud}）を、$\MSbar$ の $V$ と外線の世代の混合の有限な場の定数で一般の $R_\xi$ ゲージで求め、紫外で有限で $\xi$ によらないことを確かめる。$\delta V^{\mathrm{DS}}$ を使うと $\xi$ に依存することを示し、Feynman ゲージで FormCalc（\texttt{SM.mod} の Denner--Sack の対抗項）と比べる。(6) 外線の運動量を0にした $d_j\to u_ie\bar\nu_e$ の1ループの4フェルミの振幅を $\mu$ 崩壊（SM ノートの手順 (6)）と比べ、セミレプトニックな崩壊で測る $|V_{ij}|$ と $\MSbar$ の $V(\mu)$ の関係を求め、PDG の $|V_{ij}|$ から $V(M_Z)$、$V(M_t)$ を決める。(7) $b\to s\gamma$ の on-shell の双極子の振幅と on-shell の $Z\to b\bar s$ の振幅を一般の $R_\xi$ ゲージで求め（図~\ref{fig:fcncdiag} (a)〜(f)、(i)）、$\xi$ によらないこと、GIM の機構 \eqref{eq:gim} を確かめ、Inami--Lim の関数~\cite{il81} と比べる。(8) $\Delta F=2$ の箱の図（図~\ref{fig:fcncdiag} (g)、(h)）の和を外線の運動量を0にして一般の $R_\xi$ ゲージで求め、GIM の後で $\xi$ によらないことを確かめる。ハドロンの入力（明示した文献の値）で $B_d$、$B_s$ の質量差と $\epsilon_K$ を求め、測定値と比べる。(9) 本ノートの UFO と NLOCT の R2 の頂点で、MadLoop の世代を変えるループの振幅を手順 (7) の値と比べる。SM ノートの FJ の対抗項の UFO に CKM の対抗項を加え、$W\to u_i\bar d_j$ を手順 (5) と比べる。(10) 手順 (4) の本ノートの方程式で $V$ と $J$ を $M_{\mathrm{Pl}}$ まで走らせ、文献の値と比べる（文献の値は参照にとどめる）。

\section{検算}\label{sec:checks}

第~\ref{sec:plan} 節の手順を順に行う。スクリプトはすべて \texttt{Src/} にあり、手順 (1) は \texttt{run\_model\_check.sh}（約1分）、手順 (2) は \texttt{run\_tree\_check.sh}（約6分）、手順 (3) は \texttt{run\_oneloop\_check.sh}（約20分）、手順 (4) は \texttt{run\_rg\_check.sh}（約10分）、手順 (5) は \texttt{run\_wud\_check.sh}（約10分）、手順 (6) は \texttt{run\_sl\_check.sh}（約12分）、手順 (7) は \texttt{run\_fcnc\_check.sh}（約4分）でまとめて再実行できる。

\paragraph{手順 (1)：模型ファイル。}模型ファイル \texttt{Src/SM\_CKM.fr} は SM ノートの \texttt{Src/SM.fr} を写し、三か所だけを変えた。外部パラメータ $s_{12}$、$s_{23}$、$s_{13}$、$\delta$（ブロック \texttt{CKMBLOCK}）と、内部パラメータの複素行列 \texttt{CKM}（\texttt{Unitary -> True}、値は式~\eqref{eq:Vstd}）を加えた。左巻きの二重項の下の成分を $(Vd_L)_f=\texttt{CKM[f, f2]}\,d_{L,f2}$ にした。湯川の項を $-\bar Q_L(V\hat y_d)\Phi d_R$ にした（$\hat y_d$ は SM ノートと同じ対角の行列で、どの世代の添字も2回だけ現れる）。運動項と湯川の項を添字について展開した後、和をとる添字を持つ積 $V_{ab}V_{ac}^*$、$V_{ba}V_{ca}^*$ を $\delta_{bc}$ に置き換える規則を加えた。

スクリプト \texttt{generate\_SM\_CKM.wl} は四つのことを確かめた（出力の (a)〜(d)）。(a) FeynRules のチェックでは、エルミート性が通り、警告は SM ノートと同じ二つだけで、頂点のクラスは93個だった。(b) 93個のクラスを、表~\ref{tab:rules-ckm} と SM ノートの表~2〜5 の式と比べた。フェルミオンの頂点は世代の9個の成分ごとに比べ、$V$ は一般の複素行列の記号 \texttt{CKM[i, j]} とその複素共役のまま残した。荷電のクォークの頂点は $V_{ij}$、$V_{ij}^*$ を正しい添字の順で含み、ほかの頂点は対角で、93個すべての差が記号として0になった。$\bar u_id_jG^+$ の頂点は FeynRules の出力では和をとる世代の添字（$V_{ia}(\hat y_d)_{aj}$）を含み、それを項ごとに足してから比べた。比べ方が誤りを見逃さないことを確かめるため、SM ノートの3か所の誤り（$AW\bar W$ の符号、$G^0G^0H$ の2倍、$\bar u^Zu^ZH$ の $\xi_Z$ の欠落）に CKM の2か所（$\bar udW$ で $V$ を転置、$\bar duG^-$ で複素共役の欠落）を加えた表とも比べ、ちょうど5個の差が出た。(c) ユニタリー性の規則を確かめるため、規則を使わないラグランジアンからも頂点を作った。規則を使わないラグランジアンには $V_{ab}V_{ac}^*$ の形の項が12個あり、使ったものには0個である。標準のパラメータ表示を記号の角度で入れると（$V^\dagger V=1$ も記号として確かめた）、23個のフェルミオンの頂点の9個の成分が両方で一致した。規則を使わないラグランジアンでも中性の頂点は CKM 行列を含まず、FeynRules は \texttt{Unitary -> True} の行列のユニタリー性を自分で使っている。規則は冗長だが、誤りでもない。(d) $V=1$ を入れた93個のクラスが、SM ノートの FeynRules の出力（\texttt{models/sm/Src/generated/FeynmanRules\_SM.m}）と、フェルミオンでは成分ごとにすべて一致した。UFO の書き出しでは、93個のクラスから 163 個の世代を展開した頂点ができた。SM ノートの UFO の 139 個に、$\bar udW$、$\bar du\bar W$、$\bar udG^+$、$\bar duG^-$ の非対角の6個ずつを足した数である。FeynArts の模型では、$V$ は \texttt{CKM[i, j]} と \texttt{Conjugate[CKM[i, j]]} として、$G^\pm$ の結合では世代の和 \texttt{IndexSum} の中に現れる。

この過程で FeynRules の問題を一つ見つけた。SM ノートと同じ順に FeynArts の模型を書き出した後で UFO を書き出すと、UFO のための Feynman 則の計算が \texttt{Inner::incom} のメッセージを出し、94個のクラスを作った。94個目は電荷を保存しない $\bar duW$（電荷 $+2$）で、係数は添字が1個の \texttt{CKM[Index[Generation, 2]]} と \texttt{Index[Generation, \{\}]} を含む壊れた式だった。同じ計算を新しいカーネルで行うと93個でメッセージも出ず、FeynArts の模型を書き出すことだけが原因である。FeynArts の書き出しが残す状態が、二つの添字を持つ複素行列の世代の展開を壊すと見られるが、原因は特定していない。CKM 行列のない SM ノートの模型では起こらない。UFO を先に書き出すと、どちらも93個になり、FeynArts の模型のファイルは書き出しの順によらず同じだった（\texttt{diff} で確かめた。FeynArts の書き出しの並列のサブカーネルが \texttt{Simplify::time} を出すが、FeynRules の中の \texttt{Simplify} の時間切れで、結果は変わらない）。スクリプトはこの順にし、UFO の記録ファイルからクラスの数と頂点の数を読んで表示する。

UFO の制限カード \texttt{restrict\_massless.dat}（トップ以外のフェルミオンを質量0にし、CKM のパラメータは値を残す）は SM ノートのスクリプト \texttt{make\_restrict.py} で作った。MadGraph 3.8.1 は、その UFO から表~\ref{tab:diagrams} の数の木の図を、ユニタリーゲージと Feynman ゲージの両方で生成し、35個の過程すべてで手で数えた数と一致した。SM ノートの21個の過程は $u\bar u\to W^+W^-$ を除いて SM ノートと同じ数で、$u\bar u\to W^+W^-$ は $t$ チャンネルに $s$、$b$ が加わって3個から5個になった。残りの14個は $V$ の非対角の成分を要する過程で、$V=1$ では図がないか（$u\bar s\to W^+H$、$t\to sW^+$ など）、図が少ない（$t\bar t\to W^+W^-$ は $V=1$ では4個）。Feynman ゲージで図が増えるのは、$t\bar s\to W^+H$ の $s$ チャンネルの $G^+$（$\bar t sG^+$ の結合は $V_{ts}m_t$）だけである。世代を変える中性の過程 $d\bar s\to\mu^+\mu^-$、$s\bar b\to Z$、$b\bar s\to\gamma g$、$u\bar c\to e^+e^-$、$c\bar u\to ZH$ には、MadGraph は図を作らなかった（\texttt{NoDiagramException}。同じ設定で $u\bar s\to e^+\nu_e$ には1個の図を作った）。木のレベルに世代を変える中性の結合がないことの確かめである。MadGraph は図のない過程でスクリプトを止めるので、これらは過程ごとに別に実行した。スクリプトは \texttt{generate\_SM\_CKM.wl}（FeynRules）、\texttt{model\_mg5.mg5}（MadGraph）、\texttt{run\_model\_check.sh}（まとめと図のない過程）である。分かったツールの問題は計算環境の記録~\cite{env} の表 17 にまとめた。

\begin{table}[tbp]
\centering\small
\caption{MadGraph 3.8.1 が本ノートの UFO（トップ以外のフェルミオンは質量0、CKM のパラメータは式~\eqref{eq:bench}）から生成した木の図の数（ユニタリーゲージ / Feynman ゲージ）のうち、$V$ の非対角の成分を要する過程と、その内訳。SM ノートの表~7 の21個の過程も同じ実行で確かめ、$u\bar u\to W^+W^-$（下の表）を除いて SM ノートと同じ数だった。どれも手で数えた数と一致した。内訳の括弧は振幅に掛かる CKM の因子である。}
\label{tab:diagrams}
\begin{tabularx}{\linewidth}{>{\raggedright\arraybackslash}p{5.6cm}>{\raggedright\arraybackslash}p{1.6cm}X}
\toprule
過程 & 図の数 & 内訳 \\
\midrule
$u\bar u\to W^+W^-$、$d\bar d\to W^+W^-$ & 5 / 5 & $s$ チャンネルの $\gamma$、$Z$、$t$ チャンネルの $d,s,b$（$u,c,t$） \\
$t\bar t\to W^+W^-$ & 6 / 6 & $s$ の $\gamma$、$Z$、$H$、$t$ の $d,s,b$ \\
$u\bar s\to W^+H$ & 1 / 1 & $s$ の $W$（$V_{us}^*$） \\
$t\bar s\to W^+H$ & 2 / 3 & $s$ の $W$、$t$ の $t$（Feynman ゲージでは $s$ の $G^+$ が加わる） \\
$t\to sW^+$、$t\to dW^+$、$\bar t\to\bar bW^-$ & 1 / 1 & （$V_{ts}^*$、$V_{td}^*$、$V_{tb}$） \\
$u\bar s\to e^+\nu_e$、$u\bar b\to e^+\nu_e$、$c\bar d\to e^+\nu_e$ & 1 / 1 & $s$ の $W$ \\
$u\bar d\to c\bar s$、$u\bar s\to c\bar d$ & 1 / 1 & $s$ の $W$（$V_{ud}^*V_{cs}$、$V_{us}^*V_{cd}$） \\
$u\bar u\to d\bar s$、$u\bar c\to d\bar s$ & 1 / 1 & $t$ の $W$（$V_{ud}^*V_{us}$、$V_{ud}^*V_{cs}$） \\
\midrule
$d\bar s\to\mu^+\mu^-$、$s\bar b\to Z$、$b\bar s\to\gamma g$、$u\bar c\to e^+e^-$、$c\bar u\to ZH$ & 0 & 世代を変える中性の結合がない \\
\bottomrule
\end{tabularx}
\end{table}

\paragraph{手順 (2)：木の振幅の2乗と崩壊幅。}スクリプト \texttt{tree\_ckm\_feyncalc.wl} は、本ノートの FeynArts の模型で13個の振幅の2乗を FeynCalc で求め、手計算の式とすべて記号として一致させた。$V_{ij}$ は実の記号で $v^R_{ij}+iv^I_{ij}$ と書いた。FeynCalc の \texttt{ComplexConjugate} は記号を実数として扱い、明示的な $i$ だけを反転するので、こう書くと $|V_{ij}|^2=(v^R_{ij})^2+(v^I_{ij})^2$ が正しく出る。内訳は次の通りである（図は図~\ref{fig:treeckm}）。(i) $W^+\to u_i\bar d_j$ の4個（$u\bar d$、$u\bar s$、$c\bar b$ と、閾値より上の質量を入れた $t\bar d$）が式~\eqref{eq:widthsckm} の $|V_{ij}|^2$ 倍の式に、$t\to d_jW^+$ の3個が $|V_{tj}|^2$ 倍の式に一致した。(ii) $\bar t\to\bar bW^-$ の2乗は $t\to bW^+$ と等しく、木のレベルで CP が破れないことを確かめた。(iii) 質量0の $u_i\bar d_j\to\nu_ee^+$ の3個（$u\bar d$、$u\bar s$、$c\bar b$）と $u\bar u\to d\bar s$ が式~\eqref{eq:qqbar} に一致し、$\xi$ を含まなかった。(iv) $u\bar u\to W^+W^-$ の2乗は $\xi_A,\xi_Z,\xi_W$ によらず、$V_{uk}$ の6個の実の記号を含む。そこに標準のパラメータ表示を記号の角度で入れると、$V=1$ の値に記号として等しくなった（$\sum_k|V_{uk}|^2=1$）。(v) トップと $s$ の質量を残した $t\bar s\to W^+H$（$s$ チャンネルの $W$ と $G^+$、$t$ チャンネルのトップ、$u$ チャンネルの $s$。FeynArts は世代を変える $H$ の結合を持つ図も含めて8個を作り、その4個は対角の湯川結合で0になる）は、振幅が $\xi_W$ を含むが、2乗は $\xi_W$ によらず、$|V_{ts}|^2$ で割ると CKM の記号を含まなかった。$G^+$ の図と $W$ の伝播関数の $\xi_W$ の部分が打ち消すには、$G^+$ の頂点が $W$ の頂点と同じ $V_{ts}^*$ を持つ必要がある。$W$ と $G^\pm$ の頂点の CKM 行列の整合を、ゲージ不変性で確かめたことになる。

最初の実行では $u\bar u\to d\bar s$ だけが合わず、差は $9t^2(\cdots)/(t-\xi_WM_W^2)^2$ の形で、質量0では消えるはずの $t$ チャンネルの $G^+$ の図が残った。FeynArts の模型は $G^\pm$ の結合を世代の和（\texttt{IndexSum}）で書くが、FeynCalc の \texttt{FCFAConvert} は既定でこれを外し、和の添字を縮約されない添字（\texttt{CKM[1, Ind273]}、\texttt{yd[Ind273, 2]}）として残していた。対角の湯川結合の置き換えが効かず、$G^+$ の結合が0にならなかった。\texttt{DropIndexSum -> False} にして和を書き出して直し、振幅に \texttt{CKM} と \texttt{IndexSum} が残らないことを表示するようにした。

スクリプト \texttt{tree\_ckm\_numeric.py} は、SM ノートの手順 (2) の Dirac 行列、スピノル、偏極ベクトル（\texttt{models/sm/Src/tree\_numeric.py}）を読み込み、FeynCalc を使わずに $\sum|\mathcal M|^2$ を求めた。$t\bar s\to W^+H$ はユニタリーゲージ（$G^+$ の図がない）で計算した。FeynCalc の一般の $R_\xi$ の値と比べると、12個の点で相対 $6\times10^{-15}$ 以内で一致し、式~\eqref{eq:qqbar} の閉じた形とも $10^{-15}$ で一致した。ユニタリー性による和は、質量の分だけずれた。$W$ のハドロンへの幅は $V=1$ の値と相対 $-1.7\times10^{-6}$、$\sum_j\Gamma(t\to d_jW^+)$ は $V=1$ の $\Gamma(t\to bW^+)$ と $+2.0\times10^{-6}$ 違った。どちらも、$|V|^2$ の重みが質量の違う終状態に移るためである。

MadGraph 3.8.1 の値を表~\ref{tab:tree} に示す。12個の部分幅（$W^+$ の9個とトップの3個）は、本ノートの式と相対 $10^{-15}$ で一致した。ただし MadGraph は2体の幅を FeynRules が UFO に書いた解析的な式（\texttt{decays.py}）で計算するので、この一致は FeynRules がこの模型の頂点から作った式の確かめである。MadGraph は、色のある粒子の部分幅が QCD のスケールより小さいと0にする。そのため $t\to dW^+$、$t\to sW^+$ は幅のファイルに現れず、トップの全幅は $t\to bW^+$ だけの値になった。二つの値はログに16桁で表示されるので、そこから読んで比べた。断面積は、固定したエネルギーのパートンのビーム（$\sqrt s=200$~GeV、カットなし）で、統計誤差の0.7〜0.9倍以内で一致した。$u\bar u\to W^+W^-$（$\sqrt s=500$~GeV）は、本ノートの UFO と SM ノートの UFO（$V=1$）で表示の6桁まで同じ値になった。二つの実行は同じ乱数を使うので、被積分関数が等しいことの確かめである。手順 (2) で分かったツールの問題も、計算環境の記録~\cite{env} の表 17 に加えた。

\begin{table}[tbp]
\centering\small
\caption{手順 (2) の部分幅と断面積（基準点）。本ノートの値は式~\eqref{eq:widthsckm}、\eqref{eq:qqbar}（\texttt{tree\_ckm\_numeric.py}）、MadGraph は本ノートの UFO（幅は全質量、断面積は \texttt{restrict\_massless}、$\Gamma_W=2.085$~GeV は $s$ チャンネルだけ）。差は部分幅では相対差、断面積では MadGraph の統計誤差で割った差。$t\to dW^+$、$t\to sW^+$ の MadGraph の値は、MadGraph が幅のファイルで0にしてログに表示した値である。$W^+\to\ell^+\nu$ の3個（相対 $4\times10^{-16}$ 以内）は省いた。最後の行の本ノートの欄は、SM ノートの UFO（$V=1$）の MadGraph の値である。}
\label{tab:tree}
\begin{tabular}{lllr}
\toprule
量 & 本ノート & MadGraph & 差 \\
\midrule
$\Gamma(W^+\to u\bar d)$ [GeV] & 0.6530458862 & 0.6530458862 & $4\times10^{-16}$ \\
$\Gamma(W^+\to u\bar s)$ & 0.03482661797 & 0.03482661797 & $4\times10^{-16}$ \\
$\Gamma(W^+\to u\bar b)$ & $9.562333873\times10^{-6}$ & $9.562333873\times10^{-6}$ & $2\times10^{-16}$ \\
$\Gamma(W^+\to c\bar d)$ & 0.03478242666 & 0.03478242666 & $4\times10^{-16}$ \\
$\Gamma(W^+\to c\bar s)$ & 0.6518339581 & 0.6518339581 & $4\times10^{-16}$ \\
$\Gamma(W^+\to c\bar b)$ & $1.201188694\times10^{-3}$ & $1.201188694\times10^{-3}$ & $4\times10^{-16}$ \\
$\Gamma(t\to dW^+)$ & $9.560738356\times10^{-5}$ & $9.560738356\times10^{-5}$ & $9\times10^{-16}$ \\
$\Gamma(t\to sW^+)$ & $2.193395846\times10^{-3}$ & $2.193395846\times10^{-3}$ & $7\times10^{-16}$ \\
$\Gamma(t\to bW^+)$ & 1.294155923 & 1.294155923 & $-3\times10^{-16}$ \\
\midrule
$\sigma(u\bar d\to e^+\nu_e)$ [pb] & 1.3520376 & $1.3529\pm0.0009$ & $+0.91$ \\
$\sigma(u\bar s\to e^+\nu_e)$ & 0.072103555 & $0.072151\pm0.000050$ & $+0.94$ \\
$\sigma(u\bar u\to d\bar s)$ & 1.5884582 & $1.5893\pm0.0012$ & $+0.71$ \\
$\sigma(u\bar u\to W^+W^-)$、$\sqrt s=500$~GeV & 2.45025 & $2.45025\pm0.00667$ & 0 \\
\bottomrule
\end{tabular}
\end{table}

\paragraph{手順 (3)：クォークの自己エネルギーの行列と CKM の対抗項。}スクリプト \texttt{oneloop\_ckm\_feyncalc.wl} は、SM ノートの手順 (3) と同じ処方（FeynArts の模型の写しで光子とそのゴーストに質量 $M_A$ を与え、タドポールの付いた図は $H$ の線のものだけを残す）で、下型と上型のクォークの自己エネルギーを $3\times3$ の世代の行列として、一般の $R_\xi$ ゲージで $D$ 次元で求めた（18個の成分、図~\ref{fig:seckm}）。対角の成分の1PI の図は15個（内線の3世代の $W^\pm$、$G^\pm$ の6個、$A$、$Z$、$G$ の3個、$H$、$G^0$ の6個。そのうち世代を変える $H$、$G^0$ の4個は0）、非対角の成分は12個（$W^\pm$、$G^\pm$ の6個と、0になる $H$、$G^0$ の6個）で、タドポールの付いた図は対角の成分にだけ残る。式~\eqref{eq:Gammaq} の関数は、TID と \texttt{DiracSimplify} の後で $\slashed pP_L$、$\slashed pP_R$、$P_L$、$P_R$ の係数として読み取った。$\gamma_5$ を含むトレースは使わない。結果は次の通りである。(i) 18個の成分の1PI とタドポールの部分の36個すべてで、$P_L$ と $P_R$ のスカラーの部分が一つの関数になり、式~\eqref{eq:Gammaq} の Denner の形が成り立った。(ii) 非対角の12個の成分の紫外の極は、標準のパラメータ表示を記号の角度で入れると、式~\eqref{eq:uvoff} と記号として一致した。$\Sigma^L_{ij}$ の極は、ユニタリー性を使う前には $\xi_W$ を含み（$W$ の伝播関数の縦の部分の、$\sum_k\lambda^{(ij)}_k$ に比例する項）、使った後には含まない。(iii) 対角の6個のクォークで、タドポールを含めた $\Sigma^V+\Sigma^S$ の $p^2=m_i^2$ での極が、$|V_{ik}|^2y_{d_k}^2$ の重みを含めて式~\eqref{eq:yrun} の $\frac14\dd\ln m_i^2/\dd\ln\mu$ に等しかった。(iv) $V=1$ のトップの $\Sigma^V$、$\Sigma^S$（1PI とタドポールの部分）は、SM ノートの値と記号として一致した。(v) CKM 行列を含む $\Sigma^{WW}_T(M_W^2)$ の極（タドポールを含む）は、ユニタリー性で $V=1$ の値と同じになり、SM ノートの値とも一致した。独立な確かめとして、スクリプト \texttt{oneloop\_ckm\_smmod.wl} で FeynArts に付属の \texttt{SMQCD.mod}（Denner の規約）を \texttt{\$CKM = True} にして同じ18個の1PI の成分を求めた。光子の符号、$G^\pm$ の位相、クォークとグルーオンの頂点の符号の違いは自己エネルギーでは打ち消すので、4個の関数の72個すべてが一般の $R_\xi$ ゲージで記号として一致した。\texttt{SMQCD.mod} は湯川結合を対角に作るので、非対角の成分の図は6個、対角の成分は11個である。

有限部分の数値（\texttt{oneloop\_ckm\_numeric.wl}、LoopTools、基準点、$\mu=M_Z$）を表~\ref{tab:ckmct} に示す。非対角の成分の有限部分は $\xi_W$ で大きく変わり、$K\Sigma^L_{sb}(m_b^2)$ は $\xi_W=\frac12$ から 10 で $6.0\times10^{-3}$ から $9.7\times10^{-2}$ になる。それから式~\eqref{eq:dZos}、\eqref{eq:dVds} で作った Denner--Sack の対抗項は、9個の成分すべてで紫外の極が $\MSbar$ の $\delta V$（式~\eqref{eq:dV}、\eqref{eq:Vrun}）と30桁で一致し、$\xi_W$ によらず、$V^\dagger\delta V+\delta V^\dagger V=0$ を満たした（有限部分も $4\times10^{-20}$ 以内）。有限部分は $\xi_W$ に依存し、$K|\delta V^{\mathrm{DS}}_{\mathrm{fin}}(\xi_W)-\delta V^{\mathrm{DS}}_{\mathrm{fin}}(\frac12)|$ の最大は $\xi_W=3$ で $1.3\times10^{-5}$ である（$K|\delta V^{\mathrm{DS}}_{\mathrm{fin}}|$ の最大は $1.5\times10^{-2}$）。Gambino らと Yamada の指摘したゲージ依存~\cite{ggm99,yamada01} は本ノートの模型でも現れるが、$|V_{ij}|$ と $J$ への効果は小さい。$V_0=V^{\MSbar}(\mu)+\delta V=V^{\mathrm{DS}}+\delta V^{\mathrm{DS}}$ で紫外の極が等しいので、$V^{\mathrm{DS}}=V^{\MSbar}(\mu)-\delta V^{\mathrm{DS}}_{\mathrm{fin}}$ である。$\mu=M_Z$ では、on-shell の $|V_{ub}|$、$|V_{cb}|$、$|V_{td}|$、$|V_{ts}|$ が $\MSbar$ の値より相対 $2.30\times10^{-3}$ 大きく、$J$ は $4.60\times10^{-3}$ 大きい。$\xi_W=\frac12$〜10 での変化は、$|V|$ で $1.9\times10^{-6}$、$J$ で $3.9\times10^{-6}$ にとどまる。二つの処方の差は、$\MSbar$ の $V$ が $\mu$ を $e$ 倍にしたときに走る量（$0.9\%$、式~\eqref{eq:Vhier}）の4分の1ほどである。

最後に、本ノートの模型で $V$ を基準点の値にした $W$ とトップの自己エネルギーから、SM ノートの量への CKM 行列の効果を厳密に求めた。$V=1$ の値との差は、$\Delta\hat\Delta=-7.4129\times10^{-8}$、$\Delta M_{W,\mathrm p}=2.1815\times10^{-7}$~GeV、$\Delta m_{t,\mathrm p}=-4.3205\times10^{-8}$~GeV（$\xi_W=1$、3 で同じ）である。SM ノートの見積もり（\texttt{ckm\_estimate.wl}、第~\ref{sec:intro} 節）とは相対 $3\times10^{-12}$ 以内で一致した。見積もりは $V=1$ の式の内線の質量を $|V_{ij}|^2$ の重みで入れ替えたもので、この一致は本ノートの模型の自己エネルギーがちょうどその形であることを示す。トップの差の紫外の極は $-\frac34\sum_j(|V_{tj}|^2-\delta_{j3})y_{d_j}^2=3.652\times10^{-7}$ に等しかった。スクリプトは \texttt{oneloop\_ckm\_feyncalc.wl}（約10分）、\texttt{oneloop\_ckm\_smmod.wl}、\texttt{oneloop\_ckm\_numeric.wl} で、\texttt{run\_oneloop\_check.sh} でまとめて実行できる。

\begin{table}[tbp]
\centering\small
\caption{手順 (3) の数値（基準点、$\mu=M_Z$、$\MSbar$ の有限部分の実部）。上の3行は1PI の非対角の成分（$s\leftarrow b$ は $p^2=m_b^2$、$c\leftarrow t$ は $p^2=m_t^2$）、下の5行は Denner--Sack の対抗項の有限部分による $\delta|V_{ij}|/|V_{ij}|=\re(V_{ij}^*\delta V^{\mathrm{DS}}_{\mathrm{fin},ij})/|V_{ij}|^2$ と $\delta J/J$ である。on-shell の値と $\MSbar$ の値の差はこれの符号を変えたもので（$V^{\mathrm{DS}}=V^{\MSbar}-\delta V^{\mathrm{DS}}_{\mathrm{fin}}$）。下の5行の $\xi_W$ への依存は $2\times10^{-6}$ 以下（$\delta J/J$ では $4\times10^{-6}$ 以下）である。}
\label{tab:ckmct}
\begin{tabular}{lrrrr}
\toprule
量 & $\xi_W=\frac12$ & $\xi_W=2$ & $\xi_W=3$ & $\xi_W=10$ \\
\midrule
$K\Sigma^L_{sb}(m_b^2)$ & $6.021\times10^{-3}$ & $3.506\times10^{-2}$ & $4.758\times10^{-2}$ & $9.681\times10^{-2}$ \\
$K\Sigma^S_{sb}(m_b^2)$ & $-1.794\times10^{-2}$ & $-3.246\times10^{-2}$ & $-3.872\times10^{-2}$ & $-6.333\times10^{-2}$ \\
$K\Sigma^L_{ct}(m_t^2)$ & $1.086\times10^{-5}$ & $2.056\times10^{-5}$ & $2.634\times10^{-5}$ & $-1.683\times10^{-5}$ \\
\midrule
$\delta|V_{us}|/|V_{us}|$ & $-2.44036\times10^{-6}$ & $-2.44002\times10^{-6}$ & $-2.43980\times10^{-6}$ & $-2.43995\times10^{-6}$ \\
$\delta|V_{ub}|/|V_{ub}|$ & $-2.29848\times10^{-3}$ & $-2.29967\times10^{-3}$ & $-2.30042\times10^{-3}$ & $-2.29991\times10^{-3}$ \\
$\delta|V_{cb}|/|V_{cb}|$ & $-2.29844\times10^{-3}$ & $-2.29963\times10^{-3}$ & $-2.30038\times10^{-3}$ & $-2.29987\times10^{-3}$ \\
$\delta|V_{td}|/|V_{td}|$ & $-2.30063\times10^{-3}$ & $-2.30182\times10^{-3}$ & $-2.30257\times10^{-3}$ & $-2.30206\times10^{-3}$ \\
$\delta J/J$ & $-4.59479\times10^{-3}$ & $-4.59716\times10^{-3}$ & $-4.59866\times10^{-3}$ & $-4.59763\times10^{-3}$ \\
\bottomrule
\end{tabular}
\end{table}

\begin{figure}[tbp]
\centering
\begin{tikzpicture}[scale=0.88]
% (a) neutral exchange between the quark legs
\begin{scope}[shift={(0.6,0)}]
  \draw[scalar] (-1.0,0) -- (0,0);
  \draw[fermion] (0,0) -- (0.6,0.42); \draw[fermion] (0.6,0.42) -- (1.2,0.84);
  \draw[fermion] (1.2,-0.84) -- (0.6,-0.42); \draw[fermion] (0.6,-0.42) -- (0,0);
  \draw[photon] (0.6,0.42) -- (0.6,-0.42);
  \node[vtx] at (0,0) {}; \node[vtx] at (0.6,0.42) {}; \node[vtx] at (0.6,-0.42) {};
  \node[font=\scriptsize] at (-0.6,0.3) {$G^+$};
  \node[font=\scriptsize,anchor=west] at (1.2,0.84) {$u_i$}; \node[font=\scriptsize,anchor=west] at (1.2,-0.84) {$d_j$};
  \node[font=\scriptsize,anchor=west] at (0.72,0) {$\gamma$、$Z$};
  \node[slab] at (0.2,-1.35) {(a) 中性の場の交換\\（$H$、$G^0$ も）：$V_{ij}$};
\end{scope}
% (b) boson triangle
\begin{scope}[shift={(4.1,0)}]
  \draw[scalar] (-1.0,0) -- (0,0);
  \draw[scalar] (0,0) -- (0.75,0.5); \draw[photon] (0,0) -- (0.75,-0.5);
  \draw[fermion] (1.3,-0.85) -- (0.75,-0.5); \draw[fermion] (0.75,-0.5) -- (0.75,0.5); \draw[fermion] (0.75,0.5) -- (1.3,0.85);
  \node[vtx] at (0,0) {}; \node[vtx] at (0.75,0.5) {}; \node[vtx] at (0.75,-0.5) {};
  \node[font=\scriptsize] at (-0.6,0.3) {$G^+$};
  \node[font=\scriptsize,anchor=west] at (1.3,0.85) {$u_i$}; \node[font=\scriptsize,anchor=west] at (1.3,-0.85) {$d_j$};
  \node[font=\scriptsize] at (1.0,0) {$d_j$};
  \node[font=\scriptsize] at (0.1,0.55) {$G^+$};
  \node[font=\scriptsize] at (0.1,-0.58) {$\gamma$、$Z$};
  \node[slab] at (0.2,-1.35) {(b) ボソンの三角\\（$W^+$、$H$、$G^0$ も）：$V_{ij}$};
\end{scope}
% (c) off-diagonal self-energy on the u leg
\begin{scope}[shift={(7.6,0)}]
  \draw[scalar] (-1.0,0) -- (0,0);
  \draw[fermion] (0,0) -- (0.55,0.385); \draw[fermion] (0.75,0.525) -- (1.25,0.875);
  \draw[fermion] (1.25,-0.875) -- (0,0);
  \node[blob,minimum size=0.32cm] at (0.65,0.455) {};
  \node[vtx] at (0,0) {};
  \node[font=\scriptsize] at (-0.6,0.3) {$G^+$};
  \node[font=\scriptsize] at (0.12,0.42) {$u_k$};
  \node[font=\scriptsize,anchor=west] at (1.25,0.875) {$u_i$}; \node[font=\scriptsize,anchor=west] at (1.25,-0.875) {$d_j$};
  \node[slab] at (0.2,-1.35) {(c) $u$ の場の繰り込み\\$\delta Z^{u}_{ik}$（非対角も）};
\end{scope}
% (d) off-diagonal self-energy on the d leg
\begin{scope}[shift={(11.0,0)}]
  \draw[scalar] (-1.0,0) -- (0,0);
  \draw[fermion] (0,0) -- (1.25,0.875);
  \draw[fermion] (1.25,-0.875) -- (0.75,-0.525); \draw[fermion] (0.55,-0.385) -- (0,0);
  \node[blob,minimum size=0.32cm] at (0.65,-0.455) {};
  \node[vtx] at (0,0) {};
  \node[font=\scriptsize] at (-0.6,0.3) {$G^+$};
  \node[font=\scriptsize] at (0.12,-0.45) {$d_k$};
  \node[font=\scriptsize,anchor=west] at (1.25,0.875) {$u_i$}; \node[font=\scriptsize,anchor=west] at (1.25,-0.875) {$d_j$};
  \node[slab] at (0.2,-1.35) {(d) $d$ の場の繰り込み\\$\delta Z^{d}_{kj}$（非対角も）};
\end{scope}
% (e) G+ self-energy
\begin{scope}[shift={(14.3,0)}]
  \draw[scalar] (-1.2,0) -- (-0.75,0); \draw[scalar] (-0.35,0) -- (0,0);
  \node[blob,minimum size=0.4cm] at (-0.55,0) {};
  \draw[fermion] (0,0) -- (0.85,0.6); \draw[fermion] (0.85,-0.6) -- (0,0);
  \node[vtx] at (0,0) {};
  \node[font=\scriptsize] at (-0.95,0.3) {$G^+$};
  \node[font=\scriptsize,anchor=west] at (0.85,0.6) {$u_i$}; \node[font=\scriptsize,anchor=west] at (0.85,-0.6) {$d_j$};
  \node[slab] at (-0.1,-1.35) {(e) $\Sigma^{G}$\\$\delta Z_G$};
\end{scope}
% (f) counterterm
\begin{scope}[shift={(17.0,0)}]
  \draw[scalar] (-1.0,0) -- (0,0);
  \draw[fermion] (0,0) -- (0.85,0.6); \draw[fermion] (0.85,-0.6) -- (0,0);
  \node[ctv] at (0,0) {};
  \node[font=\scriptsize] at (-0.6,0.3) {$G^+$};
  \node[font=\scriptsize,anchor=west] at (0.85,0.6) {$u_i$}; \node[font=\scriptsize,anchor=west] at (0.85,-0.6) {$d_j$};
  \node[slab] at (0,-1.35) {(f) 対抗項\\$\delta Y_d$、$\delta Y_u^\dagger$};
\end{scope}
\end{tikzpicture}
\caption{手順 (4) で湯川行列の $\MSbar$ の対抗項を解いた $\bar u_id_jG^+$ の頂点の図の型（外線の運動量は0）。$G^+$ は破線、$\gamma$、$Z$、$W$ は波線である。(a)、(b) の1PI の頂点の補正は、$G^+$ に付く荷電の頂点が一つで残りが中性の対角な結合なので $V_{ij}$ に比例する（1組あたり47個の図）。湯川行列の $\beta$ 関数の世代を混ぜる構造 $\frac32(Y_uY_u^\dagger-Y_dY_d^\dagger)$ は、クォークの場の繰り込みの非対角の成分 (c)、(d)（図~\ref{fig:seckm} (a)、(b)）だけから来る。(e) の $G^+$ の場の繰り込みは、$G^+$ の自己エネルギー（33個の図）の極の $k^2$ の係数である。(f) の二重項の基底の対抗項で (a)〜(e) の極が消える条件から $\delta Y_d$、$\delta Y_u^\dagger$ を行列として解き、式~\eqref{eq:betaY} と比べた。}
\label{fig:ugd}
\end{figure}

\paragraph{手順 (4)：湯川行列の $\beta$ 関数と、質量、$V$、$J$ の走り。}式~\eqref{eq:betaY} を二つの独立な方法で確かめた。一つ目は本ノートの模型からの導出である（\texttt{rg\_ckm\_feyncalc.wl}）。$\bar u_id_jG^+$ の1PI の頂点の1ループの紫外の極を、9個の世代の組のすべてで一般の $R_\xi$ ゲージで求めた（図~\ref{fig:ugd}）（外線の運動量は0にした。極は運動量によらない。1組あたり47個の図）。頂点の補正はどれも $V_{ij}$ だけに比例した。$G^+$ に付く荷電の頂点が一つで、ほかの結合は中性で対角だからである（図~\ref{fig:wud} の (a)、(b) と同じ理由）。$G^+$ の場の繰り込みは、$G^+$ の自己エネルギー（33個の図）の極の $k^2$ の係数からとった。クォークの場の繰り込みは、手順 (3) の $\Sigma^L$、$\Sigma^R$ の極からエルミートな行列としてとった。そのうえで、二重項の基底の裸のラグランジアンから作った対抗項で頂点が有限になる条件から、$\delta Y_d$ と $\delta Y_u^\dagger$ を行列として解いた。正弦と余弦がどちらも有理数のユニタリーな $V$（$s_{12}=8/17$、$s_{23}=7/25$、$s_{13}=3/5$、位相は記号）を入れると、18個の成分すべてが $\delta Y=\frac{\DUV}2\beta_Y$（式~\eqref{eq:betaY}）と記号として一致し、$\xi_A,\xi_Z,\xi_W,\xi_G$ を含まなかった。基準点の $V$ でも、二組のゲージパラメータで50桁の数値として一致した（差は $10^{-49}$ 以下）。頂点の補正が $V_{ij}$ にしか比例しないので、$\beta_Y$ の世代を混ぜる構造 $\frac32(Y_uY_u^\dagger-Y_dY_d^\dagger)$ はすべて、クォークの場の繰り込みの非対角の成分から来る。実際、非対角の成分を0にすると $\delta Y_d$ の9個の成分すべてが合わなくなった。

この途中で、手順 (3) の $\Sigma^L$、$\Sigma^R$ の対角の成分の極に、$k^2$ によらないはずなのに $C_Fg_s^2m^2(\xi_G-1)/(2k^2)$ という項が現れた。FeynCalc の還元が、グルーオンの縦の部分から運動量0の辺を持つ（質量のない伝播関数の2乗を含む）スカラーの $C_0$ を残すためである。この $C_0$ は赤外発散を持つので、紫外で有限として捨てると極の一部を落とす。計算環境の記録~\cite{env} の表 6 にある既知の問題である。$C_0$ を $D$ 次元の部分積分の恒等式（SM ノートでも使ったもの）で $A_0$、$B_0$ に直すと、極は $k^2$ によらなくなり、グルーオンの部分は共変ゲージの $C_Fg_s^2\xi_G$ になった。手順 (3) の結果はこの影響を受けない。非対角の成分はグルーオンを含まず、対角の on-shell の極ではこの項が $\Sigma^V+\Sigma^S$ の中で打ち消し、数値の結果は非対角の成分と $V$ による差（グルーオンの部分は打ち消す）だけを使ったからである。

二つ目は一般の公式である（\texttt{rg\_ckm\_general.py}）。SM ノートの手順 (4) の実装（45個の Weyl フェルミオンと4個の実のスカラー）を、複素の $3\times3$ の湯川行列に広げた。一般の公式の $\beta_Y^a$ は、すべてのスカラーの成分とフェルミオンの組で、式~\eqref{eq:betaY} を Weyl の形に直したものと相対 $2\times10^{-15}$ で一致した。質量の基底にない乱数の複素行列でも、基準点でも同じである。$\beta_\lambda$ と $\gamma_{\mu^2}$ は、特異値が同じで $V=1$ の値と相対 $5\times10^{-16}$ で一致し、1ループでは $V$ によらない。

最後に、式~\eqref{eq:betaY} の行列の方程式（ゲージ結合と $\mathcal T$ の項を含む）を4次の Runge--Kutta 法で $\pm h$、$\pm h/2$ だけ走らせ、特異値分解し、Richardson の外挿で微分を求めた。式~\eqref{eq:yrun} の湯川結合の走りは相対 $7\times10^{-13}$、$J$ の走りは $3\times10^{-11}$ で一致した。$|V_{ij}|$ の走りは、倍精度では $|V_{cd}|$（いちばん小さい微分 $K\,\dd\ln|V_{cd}|/\dd\ln\mu=-7.4\times10^{-5}$）だけが相対 $8\times10^{-7}$ ずれた。このずれは $h$ によらない。LAPACK の特異値分解の軽いクォークの特異ベクトルの誤差が、$Y$ の非対角の成分（$h$ に比例）に比例し、差分の商の中で打ち消さないためである。同じ計算を50桁（mpmath）で行うと、9個の $|V_{ij}|$ すべてが式~\eqref{eq:Vrun} と相対 $5\times10^{-15}$ で一致した。$V$ の位相を任意に付け替えても、式~\eqref{eq:Vrun} から出る $|V_{ij}|$ と $J$ の走りは相対 $2\times10^{-15}$ で変わらなかった。スクリプトは \texttt{rg\_ckm\_feyncalc.wl}（約10分）、\texttt{rg\_ckm\_general.py}、\texttt{ckm\_rg\_handcheck.py} で、\texttt{run\_rg\_check.sh} でまとめて実行できる。

\paragraph{手順 (5)：$W^+\to u_i\bar d_j$ の1ループの振幅と CKM の対抗項のゲージ依存。}スクリプト \texttt{wud\_vertex\_feyncalc.wl} は、本ノートの FeynArts の模型（光子とそのゴーストに質量 $M_A$ を与えた写し）で、崩壊できる6個のチャンネル（$i=u,c$、$j=d,s,b$）の1PI の電弱の頂点を、一般の $R_\xi$ ゲージで $D$ 次元のまま求めた。グルーオンの図は除いた（QCD の補正は $V_{ij}$ に比例し、$V=1$ と同じである）。外線を on-shell にし（$k^2=M_W^2$、$\varepsilon\cdot k=0$、Dirac 方程式）、振幅を
\begin{equation}
\mathcal M=\frac{e}{\sqrt2\sw}\,\bar u(p_1)\bigl[\slashed\varepsilon\,(F_LP_L+F_RP_R)+(\varepsilon\cdot p_1)(H_LP_L+H_RP_R)\bigr]v(p_2)
\label{eq:wudff}
\end{equation}
の4個の形状因子に分けた（木では $F_L=V_{ij}$、ほかは0）。どのチャンネルでも0でない図は18個で、4個の形状因子はすべて $V_{ij}$ だけに比例した（図~\ref{fig:wud} の (a)、(b) と同じ理由）。$F_R$、$H_L$、$H_R$ には紫外の極がなく、$F_L$ の極は $\xi_A,\xi_Z,\xi_W$ を含む。スピンと偏極の和は、4次元のトレースから閉じた式で求めた（同じスクリプトの (d)）。

振幅は、場の繰り込みが S 行列の要素から落ちることを使い、裸の場の振幅に外線の on-shell の LSZ の因子を掛けて組んだ（\texttt{wud\_renormalized.wl}）。$F_L$ は
\begin{equation}
F_L=V_{ij}+F_L^{\text{頂点}}+V_{ij}\Bigl(\frac{\delta g}{g}+\frac12\delta Z_W\Bigr)+\delta V_{ij}+\frac12\bigl[\delta Z^{u,L\dagger}V+V\delta Z^{d,L}\bigr]_{ij}
\label{eq:wudfl}
\end{equation}
である。$\delta Z^{q,L}$ は Denner の on-shell の場の定数~\cite{denner93}（非対角は式~\eqref{eq:dZos}、対角は $-\Sigma^L-m^2(\Sigma^{L\prime}+\Sigma^{R\prime}+2\Sigma^{S\prime})$ を $p^2=m^2$ で）、$\delta Z_W=-\Sigma_T^{W\prime}(M_W^2)$ で、手順 (3) の自己エネルギー（1PI とタドポール）から作った。$\MSbar$ では $\delta g/g=\frac{\DUV}2\beta_g/g$（$K\beta_g=-\frac{19}6g^3$）、$\delta V$ は式~\eqref{eq:dV} である。$\MSbar$ の場を使っても、その有限な LSZ の因子は $\MSbar$ の自己エネルギーから式~\eqref{eq:dZos} の形で作られるので、同じ式になる。on-shell の処方（FormCalc との比較）では、$\delta g/g$ を $\delta Z_e-\delta\sw/\sw$（Thomson 極限の電荷と on-shell の $M_W$、$M_Z$）に、$\delta V$ を $\delta V^{\mathrm{DS}}$ に替える。

正弦と余弦が有理数のユニタリーな $V$（手順 (4)）を入れると、6個のチャンネルすべてで $F_L$ の紫外の極が記号として0になった（質量、結合、$\xi_A,\xi_Z,\xi_W$ は記号のまま）。頂点と対角の脚の極はそれぞれ $\xi_W$ を含み、和で消える。外線の非対角の場の定数を除くと、6個すべてに極が残った。$\delta V^{\mathrm{DS}}$ の極は、9個の成分すべてで $\delta V$ の極と記号として一致した（手順 (3) の30桁の数値の確かめを記号にしたもの）。

ゲージによらないことは、基準点の数値（LoopTools、$\mu=M_Z$）で確かめた。光子の質量 $M_A$ は $W$ と $G^\pm$ の部分のゲージ不変性を破る。$\xi_Z$、$\xi_W$ の依存を打ち消す恒等式には光子の逆伝播関数が入り、その質量の項を補う頂点がないからである。そこで確かめを二つに分けた。(I) $\xi_A$ だけを変える（$\xi_A=\frac12,3,10$ と1の差、$\xi_Z=\xi_W=1$）。$M_A$ を残すので赤外発散はない。$F_L$ の差は、木の振幅に対して $M_A=1$、$0.1$、$0.01$~GeV で $9.6\times10^{-7}$、$1.9\times10^{-8}$、$3.0\times10^{-10}$ だった。6個のチャンネルで同じ値で、$M_A^2$ に比例して0に向かう。$F_R$、$H_L$、$H_R$ は $\xi_A$ を含まない。(II) $\xi_A=1$ にして $M_A\to0$ とし、$\xi_Z$、$\xi_W$ を変える（$(\xi_Z,\xi_W)=(2,3)$、$(\frac13,\frac12)$、$(5,10)$ と $(1,1)$ の差）。赤外発散は LoopTools の $\lambda$（$\lambda^2=1$~GeV$^2$）で抑えた。赤外発散する積分は外線のあいだの軟らかい光子の図にだけ現れ、その係数はゲージパラメータを含まないので、この処方はゲージ依存を持ち込まない。結果を表~\ref{tab:wud} に示す。$\MSbar$ の $F_L$ の差は、$c$ のチャンネルで $10^{-15}$〜$10^{-12}$、$u$ のチャンネルで $2\times10^{-11}$ だった。個々の項の差（$10^{-2}$ ほど）より9桁以上小さい。$u$ のチャンネルの残りは倍精度の限界で、$m_u^2/M_W^2=2.5\times10^{-10}$ のためである。$m_u$、$m_d$ を10倍にすると $9\times10^{-13}$ 以下になった。$F_R$、$H_L$、$H_R$ は $D$ 次元では $\xi_Z,\xi_W$ を含むが、その項は $D-4$ に比例し、4次元では（紫外の極から来る有理数の項を含めて）記号として0になった。$F_L$ の頂点ではそうならず、外線の因子と合わせて初めて消える。なお、光子の質量を残したまま $\xi_Z$、$\xi_W$ を変えると、$c$ のチャンネルの差は $M_A=1$、$0.1$、$0.01$~GeV で $3.0\times10^{-7}$、$5.2\times10^{-9}$、$7.5\times10^{-11}$ と $M_A^2$ に比例したが、$u$ のチャンネルでは $M_A\gg m_u$ のため $2.1\times10^{-3}$ のまま変わらなかった。

表~\ref{tab:wud} の内訳のうち、CKM 行列に固有なのは非対角の脚である。$b$ を含む二つのチャンネルでは、対角の下型の脚の $\xi$ 依存（$1.54\times10^{-2}$）が $d$、$s$ のチャンネル（$9.3\times10^{-3}$）より大きく、その分を非対角の下型の脚（$6.06\times10^{-3}$）が打ち消す。$\delta V$ を $\delta V^{\mathrm{DS}}$ に替えると、振幅は紫外で有限のまま $\xi$ に依存し、その差は $u\bar b$、$c\bar b$ で $1.5\times10^{-6}$、$c\bar d$、$c\bar s$ で $2\times10^{-9}$ 台、$u\bar s$ で $1\times10^{-9}$、$u\bar d$ で $5\times10^{-11}$ だった。二つの振幅の差は $\delta V^{\mathrm{DS}}_{\mathrm{fin}}$ そのものなので、これは手順 (3) の $\delta V^{\mathrm{DS}}$ のゲージ依存（表~\ref{tab:ckmct}）が振幅にそのまま現れたもので、Gambino らと Yamada の指摘~\cite{ggm99,yamada01} の確かめである。$\MSbar$ の振幅の残りは、それより $u\bar b$ で5桁、$c\bar b$ で9桁小さい。

\begin{table}[tbp]
\centering\footnotesize
\caption{手順 (5) のゲージ依存（\texttt{wud\_renormalized.wl}、基準点、$\mu=M_Z$、$\xi_A=1$、$M_A\to0$、$\lambda^2=1$~GeV$^2$）。$(\xi_Z,\xi_W)=(2,3)$、$(\frac13,\frac12)$、$(5,10)$ と Feynman ゲージ $(1,1)$ の $F_L$ の差の最大を、木の振幅 $|V_{ij}|$ に対して示す。上の6行は式~\eqref{eq:wudfl} の各項の有限部分（$\delta g$、$\delta V$ は $\MSbar$ では極だけなので0）、「$\MSbar$」はその和、「$\times10$」は $m_u$、$m_d$ を10倍にした点の和、「Denner--Sack」は $\delta V$ を $\delta V^{\mathrm{DS}}$ に替えた振幅の差である。}
\label{tab:wud}
\begin{tabular}{lrrrrrr}
\toprule
 & $u\bar d$ & $u\bar s$ & $u\bar b$ & $c\bar d$ & $c\bar s$ & $c\bar b$ \\
\midrule
頂点 & $2.89\times10^{-2}$ & $2.89\times10^{-2}$ & $2.89\times10^{-2}$ & $2.89\times10^{-2}$ & $2.89\times10^{-2}$ & $2.89\times10^{-2}$ \\
$W$ の脚 & $1.07\times10^{-2}$ & $1.07\times10^{-2}$ & $1.07\times10^{-2}$ & $1.07\times10^{-2}$ & $1.07\times10^{-2}$ & $1.07\times10^{-2}$ \\
上型の対角の脚 & $8.88\times10^{-3}$ & $8.88\times10^{-3}$ & $8.88\times10^{-3}$ & $8.88\times10^{-3}$ & $8.88\times10^{-3}$ & $8.88\times10^{-3}$ \\
下型の対角の脚 & $9.30\times10^{-3}$ & $9.31\times10^{-3}$ & $1.54\times10^{-2}$ & $9.30\times10^{-3}$ & $9.31\times10^{-3}$ & $1.54\times10^{-2}$ \\
上型の非対角の脚 & $1.3\times10^{-10}$ & $1.3\times10^{-9}$ & $4.0\times10^{-6}$ & $8.4\times10^{-9}$ & $7.0\times10^{-9}$ & $4.0\times10^{-6}$ \\
下型の非対角の脚 & $4.6\times10^{-7}$ & $1.04\times10^{-5}$ & $6.06\times10^{-3}$ & $2.7\times10^{-7}$ & $1.02\times10^{-5}$ & $6.06\times10^{-3}$ \\
\midrule
$\MSbar$ & $2.1\times10^{-11}$ & $2.0\times10^{-11}$ & $2.0\times10^{-11}$ & $1.4\times10^{-12}$ & $4.4\times10^{-14}$ & $1.0\times10^{-15}$ \\
$\MSbar$（$\times10$） & $5.6\times10^{-14}$ & $8.6\times10^{-13}$ & $7.4\times10^{-14}$ & $4.4\times10^{-13}$ & $5.2\times10^{-14}$ & $2.2\times10^{-15}$ \\
Denner--Sack & $4.9\times10^{-11}$ & $9.9\times10^{-10}$ & $1.46\times10^{-6}$ & $2.3\times10^{-9}$ & $2.6\times10^{-9}$ & $1.46\times10^{-6}$ \\
\bottomrule
\end{tabular}
\end{table}

Feynman ゲージでは、FeynArts の \texttt{SM.mod}（\texttt{\$CKM = True}）と FormCalc 9.10 で同じ6個のチャンネルを計算した（\texttt{wud\_formcalc.wl}）。\texttt{SM.mod} は Denner の規約で本ノートとは独立に書かれた模型で、on-shell の対抗項（Thomson 極限の電荷、on-shell の $M_W$、$M_Z$、Denner--Sack の \texttt{dCKM1}）を持つ。比べた量は、スピンと偏極について和をとった $2\,\re(\mathcal M_0^*\mathcal M_1)/|\mathcal M_0|^2$ の頂点の部分と対抗項の部分（$\lambda^2=1$~GeV$^2$、$\mu=M_Z$）と、くりこみ定数である。本ノートの側は式~\eqref{eq:wudfl} を on-shell の処方にし、$\delta Z_e-\delta\sw/\sw$ を SM ノートの $\Sigma^{AA}$、$\Sigma^{AZ}$、$\Sigma^{ZZ}$（1ループでは $V$ によらない）と手順 (3) の $\Sigma^{WW}$ から作った。頂点の部分は6個のチャンネルで相対 $2.4\times10^{-15}$ 以内、対抗項の部分は $1\times10^{-9}$ 以内（$c$ のチャンネルは $2\times10^{-11}$ 以内）で一致した。くりこみ定数（$K\delta Z^{u,L}$、$K\delta Z^{d,L}$ の18個、$K\delta V^{\mathrm{DS}}$ の9個、$K\delta Z_W$、$K(\delta Z_e-\delta\sw/\sw)$）の差は最大で $6\times10^{-9}$（赤外に敏感な $K\delta Z^{u,L}_{uu}$）だった。$m_u$、$m_d$ を10倍にした点では、三つの差がそれぞれ $2\times10^{-15}$、$2\times10^{-11}$、$6\times10^{-11}$ になったので、基準点の差は軽いクォークの積分の倍精度の限界である。FormCalc の和は LoopTools の $\Delta$ を1000変えても $6\times10^{-14}$ しか変わらず、紫外で有限だった。値はたとえば $u\bar d$ で頂点が $0.06001$、対抗項が $0.04058$ である（実の光子の放出を含まないので $\lambda$ による）。

この手順で、ツールの使い方の問題を二つ見つけた。第一に、1ループの補正の符号は、同じ過程の FeynArts の木の振幅に対する比として決めなければならない。FeynArts の木の振幅と表~\ref{tab:rules-ckm} の規則の比は過程によって違い、$W^+\to u\bar d$ では $i$、手順 (4) の $d\,G^+\to u$ では $-i$ である。最初は手順 (4) の「規則の単位の補正は FeynArts の振幅の $i$ 倍」をそのまま使ったので、頂点の符号が逆になった。$F_L$ に頂点の極の2倍が残ったことと、FormCalc の頂点の値が同じ大きさで逆の符号だったことから気づいた。第二に、FormCalc の \texttt{CalcRenConst} の結果は \texttt{RenConstList} に包まれ、\texttt{dSW1} が \texttt{dMWsq1} で、\texttt{dCKM1} が \texttt{dZfL1} で書かれるので、くりこみ定数を互いに代入しきってから数値にする必要があった。また、対抗項の頂点の世代の和は \texttt{SumOver} のまま残り、\texttt{HelicityME} は既定の \texttt{FermionChains -> Weyl} では働かず（\texttt{VA} にした）、$V^*$ は \texttt{CKMC} という記号で現れる。スクリプトは \texttt{wud\_vertex\_feyncalc.wl}（1チャンネル約6分、6個を並列に実行）、\texttt{wud\_renormalized.wl}（約2分）、\texttt{wud\_formcalc.wl} で、\texttt{run\_wud\_check.sh} でまとめて実行できる。

\begin{figure}[tbp]
\centering
\begin{tikzpicture}[scale=0.85]
% helper layout per panel: quark vertex Q = (0,0), lepton vertex L = (1.5,0)
% (a) tree
\begin{scope}[shift={(0.8,0)}]
  \draw[fermion] (-0.75,-0.75) -- (0,0); \draw[fermion] (0,0) -- (-0.75,0.75);
  \draw[photon] (0,0) -- (1.5,0);
  \draw[fermion] (1.5,0) -- (2.25,0.75); \draw[fermion] (2.25,-0.75) -- (1.5,0);
  \node[vtx] at (0,0) {}; \node[vtx] at (1.5,0) {};
  \node[font=\scriptsize,anchor=east] at (-0.75,-0.75) {$d_j$}; \node[font=\scriptsize,anchor=east] at (-0.75,0.75) {$u_i$};
  \node[font=\scriptsize,anchor=west] at (2.25,0.75) {$e^-$}; \node[font=\scriptsize,anchor=west] at (2.25,-0.75) {$\bar\nu_e$};
  \node[font=\scriptsize] at (0.75,0.28) {$W$};
  \node[slab] at (0.75,-1.25) {(a) 木：$V_{ij}$};
\end{scope}
% (b) quark-side neutral exchange
\begin{scope}[shift={(5.4,0)}]
  \draw[fermion] (-0.75,-0.75) -- (-0.4,-0.4); \draw[fermion] (-0.4,-0.4) -- (0,0); \draw[fermion] (0,0) -- (-0.4,0.4); \draw[fermion] (-0.4,0.4) -- (-0.75,0.75);
  \draw[photon] (-0.4,-0.4) -- (-0.4,0.4);
  \draw[photon] (0,0) -- (1.5,0);
  \draw[fermion] (1.5,0) -- (2.25,0.75); \draw[fermion] (2.25,-0.75) -- (1.5,0);
  \node[vtx] at (0,0) {}; \node[vtx] at (1.5,0) {}; \node[vtx] at (-0.4,0.4) {}; \node[vtx] at (-0.4,-0.4) {};
  \node[font=\scriptsize,anchor=east] at (-0.75,-0.75) {$d_j$}; \node[font=\scriptsize,anchor=east] at (-0.75,0.75) {$u_i$};
  \node[font=\scriptsize,anchor=west] at (2.25,0.75) {$e^-$}; \node[font=\scriptsize,anchor=west] at (2.25,-0.75) {$\bar\nu_e$};
  \node[font=\scriptsize,anchor=east] at (-0.48,0) {$\gamma$、$Z$};
  \node[font=\scriptsize] at (0.75,0.28) {$W$};
  \node[slab] at (0.75,-1.25) {(b) クォークの頂点\\$\gamma$、$Z$ の交換：$V_{ij}$};
\end{scope}
% (c) quark-side triangle: gamma or Z from the W to the u leg
\begin{scope}[shift={(10.0,0)}]
  \draw[fermion] (-0.75,-0.75) -- (0,0); \draw[fermion] (0,0) -- (-0.4,0.4); \draw[fermion] (-0.4,0.4) -- (-0.75,0.75);
  \draw[photon] (0,0) -- (0.65,0); \draw[photon] (0.65,0) -- (1.5,0);
  \draw[photon] (0.65,0) -- (-0.4,0.4);
  \draw[fermion] (1.5,0) -- (2.25,0.75); \draw[fermion] (2.25,-0.75) -- (1.5,0);
  \node[vtx] at (0,0) {}; \node[vtx] at (0.65,0) {}; \node[vtx] at (1.5,0) {}; \node[vtx] at (-0.4,0.4) {};
  \node[font=\scriptsize,anchor=east] at (-0.75,-0.75) {$d_j$}; \node[font=\scriptsize,anchor=east] at (-0.75,0.75) {$u_i$};
  \node[font=\scriptsize,anchor=west] at (2.25,0.75) {$e^-$}; \node[font=\scriptsize,anchor=west] at (2.25,-0.75) {$\bar\nu_e$};
  \node[font=\scriptsize] at (0.62,0.42) {$\gamma$、$Z$};
  \node[font=\scriptsize] at (1.08,-0.25) {$W$};
  \node[slab] at (0.75,-1.25) {(c) $W$ に付く $\gamma$、$Z$\\（$d_j$ の側も）：$V_{ij}$};
\end{scope}
% (d) mu decay for comparison
\begin{scope}[shift={(14.6,0)}]
  \draw[fermion] (-0.75,-0.75) -- (0,0); \draw[fermion] (0,0) -- (-0.75,0.75);
  \draw[photon] (0,0) -- (1.5,0);
  \draw[fermion] (1.5,0) -- (2.25,0.75); \draw[fermion] (2.25,-0.75) -- (1.5,0);
  \node[vtx] at (0,0) {}; \node[vtx] at (1.5,0) {};
  \node[font=\scriptsize,anchor=east] at (-0.75,-0.75) {$\mu^-$}; \node[font=\scriptsize,anchor=east] at (-0.75,0.75) {$\nu_\mu$};
  \node[font=\scriptsize,anchor=west] at (2.25,0.75) {$e^-$}; \node[font=\scriptsize,anchor=west] at (2.25,-0.75) {$\bar\nu_e$};
  \node[font=\scriptsize] at (0.75,0.28) {$W$};
  \node[slab] at (0.75,-1.25) {(d) $\mu$ 崩壊（比べる相手）};
\end{scope}
% (e) box, direct
\begin{scope}[shift={(0.8,-3.3)}]
  \draw[fermion] (-0.75,-0.95) -- (0,-0.45); \draw[fermion] (0,-0.45) -- (0,0.45); \draw[fermion] (0,0.45) -- (-0.75,0.95);
  \draw[photon] (0,-0.45) -- (1.5,-0.45); \draw[photon] (0,0.45) -- (1.5,0.45);
  \draw[fermion] (2.25,-0.95) -- (1.5,-0.45); \draw[fermion] (1.5,-0.45) -- (1.5,0.45); \draw[fermion] (1.5,0.45) -- (2.25,0.95);
  \node[vtx] at (0,-0.45) {}; \node[vtx] at (0,0.45) {}; \node[vtx] at (1.5,-0.45) {}; \node[vtx] at (1.5,0.45) {};
  \node[font=\scriptsize,anchor=east] at (-0.75,-0.95) {$d_j$}; \node[font=\scriptsize,anchor=east] at (-0.75,0.95) {$u_i$};
  \node[font=\scriptsize,anchor=west] at (2.25,0.95) {$e^-$}; \node[font=\scriptsize,anchor=west] at (2.25,-0.95) {$\bar\nu_e$};
  \node[font=\scriptsize] at (0.75,-0.7) {$W$}; \node[font=\scriptsize] at (0.75,0.7) {$\gamma$、$Z$};
  \node[font=\scriptsize] at (-0.25,0) {$u_i$}; \node[font=\scriptsize] at (1.8,0) {$e$};
  \node[slab] at (0.75,-1.45) {(e) 箱：$V_{ij}$};
\end{scope}
% (f) box, crossed
\begin{scope}[shift={(5.4,-3.3)}]
  \draw[fermion] (-0.75,-0.95) -- (0,-0.45); \draw[fermion] (0,-0.45) -- (0,0.45); \draw[fermion] (0,0.45) -- (-0.75,0.95);
  % gamma or Z from the d_j vertex to the e vertex, W from the u_i vertex to the nubar vertex
  \draw[photon] (0,-0.45) -- (1.5,0.45); \draw[photon] (0,0.45) -- (1.5,-0.45);
  \draw[fermion] (2.25,-0.95) -- (1.5,-0.45); \draw[fermion] (1.5,-0.45) -- (1.5,0.45); \draw[fermion] (1.5,0.45) -- (2.25,0.95);
  \node[vtx] at (0,-0.45) {}; \node[vtx] at (0,0.45) {}; \node[vtx] at (1.5,-0.45) {}; \node[vtx] at (1.5,0.45) {};
  \node[font=\scriptsize,anchor=east] at (-0.75,-0.95) {$d_j$}; \node[font=\scriptsize,anchor=east] at (-0.75,0.95) {$u_i$};
  \node[font=\scriptsize,anchor=west] at (2.25,0.95) {$e^-$}; \node[font=\scriptsize,anchor=west] at (2.25,-0.95) {$\bar\nu_e$};
  \node[font=\scriptsize] at (0.42,0.7) {$W$}; \node[font=\scriptsize] at (0.45,-0.72) {$\gamma$、$Z$};
  \node[font=\scriptsize] at (-0.25,0) {$d_j$}; \node[font=\scriptsize] at (1.75,0) {$e$};
  \node[slab] at (0.75,-1.45) {(f) 交差した箱：$V_{ij}$};
\end{scope}
% (g) d leg with the top loop: diagonal and off-diagonal
\begin{scope}[shift={(10.0,-3.3)}]
  \draw[fermion] (-0.95,-0.95) -- (-0.62,-0.62); \draw[fermion] (-0.38,-0.38) -- (0,0); \draw[fermion] (0,0) -- (-0.75,0.75);
  \node[blob,minimum size=0.36cm] at (-0.5,-0.5) {};
  \draw[photon] (0,0) -- (1.5,0);
  \draw[fermion] (1.5,0) -- (2.25,0.75); \draw[fermion] (2.25,-0.75) -- (1.5,0);
  \node[vtx] at (0,0) {}; \node[vtx] at (1.5,0) {};
  \node[font=\scriptsize,anchor=east] at (-0.95,-0.95) {$d_j$}; \node[font=\scriptsize,anchor=east] at (-0.75,0.75) {$u_i$};
  \node[font=\scriptsize] at (0.05,-0.42) {$d_k$};
  \node[font=\scriptsize,anchor=west] at (2.25,0.75) {$e^-$}; \node[font=\scriptsize,anchor=west] at (2.25,-0.75) {$\bar\nu_e$};
  \node[font=\scriptsize] at (0.75,0.28) {$W$};
  \node[slab] at (0.75,-1.45) {(g) $d$ の脚：$t$ が入る\\$V_{ik}\Sigma^d_{kj}$（$k=j$ も）};
\end{scope}
% (h) counterterm
\begin{scope}[shift={(14.6,-3.3)}]
  \draw[fermion] (-0.75,-0.75) -- (0,0); \draw[fermion] (0,0) -- (-0.75,0.75);
  \draw[photon] (0,0) -- (1.5,0);
  \draw[fermion] (1.5,0) -- (2.25,0.75); \draw[fermion] (2.25,-0.75) -- (1.5,0);
  \node[ctv] at (0,0) {}; \node[vtx] at (1.5,0) {};
  \node[font=\scriptsize,anchor=east] at (-0.75,-0.75) {$d_j$}; \node[font=\scriptsize,anchor=east] at (-0.75,0.75) {$u_i$};
  \node[font=\scriptsize,anchor=west] at (2.25,0.75) {$e^-$}; \node[font=\scriptsize,anchor=west] at (2.25,-0.75) {$\bar\nu_e$};
  \node[font=\scriptsize] at (0.75,0.28) {$W$};
  \node[slab] at (0.75,-1.45) {(h) $\delta V_{ij}$、$\delta g$};
\end{scope}
\end{tikzpicture}
\caption{手順 (6) の運動量0の $d_j\to u_ie^-\bar\nu_e$ の1ループの図の型（外線のフェルミオンと軽いクォークの質量は0）。(a) 木。(b)、(c) クォークの側の頂点の補正（レプトンの側の頂点は (d) の $\mu$ 崩壊と同じで、比 \eqref{eq:Rij} で落ちる）。(e)、(f) $W$ と $\gamma$ または $Z$ の箱で、$\gamma$、$Z$ はクォークの線のどちらにもレプトンの線のどちらにも付く（$\gamma$ は電荷を持つ線だけ）。(b)、(c)、(e)、(f) の内線のフェルミオンは外線の世代そのもので、どれも $V_{ij}$ にだけ比例し、$t$ を含まない。$t$ が入るのは (g) の下型の脚の自己エネルギー（対角と非対角、図~\ref{fig:seckm} (a)、(b)）だけで、それと (h) の $\delta V$ の和が式~\eqref{eq:Deltaij}、\eqref{eq:Gmu} の $G$ を与える。$W$ の自己エネルギー、$\delta g$、$e\nu_e$ の頂点と脚は (d) と共通で比から落ちる。セミレプトニックで頂点10個、箱6個、$\mu$ 崩壊で頂点8個、箱5個の図が0でない（一般の $R_\xi$）。}
\label{fig:sldiag}
\end{figure}

\paragraph{手順 (6)：セミレプトニックな崩壊で測る $|V_{ij}|$ と $\MSbar$ の $V(\mu)$。}SM ノートの手順 (6) と同じく、外線の運動量を0、外線のフェルミオンの質量を0にして、$d_j\to u_ie^-\bar\nu_e$（6個のチャンネル）と $\mu\to\nu_\mu e\bar\nu_e$ の1ループの振幅を一般の $R_\xi$ ゲージで求めた（図~\ref{fig:sldiag}）（\texttt{sl\_ckm\_feyncalc.wl}。光子の質量 $M_A$ を残し、グルーオンは除く）。運動量0の照合では、外線の世代の内線にも外線と同じ質量を使わなければならないので、トップ以外のクォークの質量は内線でも0にした。木、頂点、箱の図の数は、セミレプトニックで1、10、6個、$\mu$ 崩壊で1、8、5個だった。二つの Dirac の鎖をトレースで木の構造に射影すると、セミレプトニックの頂点と箱の補正は、6個のチャンネルで同じ関数の $V_{ij}$ 倍になった。頂点と箱の図の内線のフェルミオンは外線の世代そのもので、トップが入らないからである。二つの振幅の比
\begin{equation}
\frac{\mathcal M(d_j\to u_ie\bar\nu_e)}{V_{ij}\,\mathcal M(\mu\to\nu_\mu e\bar\nu_e)}=1+R_{ij}
\label{eq:Rij}
\end{equation}
では、$W$ の自己エネルギー、$\delta g$、$e\nu_e$ の頂点と脚が落ち、$R_{ij}$ は頂点と箱の差、外線の場の定数の差、$\delta V_{ij}/V_{ij}$ だけになる（\texttt{sl\_ckm\_check.wl}）。クォークの脚は手順 (3) の自己エネルギーから作った。軽いクォークの質量を0にする極限では、非対角の場の定数（式~\eqref{eq:dZos}）は質量の大小の向きだけで決まり、$m_k\gg m_j$ なら $\delta Z^L_{kj}=2\Sigma^S_{kj}(0)$、$m_k\ll m_j$ なら $-2[\Sigma^L_{kj}(0)+\Sigma^S_{kj}(0)]$ になる。$\MSbar$ の $\delta V$ も同じ極限で、$A_u=0$、$K(A_d)_{kj}=-\frac32s_{kj}(V^\dagger D_uV)_{kj}$（$m_{d_j}>m_{d_k}$ なら $s_{kj}=1$、逆なら $-1$）になる。

結果は次の通りである。(a) 有理数のユニタリーな $V$ で、6個のチャンネルの $R_{ij}$ の紫外の極は記号として0だった（光子の質量の項も残らない）。CKM の対抗項を入れた形の Sirlin の普遍性~\cite{sirlin78} である。(b) $\ln M_A$ の係数は6個のチャンネルで同じ $-4e^2/K$ で、ほかに $M_A$ は残らない。(c) $R_{ij}$ は $\xi_A,\xi_Z,\xi_W$ によらない（下型の脚と $\delta V$ だけをとると $\xi_W$ による）。(d) トップの質量を0にした極限では、$R_{ij}$ は6個のチャンネルで同じで、$V$ を含まず、
\begin{equation}
R_{ij}\big|_{M_T\to0}=\frac{\alpha}{\pi}\ln\frac{M_Z}{M_A}
\label{eq:sirlinSD}
\end{equation}
にちょうど等しかった（定数項がなく、ゲージと $\mu$ によらない）。これは Sirlin の短距離の因子~\cite{sirlin82}（崩壊率では $1+\frac{2\alpha}\pi\ln\frac{M_Z}\mu$）を、運動量0で光子の質量を赤外の切断にした形で再現したものである。(e) 残りのトップの寄与 $\Delta_{ij}=R_{ij}-R_{ij}|_{M_T\to0}$ は、6個のチャンネルで一つの関数 $G(\mu)$ と CKM の因子の積
\begin{equation}
\Delta_{ib}=|V_{tb}|^2G,\qquad \Delta_{id}=-|V_{td}|^2G,\qquad \Delta_{is}=\Bigl(|V_{ts}|^2+2\frac{V_{ib}V_{tb}^*V_{ts}}{V_{is}}\Bigr)G
\label{eq:Deltaij}
\end{equation}
に記号として一致した。$G$ は
\begin{equation}
KG(\mu)=\frac{y_t^2}8\Bigl[-\frac{5x_t-11}{x_t-1}+\frac{6x_t(x_t-2)\ln x_t}{(x_t-1)^2}\Bigr]-\frac34y_t^2\ln\frac{\mu^2}{M_W^2},\qquad x_t=\frac{M_T^2}{M_W^2}
\label{eq:Gmu}
\end{equation}
である（$M_T$、$M_W$、$y_t$ は $\MSbar$）。式~\eqref{eq:Deltaij} の形は手で導ける。下型の自己エネルギーのトップの部分は、ユニタリー性で $V^*_{tk}V_{tj}$ に比例する。非対角の場の定数と $\delta V$ は世代の向きの符号だけが違い、$\frac12\Sigma^L+\Sigma^S$ のトップの部分と $\delta V$ の極の和が $G$ になる。$\Sigma^L$、$\Sigma^S$ はそれぞれ $\xi_W$ に依存するが、$G$ はゲージによらない。$K\,\dd G/\dd\ln\mu=-\frac32y_t^2$ なので、$|V_{ij}(\mu)|(1+\re\Delta_{ij})$ は式~\eqref{eq:Vhier} の走りと打ち消して $\mu$ によらない。

$|V_{ij}|$ の測定値は、$\mu$ 崩壊の $G_F$ で規格化し、Sirlin の普遍的な補正、QCD とハドロンの補正を除いて求められる。トップによる世代ごとに違う電弱の補正は除かれないので、本ノートでは PDG の $|V_{ij}|$ を $V^{\mathrm{phys}}_{ij}=V_{ij}(\mu)(1+\Delta_{ij}(\mu))$ の絶対値とみなす。基準点の $\mu=M_Z$ では $G=2.3028\times10^{-3}$ で、$\Delta_{ub}=\Delta_{cb}=2.299\times10^{-3}$、$\Delta_{cs}=-4.2\times10^{-6}$、$\Delta_{us}=(2.5+2.8i)\times10^{-6}$、$\Delta_{ud}=\Delta_{cd}=-1.7\times10^{-7}$ である。$\mu$ を2倍にすると $G$ は $6.05\times10^{-3}$ 減り、$\mu=M_t$ では $-3.1\times10^{-3}$ になる。この $\re\Delta_{ij}$ を手順 (3) の Denner--Sack の $|V^{\mathrm{DS}}_{ij}|/|V^{\MSbar}_{ij}|-1$（表~\ref{tab:ckmct}）と比べると、$u\bar b$、$c\bar b$ では $\xi_W$ による幅（$2.2985$〜$2.3004\times10^{-3}$）の中にあり、$c\bar s$ では相対 $2\times10^{-3}$ 以内、$u\bar d$、$u\bar s$、$c\bar d$ では差が $1.1\times10^{-7}$ 以下だった。軽いクォークの極限では、Denner--Sack の on-shell の $V$ は、そのゲージ依存の範囲で、セミレプトニックな崩壊で測る $V$ と同じになる。

\texttt{ckm\_pdg\_msbar.py} は、PDG~2024 の標準のパラメータ~\cite{pdg24}（式~\eqref{eq:bench} の値）を $V^{\mathrm{phys}}$ として、$|V_{us}|$、$|V_{cb}|$、$|V_{ub}|$ と角 $\gamma=\arg(-V_{ud}V_{ub}^*/V_{cd}V_{cb}^*)$（$\gamma$ に入る4個の $\Delta$ は実数なので変わらない）から $\MSbar$ の $V(M_Z)$ を解き、式~\eqref{eq:betaY} の行列の方程式（ゲージ結合も1ループ）で $M_t=172.57$~GeV まで走らせた（表~\ref{tab:pdgms}）。$\MSbar$ の $s_{23}$、$s_{13}$ は $M_Z$ で PDG の値より $0.229\%$ 小さく、$J$ は $0.458\%$ 小さい。$s_{12}$、$\delta$ の変化は $3\times10^{-6}$ 以下である。差は PDG の誤差（$s_{23}$ で1.9\%、$s_{13}$ で2.4\%）より1桁小さい。$M_t$ では、換算した $s_{23}$ は 0.041960 で、PDG の値を $M_Z$ の $\MSbar$ の値とみなして走らせた値（0.042056。本ノートのこれまでの基準点の扱い）より $0.23\%$ 小さい。$\mu=M_t$ で式~\eqref{eq:Gmu} を直接使った換算は、走らせた値と相対 $1.4\times10^{-6}$ で一致した。手順 (1)〜(5) の数値は、基準点の $V$ をこの分だけ変えても結論は変わらない（$|V_{cb}|^2$ に比例する幅で $0.46\%$）。スクリプトは \texttt{sl\_ckm\_feyncalc.wl}（約12分）、\texttt{sl\_ckm\_check.wl}、\texttt{ckm\_pdg\_msbar.py} で、\texttt{run\_sl\_check.sh} でまとめて実行できる。

\begin{table}[tbp]
\centering\small
\caption{手順 (6) の CKM 行列の換算（\texttt{ckm\_pdg\_msbar.py}）。PDG~2024 の値を $V^{\mathrm{phys}}$ とし、式~\eqref{eq:Deltaij}、\eqref{eq:Gmu} で $\MSbar$ の $V(M_Z)$ に直し、式~\eqref{eq:betaY} で $M_t=172.57$~GeV まで走らせた。最後の列は、PDG の値をそのまま $M_Z$ の $\MSbar$ の値として走らせたもの（基準点の扱い）。}
\label{tab:pdgms}
\begin{tabular}{lllll}
\toprule
 & PDG（$V^{\mathrm{phys}}$） & $\MSbar$、$M_Z$ & $\MSbar$、$M_t$ & PDG を $M_Z$ の $\MSbar$ として、$M_t$ \\
\midrule
$s_{12}$ & 0.22501 & 0.225009 & 0.225011 & 0.225011 \\
$s_{23}$ & 0.04183 & 0.041734 & 0.041960 & 0.042056 \\
$s_{13}$ & 0.003732 & 0.0037234 & 0.0037436 & 0.0037522 \\
$\delta$ & 1.147 & 1.146997 & 1.146995 & 1.146998 \\
$J$ [$10^{-5}$] & 3.1170 & 3.1027 & 3.1363 & 3.1508 \\
\bottomrule
\end{tabular}
\end{table}

\begin{figure}[tbp]
\centering
\begin{tikzpicture}[scale=0.88]
% (a), (b): Z on the internal u_k pair, W or G rung
\foreach \x/\sty/\rung/\lbl in {0.6/photon/{$W$}/{(a) $Z$ が $u_k$ に付く\\$W$ の交換}, 3.9/scalar/{$G^\pm$}/{(b) $Z$ が $u_k$ に付く\\$G^\pm$ の交換}} {
  \begin{scope}[shift={(\x,0)}]
    \draw[photon] (-1.0,0) -- (0,0);
    \draw[fermion] (0,0) -- (0.65,0.45); \draw[fermion] (0.65,0.45) -- (1.25,0.85);
    \draw[fermion] (1.25,-0.85) -- (0.65,-0.45); \draw[fermion] (0.65,-0.45) -- (0,0);
    \draw[\sty] (0.65,0.45) -- (0.65,-0.45);
    \node[vtx] at (0,0) {}; \node[vtx] at (0.65,0.45) {}; \node[vtx] at (0.65,-0.45) {};
    \node[font=\scriptsize] at (-0.6,0.3) {$Z$};
    \node[font=\scriptsize,anchor=west] at (1.25,0.85) {$b$}; \node[font=\scriptsize,anchor=west] at (1.25,-0.85) {$\bar s$};
    \node[font=\scriptsize] at (0.18,0.42) {$u_k$}; \node[font=\scriptsize] at (0.18,-0.42) {$u_k$};
    \node[font=\scriptsize,anchor=west] at (0.72,0) {\rung};
    \node[slab] at (0.2,-1.35) {\lbl};
  \end{scope}}
% (c), (d): Z on the boson pair, internal u_k
\foreach \x/\sty/\lbl in {7.2/photon/{(c) $ZW^+W^-$}, 10.5/scalar/{(d) $ZG^+G^-$}} {
  \begin{scope}[shift={(\x,0)}]
    \draw[photon] (-1.0,0) -- (0,0);
    \draw[\sty] (0,0) -- (0.75,0.5); \draw[\sty] (0,0) -- (0.75,-0.5);
    \draw[fermion] (1.3,-0.85) -- (0.75,-0.5); \draw[fermion] (0.75,-0.5) -- (0.75,0.5); \draw[fermion] (0.75,0.5) -- (1.3,0.85);
    \node[vtx] at (0,0) {}; \node[vtx] at (0.75,0.5) {}; \node[vtx] at (0.75,-0.5) {};
    \node[font=\scriptsize] at (-0.6,0.3) {$Z$};
    \node[font=\scriptsize,anchor=west] at (1.3,0.85) {$b$}; \node[font=\scriptsize,anchor=west] at (1.3,-0.85) {$\bar s$};
    \node[font=\scriptsize] at (1.0,0) {$u_k$};
    \node[slab] at (0.2,-1.35) {\lbl};
  \end{scope}}
% (e): Z W G
\begin{scope}[shift={(13.8,0)}]
  \draw[photon] (-1.0,0) -- (0,0);
  \draw[photon] (0,0) -- (0.75,0.5); \draw[scalar] (0,0) -- (0.75,-0.5);
  \draw[fermion] (1.3,-0.85) -- (0.75,-0.5); \draw[fermion] (0.75,-0.5) -- (0.75,0.5); \draw[fermion] (0.75,0.5) -- (1.3,0.85);
  \node[vtx] at (0,0) {}; \node[vtx] at (0.75,0.5) {}; \node[vtx] at (0.75,-0.5) {};
  \node[font=\scriptsize] at (-0.6,0.3) {$Z$};
  \node[font=\scriptsize,anchor=west] at (1.3,0.85) {$b$}; \node[font=\scriptsize,anchor=west] at (1.3,-0.85) {$\bar s$};
  \node[font=\scriptsize] at (1.0,0) {$u_k$};
  \node[font=\scriptsize] at (0.12,0.55) {$W$}; \node[font=\scriptsize] at (0.12,-0.58) {$G$};
  \node[slab] at (0.2,-1.35) {(e) $ZW^\mp G^\pm$\\（逆の組も）};
\end{scope}
% (f): external mixing on the b leg
\begin{scope}[shift={(17.0,0)}]
  \draw[photon] (-1.0,0) -- (0,0);
  \draw[fermion] (0,0) -- (0.55,0.385); \draw[fermion] (0.75,0.525) -- (1.25,0.875);
  \draw[fermion] (1.25,-0.875) -- (0,0);
  \node[blob,minimum size=0.32cm] at (0.65,0.455) {};
  \node[vtx] at (0,0) {};
  \node[font=\scriptsize] at (-0.6,0.3) {$Z$};
  \node[font=\scriptsize] at (0.15,0.42) {$s$};
  \node[font=\scriptsize,anchor=west] at (1.25,0.875) {$b$}; \node[font=\scriptsize,anchor=west] at (1.25,-0.875) {$\bar s$};
  \node[slab] at (0.2,-1.35) {(f) 外線の混合\\（$\bar s$ の側も）};
\end{scope}
\end{tikzpicture}
\caption{手順 (7) の $Z\to b\bar s$ の1ループの図の型。内線の $u_k$ について $\lambda_k^*=V^*_{kb}V_{ks}$ を掛けて足す。(a)〜(e) が1PI の頂点（一般の $R_\xi$ で0でない図は18個）で、$Z$ は内線の $u_k$（(a)、(b)）、荷電のボソンの対（(c)、(d)）、$ZW^\mp G^\pm$ の頂点（(e)、$M_Z\sw^2$ に比例）に付く。(f) は外線の世代の混合で、対角の $Z$ の結合 $g_L,g_R$ に on-shell の場の定数（式~\eqref{eq:dZos}、\eqref{eq:dZRos}）を掛けた形で入る。個々の図は $\xi_W$ に依存し、on-shell（$k^2=M_Z^2$）の和は $\xi$ によらない。運動量0の Feynman ゲージの和が Inami--Lim の $C_0(x)$（式~\eqref{eq:C0IL}）である。$b\to s\gamma$ は $Z$ を光子に替えた図で、(e) が $\gamma W^\mp G^\pm$（$eM_W$ に比例）になる（図~\ref{fig:fcncdiag} (a)〜(f)）。}
\label{fig:zbsdiag}
\end{figure}

\paragraph{手順 (7)：$b\to s\gamma$ と $Z\to b\bar s$。}スクリプト \texttt{fcnc\_ckm\_feyncalc.wl} は、本ノートの FeynArts の模型で、on-shell の $b\to s\gamma$（$p_1^2=m_b^2$、$p_2^2=0$、$q^2=0$）と $Z(k)\to b\bar s$（$k^2$ は記号のまま）の1PI の1ループの頂点を、一般の $R_\xi$ ゲージで求めた（$s$ の質量は0、グルーオンは除く。図~\ref{fig:fcncdiag} (a)〜(f)、図~\ref{fig:zbsdiag}）。どちらも0でない図は18個で、内線の上型のクォーク $u_k$ について $\lambda_k=V^*_{ks}V_{kb}$（$Z\to b\bar s$ では $\lambda_k^*$）を掛けて足す形になった。振幅は手順 (5) の4個の構造（式~\eqref{eq:wudff} の括弧）の係数 $F_L,F_R,H_L,H_R$ に分けた。木の結合がないので対抗項はなく、外線の世代の混合は on-shell の場の定数で入る（\texttt{fcnc\_ckm\_numeric.wl}）。光子と $Z$ の結合 $\bar\psi_0\gamma^\mu(g_LP_L+g_RP_R)\psi_0$ から、$b\to s\gamma$ の脚の項は $g_{L,R}\frac12(\delta Z^{L,R}_{sb}+\delta Z^{L,R*}_{bs})$ で、$m_s=0$ では式~\eqref{eq:dZos}、\eqref{eq:dZRos} から $\delta Z^L_{sb}=-2(\Sigma^L+\Sigma^S)_{sb}(m_b^2)$、$\delta Z^L_{bs}=2\Sigma^S_{bs}(0)$、$\delta Z^R_{sb}=-2\Sigma^R_{sb}(m_b^2)$、$\delta Z^R_{bs}=0$ になる（$Z\to b\bar s$ では添字を入れ替える）。$g_{L,R}$ は、同じ外線の配置の対角の木（$b\to b\gamma$、$Z\to b\bar b$）の FeynArts の振幅からとった。

結果は次の通りである。(a) 有理数のユニタリーな $V$ で、二つの過程の4個の形状因子の紫外の極は、外線の項を含めると記号として0になった（$Z\to b\bar s$ は任意の $k^2$ で）。1PI の部分だけでは $F_L$ に極が残る。(b) $b\to s\gamma$ を基準点で、Feynman パラメータの表示（内側の積分を閉じた形にした1次元の積分）で40桁で求めた。この運動学には閾値がない。PaVe の還元の Gram 行列式が $m_b^4$ に比例するので、LoopTools の倍精度では6桁を失う（使った積分は LoopTools の値と相対 $4\times10^{-11}$ 以内で一致した）。電磁気のゲージ不変性 $F_L+\frac{m_b}2H_R=0$、$F_R+\frac{m_b}2H_L=0$ は作業精度で成り立ち（$F_R=H_L=0$）、$\xi_W=\frac12,2,3,10$ の値は $\xi_W=1$ と作業精度で一致した。1PI の $F_L$ は $\xi_W$ を3にすると 163\% 変わるが、双極子 $H_R$ は1PI だけで $\xi_W$ によらない（脚は $H$ に寄与しない）。上型のクォークの質量をすべて等しくすると、振幅は0になった（GIM）。(c) 双極子は、Inami と Lim の関数~\cite{il81,bbl96}
\begin{equation}
D_0'(x)=-\frac{8x^3+5x^2-7x}{12(1-x)^3}+\frac{x^2(2-3x)\ln x}{2(1-x)^4}\qquad\bigl(C_7(M_W)=-\tfrac12D_0'(x_t)\bigr)
\label{eq:D0p}
\end{equation}
で $KH_R=-\frac{g^2em_b}{M_W^2}\sum_k\lambda_kD_0'(x_k)\,[1+O(m_b^2/M_W^2)]$ になった。比は $m_t=100$、165、300~GeV で、$m_b=2.86$~GeV では $-1.00080$、$-1.00065$、$-1.00054$、$m_b$ を10分の1にすると $-1.0000080$、$-1.0000065$、$-1.0000054$ で、1からのずれは $m_b^2$ に比例する（全体の符号は FeynArts の振幅の規約による）。(d) $Z\to b\bar s$ の on-shell（$k^2=M_Z^2$）の振幅は、$Z\to u\bar u$、$c\bar c$ が開いているので LoopTools で求めた。4個の形状因子は $\xi_W$ によらず（差は $1.7\times10^{-13}$）、1PI の $F_L$ だけでは 42\% 変わった。$k^2\to0$ の Feynman ゲージの $F_L$ は、Inami--Lim の $Z$ ペンギンの関数
\begin{equation}
C_0(x)=\frac x8\Bigl[\frac{x-6}{x-1}+\frac{(3x+2)\ln x}{(x-1)^2}\Bigr]
\label{eq:C0IL}
\end{equation}
で $KF_L=-\frac{2g^3}{\cw}\sum_k\lambda_k^*C_0(x_k)\,[1+O(m_b^2/M_W^2)]$ になった（比は $m_t=165$~GeV で $m_b=2.86$~GeV では $-1.000226$、10分の1では $-1.0000023$。$k^2=10^{-16}M_Z^2$ で、$1/k^2$ の項は60桁の計算で打ち消した）。同じ $k^2\to0$ の値は、$\xi_W=3$ にすると 24\% 変わった。$C_0$ は運動量0の外れた $Z$ の量で、ゲージに依存し、箱の図と組んで初めてゲージによらなくなる~\cite{bbl96}。物理的な on-shell の値はゲージによらず、$C_0$ の単位で $0.760-0.234i$ である（虚部は $Z\to u\bar u$、$c\bar c$ の吸収部分）。

この手順で、式~\eqref{eq:dZRos} の右巻きの場の定数の形の誤りに気づいた。最初は $\delta Z^R$ も $\delta Z^L$ と同じ形（$\Sigma^S$ の係数 $m_i^2+m_j^2$）と思い込み、$b\to s\gamma$ の右巻きの脚が0にならずに電磁気のゲージ不変性が 31\% 破れ、$Z\to b\bar s$ に $\xi_W$ 依存（$7\times10^{-6}$）が残った。on-shell の条件を解き直すと係数は $2m_im_j$ で、$m_s=0$ では $\delta Z^R_{bs}=0$ になる。手順 (3)〜(6) は $\delta Z^L$ だけを使ったので影響を受けない。スクリプトは \texttt{fcnc\_ckm\_feyncalc.wl}（2個の過程を並列に約2分）、\texttt{fcnc\_ckm\_numeric.wl} で、\texttt{run\_fcnc\_check.sh} でまとめて実行できる（約4分）。

\section{議論と今後}\label{sec:discussion}

\paragraph{現時点で分かったこと。}CKM 行列は、二重項の下の成分を $Vd_L$ にするだけで模型に入り、頂点の式が変わるのは荷電のクォークの4個のクラスだけである。中性の結合が対角になることは、FeynRules の出力、標準のパラメータ表示を入れた値、MadGraph が世代を変える中性の過程に図を作らないことの三つで確かめた。$V=1$ で SM ノートの模型に戻るので、SM ノートの結果のうち $V$ によらないものは、本ノートの模型でもそのまま成り立つ。

手計算では、$\MSbar$ の CKM の対抗項が湯川行列の $\beta$ 関数だけで決まり（式~\eqref{eq:dV}）、FJ のタドポールが入らないことが分かった。SM ノートの質量の対抗項がタドポールを通じて $1/\lambda$ の大きな項を持ったのと違い、CKM 行列の走りには $v$ の走りが入らない。自己エネルギーの非対角の紫外の極は $G^\pm$ のループの $m_{u_k}^2$ の部分だけで、$\xi$ によらない（式~\eqref{eq:uvoff}）。on-shell の Denner--Sack の対抗項の紫外の極は $\MSbar$ の $\delta V$ と同じで、二つの処方の違いは有限部分にあり、そこに $\xi$ 依存が入る。$V$ の走り（式~\eqref{eq:Vrun}）は、基準点で $|V_{ub}|$、$|V_{cb}|$、$|V_{td}|$、$|V_{ts}|$ を $\mu$ が $e$ 倍になるごとに $1.38/K\approx0.9\%$ 大きくする。

木のレベルでは、振幅の2乗が $|V_{ij}|^2$ だけで書け、CP が破れず、ユニタリー性で $V=1$ の値に戻る過程があることを、FeynCalc、明示的な行列、MadGraph の三つで確かめた（手順 (2)）。$t\bar s\to W^+H$ が $\xi_W$ によらないことは、$W$ と $G^\pm$ の頂点の CKM 行列が整合していることのゲージ不変性による確かめである。

\paragraph{自己エネルギーと CKM の対抗項。}手順 (3) で、第~\ref{sec:se} 節の構造がすべて確かめられた。世代を変える自己エネルギーの紫外の極は $G^\pm$ のループの $m_{u_k}^2$ の部分だけで、$\xi$ によらず、$\MSbar$ の CKM の対抗項は湯川行列の $\beta$ 関数で決まる。有限部分は $\xi_W$ で大きく変わり（$\Sigma^L_{sb}$ は $\xi_W=\frac12$ と 10 で16倍違う）、それから作る Denner--Sack の対抗項の有限部分も $\xi_W$ による。ただしその依存は対抗項の大きさの0.1\% 程度で、$|V_{ij}|$ では $2\times10^{-6}$ 以下である。on-shell の処方の不整合は原理の問題で、数値としては、$\MSbar$ と on-shell の処方の差（$|V_{cb}|$ で $2.3\times10^{-3}$）や現在の測定の誤差（$|V_{cb}|$ で約3\%~\cite{pdg24}）より十分小さい。

\paragraph{湯川行列の $\beta$ 関数。}手順 (4) で、式~\eqref{eq:betaY} を本ノートの模型の頂点と場の繰り込みから導き、一般の公式とも一致させた。$\bar u_id_jG^+$ の頂点の補正は $V_{ij}$ にしか比例しないので、CKM 行列の走りを生む $\frac32(Y_uY_u^\dagger-Y_dY_d^\dagger)$ は、クォークの場の繰り込みの非対角の成分（手順 (3) の式~\eqref{eq:uvoff}）だけから来る。手順 (3) の非対角の自己エネルギーの極と、CKM 行列の $\MSbar$ の対抗項と、行列の $\beta$ 関数は、この意味で一つの構造の三つの面である。

\paragraph{$W\to u_i\bar d_j$ の振幅とゲージ依存。}手順 (5) で、$\MSbar$ の $V$ と外線の on-shell の LSZ の因子で組んだ $W^+\to u_i\bar d_j$ の振幅が、紫外で有限で（記号として）、$\xi_Z$、$\xi_W$ によらないこと（木の振幅の $2\times10^{-11}$ 以下、倍精度の限界）を確かめた。外線の非対角の場の定数は、紫外の極を消すのにも、$\xi$ 依存を消すのにも要る。$b$ を含むチャンネルでは、その $\xi$ 依存（$6\times10^{-3}$）が対角の脚の $\xi$ 依存の世代による違いをちょうど打ち消す。$\delta V$ を Denner--Sack の対抗項に替えると、振幅に $\xi$ 依存が残る（$u\bar b$、$c\bar b$ で $1.5\times10^{-6}$）。これは $\delta V^{\mathrm{DS}}_{\mathrm{fin}}$ のゲージ依存（手順 (3)）がそのまま現れたもので、$\MSbar$ の処方ではこの問題が起きない。検算の方法としては、光子の質量で赤外発散を抑えると $W$ の部分のゲージ不変性が $M_A^2$ の程度で破れ、$M_A$ が外線のクォークの質量より大きいと残りが消えないことが分かった。$\xi_Z$、$\xi_W$ の確かめには、光子の質量を0にして赤外発散する積分にだけ $\lambda$ を入れる処方が要る。

\paragraph{測った $|V_{ij}|$ と $\MSbar$ の $V$。}手順 (6) で、セミレプトニックな崩壊と $\mu$ 崩壊の運動量0の振幅の比が、$\MSbar$ の $V$ で紫外で有限、ゲージによらないことを確かめた。比は、すべてのクォークの質量を0にした普遍的な部分 $\frac\alpha\pi\ln\frac{M_Z}{M_A}$（Sirlin の因子そのもの）と、トップによる部分 $\Delta_{ij}$（式~\eqref{eq:Deltaij}）に分かれる。CKM 行列に固有なのは後者だけで、一つのゲージによらない関数 $G$ と CKM の因子で書ける。$G(M_Z)=2.3\times10^{-3}$ なので、$\MSbar$ の $|V_{ub}|$、$|V_{cb}|$ は $M_Z$ で測定値より $0.23\%$ 小さく、$M_t$ では $0.31\%$ 大きい。これは今の測定の誤差（2〜3\%）より小さいが、$V$ の $\MSbar$ の値を高いスケールまで走らせるとき（手順 (10)）の出発点を決める。Denner--Sack の on-shell の $V$ は、軽いクォークの極限では、そのゲージ依存（$10^{-6}$）の範囲でこの物理的な $V$ と一致する。

\paragraph{世代を変える過程。}手順 (7) で、$b\to s\gamma$ と $Z\to b\bar s$ の on-shell の振幅が、外線の世代の混合を入れると紫外で有限で、ゲージによらず、GIM の機構で上型のクォークの質量の差だけによることを確かめた。$b\to s\gamma$ では外線の項は電荷の形状因子だけに入り、双極子は1PI の図だけで決まって、Inami--Lim の関数に $m_b^2/M_W^2$ の補正の範囲で一致する。$Z\to b\bar s$ では、運動量0の Inami--Lim の関数は Feynman ゲージの外れた量で、ゲージによらない on-shell の振幅とは $24\%$ 違う。外線の混合を on-shell の場の定数で扱うとき、右巻きの定数（式~\eqref{eq:dZRos}）は左巻きと形が違い、それを間違えると電磁気のゲージ不変性が破れる。

\paragraph{残る手順。}手順 (8)〜(10) はまだ行っていない。とくに、$\Delta F=2$ の箱の図の和のゲージ不変性（手順 (8)）が、本ノートの中心の確かめとして残る。本ノートで扱わないのは、ニュートリノの質量と PMNS 行列、世代を変える過程の QCD の補正の繰り込み群による改良（$b\to s\gamma$ の大きな QCD の補正）と、ハドロンの行列要素の計算である。

\appendix

\section{Feynman 則の導出}\label{app:rules}

\paragraph{荷電の電流。}SM ノートの付録~A のとおり、共変微分の荷電の部分は $T^1W^1+T^2W^2=\frac1{\sqrt2}(\sigma^+W^++\sigma^-W^-)$ である。式~\eqref{eq:QL} の二重項では $\bar Q_L\gamma^\mu\sigma^+Q_L=\bar u_L\gamma^\mu Vd_L$ なので、$\bar u_i\gamma^\mu d_jW^+_\mu$ の係数は $\frac g{\sqrt2}V_{ij}$ で、頂点は $\frac{ie}{\sqrt2\sw}V_{ij}\gamma^\mu P_L$ になる。エルミート共役から $\bar d_j\gamma^\mu u_iW^-_\mu$ の係数は $\frac g{\sqrt2}V_{ij}^*$ である。中性の部分（$T^3$、$Y$、$T^a_{\mathrm{color}}$）は二重項の各成分の中で閉じ、下の成分では $\overline{(Vd_L)}\gamma^\mu(Vd_L)=\bar d_LV^\dagger V\gamma^\mu d_L=\bar d_L\gamma^\mu d_L$ になる。

\paragraph{湯川結合。}式~\eqref{eq:Y} に $Y_d=V\hat y_d$、$Y_u=\hat y_u$ を入れる。$d$ 型の項は、二重項の上の成分から $-\bar u_{Li}(V\hat y_d)_{ij}(-iG^+)d_{Rj}=iV_{ij}y_{d_j}\bar u_{Li}d_{Rj}G^+$、下の成分から $-\overline{(Vd_L)}V\hat y_d\frac{v+H+iG^0}{\sqrt2}d_R=-\bar d_L\hat y_d\frac{v+H+iG^0}{\sqrt2}d_R$（$V^\dagger V=1$）で、中性の部分は SM ノートと同じ対角の式になる。$u$ 型の項は $\tilde\Phi=((v+H-iG^0)/\sqrt2,-iG^-)$ の下の成分から $-\overline{(Vd_L)}\hat y_u(-iG^-)u_R=iG^-\bar d_{Lj}V_{ij}^*y_{u_i}u_{Ri}$ を与える。したがって $\bar u_id_jG^+$ の係数は、$d$ 型の項の $iV_{ij}y_{d_j}P_R$ と、$u$ 型の項のエルミート共役 $-iV_{ij}y_{u_i}P_L$ の和で、頂点（$i$ を掛ける）は $V_{ij}(y_{u_i}P_L-y_{d_j}P_R)=\frac{\sqrt2}vV_{ij}(m_{u_i}P_L-m_{d_j}P_R)$ になる。$\bar d_ju_iG^-$ はそのエルミート共役から $\frac{\sqrt2}vV_{ij}^*(m_{d_j}P_L-m_{u_i}P_R)$ である。

\section{木のレベルの導出}\label{app:tree}

$W^+\to u_i\bar d_j$ と $t\to d_jW^+$ の振幅は、SM ノートの付録~B の振幅に $V_{ij}$（$V_{tj}$）を掛けたものである。スピンと偏極の和の中で $V$ は $|V_{ij}|^2$ としてだけ現れるので、式~\eqref{eq:widthsckm} になる。$u\bar u\to W^+W^-$ の $t$ チャンネルの図は、$u\to d_kW^+$ の頂点 $V_{uk}^*$（$\bar d_ku\bar W$）と $\bar u\to\bar d_kW^-$ の頂点 $V_{uk}$（$\bar ud_kW$）を持ち、内線の $d_k$ の伝播関数 $i(\slashed q+m_{d_k})/(q^2-m_{d_k}^2)$ を挟む。$m_{d_k}=0$ では内線によらないので、三つの図の和は $\sum_k|V_{uk}|^2=1$ 倍の1個の図に等しい。

式~\eqref{eq:qqbar} は $\mu$ 崩壊と同じトレースから出る。$u_i\bar d_j\to\nu_ee^+$ の振幅は
\begin{equation*}
i\mathcal M=\Bigl(\frac{ig}{\sqrt2}\Bigr)^2V_{ij}^*\frac{-i}{s-M_W^2}\bigl[\bar u(k_1)\gamma^\mu P_Lv(k_2)\bigr]\bigl[\bar v(p_2)\gamma_\mu P_Lu(p_1)\bigr]
\end{equation*}
で（質量0の流れでは伝播関数の $k_\mu k_\nu$ の部分が消える）、二つの流れの積のスピンの和は $16(p_1\cdot k_2)(p_2\cdot k_1)=4u^2$ である（$\mu$ 崩壊の $\sum|\mathcal M|^2=128G_F^2(p_\mu\cdot p_{\bar\nu_e})(p_e\cdot p_{\nu_\mu})$ と同じ形で、粒子と反粒子の組み合わせを入れ替えたもの）。$(g^2/2)^2\cdot4u^2$ に色の和 $N_c$ を掛けて1つ目の式になる。$u\bar u\to d\bar s$ は $[\bar u(k_1)\gamma^\mu P_Lu(p_1)][\bar v(p_2)\gamma_\mu P_Lv(k_2)]$ で、同じく $16(p_1\cdot k_2)(p_2\cdot k_1)=4u^2$（$u$ と $\bar s$、$\bar u$ と $d$ の組）と色の和 $N_c^2$ から2つ目の式になる。$\int\dd\Omega\,u^2=\frac{16\pi}3\frac{s^2}4$、平均 $\frac1{36}$、位相空間 $1/(64\pi^2s)$ から断面積が出る。

\section{湯川行列の繰り込みと CKM 行列の走りの導出}\label{app:rg}

\paragraph{行列の $\beta$ 関数。}SM ノートの付録~C の湯川結合の導出を行列に広げる。場の異常次元は $\Phi$ で $\mathcal T$、$Q_L$ で $\frac12(Y_uY_u^\dagger+Y_dY_d^\dagger)$（左から掛かる）、$u_R$ で $Y_u^\dagger Y_u$（右から掛かる）、$d_R$ で $Y_d^\dagger Y_d$ である（いずれも $K$ で割る。一般の公式の $\frac12Y_2(F)$~\cite{mv1}）。$\Phi$ の内線が $Q_L$ を $d_R$ に変えて $\tilde\Phi$ の頂点を挟む図は $-2Y_dY_d^\dagger Y_u$ を与える。和は
\begin{equation*}
\frac12\bigl(H_u+H_d\bigr)Y_u+Y_u\bigl(Y_u^\dagger Y_u\bigr)+\mathcal TY_u-2H_dY_u=\Bigl[\frac32\bigl(H_u-H_d\bigr)+\mathcal T\Bigr]Y_u
\end{equation*}
で（$Y_uY_u^\dagger Y_u=H_uY_u$）、ゲージ場の部分は SM ノートと同じく単位行列の倍である。これが式~\eqref{eq:betaY} の1行目で、2行目は $u$ と $d$ を入れ替えたものである。

\paragraph{質量と $V$ の走り。}$H_u=Y_uY_u^\dagger$ の走りは $\dd H_u/\dd\ln\mu=\beta_{Y_u}Y_u^\dagger+Y_u\beta_{Y_u}^\dagger$ で、式~\eqref{eq:betaY} から
\begin{equation*}
K\frac{\dd H_u}{\dd\ln\mu}=3H_u^2-\frac32\bigl(H_dH_u+H_uH_d\bigr)+2a_uH_u,\qquad
K\frac{\dd H_d}{\dd\ln\mu}=3H_d^2-\frac32\bigl(H_uH_d+H_dH_u\bigr)+2a_dH_d
\end{equation*}
である（$a_u$、$a_d$ はゲージ結合と $\mathcal T$ の項）。$H_u=U\,D_u\,U^\dagger$ と書き、$\dd U/\dd\ln\mu=UA_u$（$A_u$ は反エルミート、対角の成分は0とする）とすると、$U^\dagger(\dd H_u/\dd\ln\mu)U=\dd D_u/\dd\ln\mu+[A_u,D_u]$ なので、上型の質量の基底（$U=1$、$H_d=VD_dV^\dagger$）で対角の成分から
\begin{equation*}
K\frac{\dd y_{u_i}^2}{\dd\ln\mu}=3y_{u_i}^4-3y_{u_i}^2(VD_dV^\dagger)_{ii}+2a_uy_{u_i}^2
\end{equation*}
が、非対角の成分から $(A_u)_{ik}=-\frac3{2K}\frac{y_{u_i}^2+y_{u_k}^2}{y_{u_k}^2-y_{u_i}^2}(VD_dV^\dagger)_{ik}$（$i\ne k$）が出る。下型の質量の基底では同様に $(A_d)_{kj}=-\frac3{2K}\frac{y_{d_k}^2+y_{d_j}^2}{y_{d_j}^2-y_{d_k}^2}(V^\dagger D_uV)_{kj}$ である。$V=U_u^\dagger U_d$ から $\dd V/\dd\ln\mu=-A_uV+VA_d$ で、式~\eqref{eq:yrun}、\eqref{eq:Vrun} になる。式~\eqref{eq:Vhier} は、$V_{cb}$ なら1項目の $k=t$ の項 $\frac32|V_{tb}|^2y_b^2V_{cb}$ と、2項目の $k=d,s$ の項 $-\frac32y_t^2V_{tb}(V_{cd}V_{td}^*+V_{cs}V_{ts}^*)=\frac32y_t^2|V_{tb}|^2V_{cb}$（行のユニタリー性）から来る。

\paragraph{自己エネルギーの紫外の極。}対称な相の $\MSbar$ の場の繰り込みをエルミートにとると、湯川の部分は $\delta Z_Q=-\frac{\DUV}{2K}(H_u+H_d)$、$\delta Z_{u_R}=-\frac{\DUV}KY_u^\dagger Y_u$、$\delta Z_{d_R}=-\frac{\DUV}KY_d^\dagger Y_d$ で、ゲージ場の部分は単位行列の倍である。下型の質量の基底では $d_L$ の場の定数は $V^\dagger\delta Z_QV$ の下の成分で、その非対角の成分は $-\frac{\DUV}{2K}(V^\dagger D_uV)_{ij}$、$d_R$ の非対角の成分は0である。質量の項の対抗項は $\frac v{\sqrt2}V^\dagger\delta Y_d$ で、式~\eqref{eq:betaY} から非対角の成分は $-\frac34\frac{\DUV}K(V^\dagger D_uV)_{ij}m_{d_j}$ になる。式~\eqref{eq:Gammaq} の自己エネルギーにこれらの対抗項を加えたものが有限であるという条件から
\begin{equation*}
\Sigma^L_{ij}\big|_{\mathrm{UV}}=-\delta Z^L_{ij},\qquad
m_j\Sigma^S_{ij}\big|_{\mathrm{UV}}=\delta M_{ij}+\frac12\delta Z^{L\dagger}_{ij}m_j+\frac12m_i\delta Z^R_{ij}=-\frac{\DUV}K(V^\dagger D_uV)_{ij}m_j
\end{equation*}
で、式~\eqref{eq:uvoff} になる（$P_L$ の側も同じ $\Sigma^S$ を与える）。$G^\pm$ のループを直接求めると、$\Sigma_{ij}=-\frac2{v^2K}\sum_k\lambda^{(ij)}_k\bigl[\slashed pP_Lm_{u_k}^2B_1+\slashed pP_Rm_{d_i}m_{d_j}B_1+(m_{d_i}P_L+m_{d_j}P_R)m_{u_k}^2B_0\bigr]$（$B$ の引数は $(p^2,m_{u_k}^2,\xi_WM_W^2)$）で、$B_1|_{\mathrm{UV}}=-\DUV/2$、$B_0|_{\mathrm{UV}}=\DUV$ から同じ式になる（$\slashed pP_R$ の極はユニタリー性で非対角の成分が0）。

\paragraph{$\delta V^{\mathrm{DS}}$ の紫外の極。}式~\eqref{eq:uvoff} を式~\eqref{eq:dZos} に入れると、$i\ne j$ で $\delta Z^L_{ij}|_{\mathrm{UV}}=-\frac{m_j^2+2m_i^2}{m_i^2-m_j^2}X_{ij}\frac{\DUV}K$、$X$ がエルミートなので $(\delta Z^L-\delta Z^{L\dagger})_{ij}|_{\mathrm{UV}}=-3\frac{m_i^2+m_j^2}{m_i^2-m_j^2}X_{ij}\frac{\DUV}K$ である。式~\eqref{eq:dVds} の $d$ 型の部分は $\frac34\frac{\DUV}K\sum_{k\ne j}V_{ik}\frac{m_{d_k}^2+m_{d_j}^2}{m_{d_k}^2-m_{d_j}^2}(V^\dagger D_uV)_{kj}$ で、これは $\frac{\DUV}2(VA_d)_{ij}$ に等しい。$u$ 型の部分も同様に $-\frac{\DUV}2(A_uV)_{ij}$ になり、和が式~\eqref{eq:dV} である。

\section{検算の細部}\label{app:details}

\paragraph{模型ファイルのパラメータ。}\texttt{SM\_CKM.fr} の外部パラメータは、SM ノートの \texttt{SM.fr} のもの（ブロック \texttt{SMINPUTS}、\texttt{MASS}、\texttt{DECAY}、\texttt{GAUGE}）に、ブロック \texttt{CKMBLOCK} の \texttt{s12}（1番）、\texttt{s23}（2番）、\texttt{s13}（3番）、\texttt{dCKM}（4番）を加えたものである。内部パラメータ \texttt{c12}、\texttt{c23}、\texttt{c13} は $\sqrt{1-s^2}$、\texttt{CKM} は式~\eqref{eq:Vstd} の複素行列で、UFO では \texttt{CKM1x1}〜\texttt{CKM3x3}（\texttt{cmath.exp} を含む複素数）になる。結合の次数、粒子の名前と PDG の番号は SM ノートと同じである。

\paragraph{手計算の数値の確かめ。}第~\ref{sec:rg} 節の式~\eqref{eq:Vrun} の確かめと基準点の数値（$|V|$、$J$、$K\,\dd\ln|V_{ij}|/\dd\ln\mu$、$\lambda^{(ij)}_k$）は \texttt{ckm\_rg\_handcheck.py}（numpy だけを使う）が出し、\texttt{results/ckm\_rg\_handcheck.csv} に書く。模型ファイルは使わない（式~\eqref{eq:betaY} を模型から導くのは手順 (4)）。

\paragraph{比べ方。}\texttt{generate\_SM\_CKM.wl} は SM ノートの \texttt{generate\_SM.wl} と同じく、Dirac 構造、計量、運動量、色の行列を評価されない記号に置き換えてから表と比べる。フェルミオンの頂点は、外線の世代の添字に $1,2,3$ を入れ、和をとる世代の添字（\texttt{Index[Generation, Generation\$1]} の形）を項ごとに足してから比べる（その添字を持たない項を3倍しないため）。標準のパラメータ表示の値は \texttt{ComplexExpand} で角度と位相を実数として整理する。

\begin{thebibliography}{99}
\bibitem{smnote} 標準模型：$\MSbar$ の電弱の1ループの繰り込みと検算（$V_{\mathrm{CKM}}=1$、FJ のタドポール、一般の $R_\xi$ ゲージ、$d=4$）（本リポジトリ \nolinkurl{models/sm/ResearchNote/note.pdf}、2026年10月17日の版）。式、表、手順の番号はこの版のもの。
\bibitem{env} ループ計算環境の構築記録（本リポジトリ \nolinkurl{environment/environment.pdf}）。
\bibitem{cabibbo63} N. Cabibbo, ``Unitary symmetry and leptonic decays,'' Phys. Rev. Lett. \textbf{10}, 531 (1963)。
\bibitem{km73} M. Kobayashi, T. Maskawa, ``CP-violation in the renormalizable theory of weak interaction,'' Prog. Theor. Phys. \textbf{49}, 652 (1973)。
\bibitem{gim70} S. L. Glashow, J. Iliopoulos, L. Maiani, ``Weak interactions with lepton-hadron symmetry,'' Phys. Rev. D \textbf{2}, 1285 (1970)。
\bibitem{ck84} L.-L. Chau, W.-Y. Keung, ``Comments on the parametrization of the Kobayashi-Maskawa matrix,'' Phys. Rev. Lett. \textbf{53}, 1802 (1984)。
\bibitem{jarlskog85} C. Jarlskog, ``Commutator of the quark mass matrices in the standard electroweak model and a measure of maximal CP nonconservation,'' Phys. Rev. Lett. \textbf{55}, 1039 (1985)。
\bibitem{il81} T. Inami, C. S. Lim, ``Effects of superheavy quarks and leptons in low-energy weak processes $K_L\to\mu\bar\mu$, $K^+\to\pi^+\nu\bar\nu$ and $K^0\leftrightarrow\bar K^0$,'' Prog. Theor. Phys. \textbf{65}, 297 (1981)。
\bibitem{ds90} A. Denner, T. Sack, ``Renormalization of the quark mixing matrix,'' Nucl. Phys. B \textbf{347}, 203 (1990)。
\bibitem{ggm99} P. Gambino, P. A. Grassi, F. Madricardo, ``Fermion mixing renormalization and gauge invariance,'' Phys. Lett. B \textbf{454}, 98 (1999), arXiv:hep-ph/9811470。
\bibitem{yamada01} Y. Yamada, ``Gauge dependence of the on-shell renormalized mixing matrices,'' Phys. Rev. D \textbf{64}, 036008 (2001), arXiv:hep-ph/0103046。
\bibitem{dkr04} A. Denner, E. Kraus, M. Roth, ``Physical renormalization condition for the quark-mixing matrix,'' Phys. Rev. D \textbf{70}, 033002 (2004), arXiv:hep-ph/0402130。
\bibitem{ks06} B. A. Kniehl, A. Sirlin, ``Simple approach to renormalize the Cabibbo-Kobayashi-Maskawa matrix,'' Phys. Rev. Lett. \textbf{97}, 221801 (2006), arXiv:hep-ph/0608306。
\bibitem{babu87} K. S. Babu, ``Renormalization-group analysis of the Kobayashi-Maskawa matrix,'' Z. Phys. C \textbf{35}, 69 (1987)。
\bibitem{bhmp99} C. Balzereit, T. Hansmann, T. Mannel, B. Plümper, ``The renormalization group evolution of the CKM matrix,'' Eur. Phys. J. C \textbf{9}, 197 (1999), arXiv:hep-ph/9810350。
\bibitem{denner93} A. Denner, ``Techniques for the calculation of electroweak radiative corrections at the one-loop level and results for W-physics at LEP200,'' Fortsch. Phys. \textbf{41}, 307 (1993), arXiv:0709.1075。
\bibitem{mv1} M. E. Machacek, M. T. Vaughn, ``Two-loop renormalization group equations in a general quantum field theory (I). Wave function renormalization,'' Nucl. Phys. B \textbf{222}, 83 (1983)。
\bibitem{mv2} M. E. Machacek, M. T. Vaughn, ``(II). Yukawa couplings,'' Nucl. Phys. B \textbf{236}, 221 (1984)。
\bibitem{luo03} M.-X. Luo, H.-W. Wang, Y. Xiao, ``Two-loop renormalization group equations in general gauge field theories,'' Phys. Rev. D \textbf{67}, 065019 (2003), arXiv:hep-ph/0211440。
\bibitem{arason92} H. Arason \textit{et al.}, ``Renormalization-group study of the standard model and its extensions: The standard model,'' Phys. Rev. D \textbf{46}, 3945 (1992)。
\bibitem{sirlin78} A. Sirlin, ``Current algebra formulation of radiative corrections in gauge theories and the universality of the weak interactions,'' Rev. Mod. Phys. \textbf{50}, 573 (1978)。
\bibitem{sirlin82} A. Sirlin, ``Large $m_W$, $m_Z$ behaviour of the $O(\alpha)$ corrections to semileptonic processes mediated by $W$,'' Nucl. Phys. B \textbf{196}, 83 (1982)。
\bibitem{bbl96} G. Buchalla, A. J. Buras, M. E. Lautenbacher, ``Weak decays beyond leading logarithms,'' Rev. Mod. Phys. \textbf{68}, 1125 (1996), arXiv:hep-ph/9512380。
\bibitem{pdg24} S. Navas \textit{et al.} (Particle Data Group), Phys. Rev. D \textbf{110}, 030001 (2024)。CKM 行列のレビュー（標準のパラメータ表示と大域的な当てはめ）。
\end{thebibliography}

\end{document}
